% Public Finance — Economics Upper Sixth — Cours signé OnBuch+
% Official MINESEC syllabus. Compiler avec : tectonic ECO05-public-finance.tex
\newcommand{\DOCMATIERE}{Economics}
\newcommand{\DOCNIVEAU}{Upper Sixth}
\newcommand{\DOCMODULE}{Upper Sixth — Economics}
\newcommand{\DOCLECON}{5}
\newcommand{\DOCTITRE}{Public Finance}
% OnBuch+ — préambule « cours signés » ENRICHI (compilé avec tectonic / XeLaTeX)
% Système visuel « Pop » d'OnBuch+ : fond crème (jamais blanc), bordures encre
% épaisses, ombres « dures » décalées, titres Archivo Black en capitales.
% Version MATHÉMATIQUES : graphes (pgfplots/tikz), boîtes pédagogiques
% (définition, propriété, méthode, attention, le savais-tu, exercice résolu,
% exercice, TP, bilan), page de garde et signature OnBuch+ ancrée en bas.
%
% Macros attendues dans main.tex AVANT \input{preamble} :
%   \DOCMATIERE  \DOCNIVEAU  \DOCMODULE  \DOCLECON (numéro)  \DOCTITRE
\documentclass[11pt]{article}

\usepackage{fontspec}
\usepackage[english]{babel}
\usepackage[a4paper,margin=2cm,top=3.4cm,bottom=2.8cm,headheight=1.4cm,footskip=1.3cm]{geometry}
\usepackage{amsmath,amssymb,amsfonts}
\usepackage{mathtools} % \xmapsto, \coloneqq... souvent utilisés par les modèles
\usepackage{cancel} % simplifications barrées (bilans de forces)
% \abs{z} (module, valeur absolue) et \norm{u} (norme), utilisés par les modèles
\providecommand{\abs}{}\let\abs\relax\DeclarePairedDelimiter{\abs}{\lvert}{\rvert}
\providecommand{\norm}{}\let\norm\relax\DeclarePairedDelimiter{\norm}{\lVert}{\rVert}
\usepackage[table]{xcolor} % [table] : fournit aussi \rowcolors (alternance de lignes)
\usepackage{pagecolor}
\usepackage{tikz}
\usetikzlibrary{shapes.symbols,circuits.logic.US,circuits.logic.IEC,arrows.meta,positioning,calc,shapes.geometric,shapes.misc,shapes.arrows,decorations.pathmorphing,decorations.pathreplacing,decorations.markings,patterns,shadows,mindmap,trees,fit,backgrounds,matrix,intersections,angles,quotes}
\usepackage{pgfplots}
\usepackage[version=4]{mhchem} % équations nucléaires \ce{^{235}_{92}U}
\usepackage[european,siunitx]{circuitikz} % schémas électriques
\pgfplotsset{compat=1.18}
\usepgfplotslibrary{fillbetween,groupplots,polar,statistics,dateplot} % aires entre courbes (calcul intégral)
\usepackage{enumitem}
\usepackage{fancyhdr}
% [most] charge toutes les bibliothèques tcolorbox (listings, minted, raster,
% documentation...) et gonfle inutilement la mémoire TeX sur un document long
% (~300 boîtes) : on ne charge que ce qui est réellement utilisé.
\usepackage[skins,breakable]{tcolorbox}
\usepackage{graphicx}
\usepackage{colortbl}
\usepackage{booktabs}
\usepackage{tabularx}
\usepackage{diagbox} % case d'en-tête diagonale des tableaux à double entrée
% Colonnes extensibles centrée / gauche / droite (C, L, R), souvent utilisées par les modèles
\newcolumntype{C}{>{\centering\arraybackslash}X}
\newcolumntype{L}{>{\raggedright\arraybackslash}X}
\newcolumntype{R}{>{\raggedleft\arraybackslash}X}
\usepackage{array}
\usepackage{multicol}
\usepackage{siunitx}
% parse-numbers=false : accepte aussi \SI{2,5 \times 10^{-3}}{...}
\sisetup{locale=UK,output-decimal-marker={.},group-minimum-digits=4,parse-numbers=false}
\DeclareSIUnit\ohm{\ensuremath{\symup{\Omega}}} % U+2126 absent de la police
\DeclareSIUnit\division{div} % réglages d'oscilloscope (ms/div, V/div)
\DeclareSIUnit\year{yr}\DeclareSIUnit\annum{yr}\DeclareSIUnit\Ma{Ma}\DeclareSIUnit\tr{tr} % tours (vitesse de rotation, ex. \SI{1500}{\tr\per\minute})
\usepackage{qrcode}
% \degree : commande fréquemment utilisée par les modèles pour un angle ou une
% température en dehors de \SI{}{} (non fournie par les paquets chargés).
\providecommand{\degree}{^{\circ}}
% \qed (fin de démonstration), utilisé par les modèles sans amsthm
\providecommand{\qed}{\hfill$\square$}
% \barre : double barre des valeurs interdites dans un tableau de variations (tkz-tab)
\providecommand{\barre}{\ensuremath{\|}}
% environnement proof (amsthm non chargé) utilisé par les modèles
\ifdefined\proof\else\newenvironment{proof}[1][Proof]{\par\noindent\textit{#1.}\ }{\qed\par}\fi
\usepackage{lastpage}
\usepackage{hyperref}
\hypersetup{hidelinks}

% ============================================================
% Palette « Pop » OnBuch+ — mêmes hex que la classe P (plus_ui.dart)
% ============================================================
\definecolor{poporange}{HTML}{F59321}
\definecolor{popdark}{HTML}{E07A0C}
\definecolor{popcream}{HTML}{FFF6E8}
\definecolor{popink}{HTML}{1C1714}
\definecolor{poppurple}{HTML}{7A5AE0}
\definecolor{popgreen}{HTML}{2E9E6A}
\definecolor{poppink}{HTML}{E0457B}
\definecolor{popblue}{HTML}{2F7DE1}
\definecolor{popgold}{HTML}{C98A12}
\definecolor{popmuted}{HTML}{8A7F76}
\definecolor{popline}{HTML}{EADFD2}
\definecolor{popgray}{HTML}{8A7F76}
\colorlet{popgrey}{popgray}
% Alias tolérants (le modèle nomme parfois les couleurs différemment)
\colorlet{popwhite}{white}
\colorlet{popblack}{popink}
\colorlet{popred}{poppink}
\colorlet{popyellow}{popgold}
% Teintes claires pour fonds de tableaux / zones
\colorlet{popblueL}{popblue!10!white}
\colorlet{poporangeL}{poporange!14!white}
\colorlet{popgreenL}{popgreen!12!white}
\colorlet{poppurpleL}{poppurple!12!white}
\colorlet{poppinkL}{poppink!12!white}
\colorlet{popgoldL}{popgold!14!white}

\pagecolor{popcream}

% ============================================================
% Typographie
% ============================================================
\defaultfontfeatures{Ligatures=TeX}
\setmainfont{PlusJakartaSans-Medium.ttf}[Path=../../../fonts/,BoldFont=PlusJakartaSans-Bold.ttf]
% Maths en sans-serif (Fira Math) avec chiffres et lettres droites de Plus
% Jakarta, pour que les formules se fondent dans le texte.
\usepackage[math-style=ISO,bold-style=ISO]{unicode-math}
\setmathfont{FiraMath-Regular.otf}[Path=../../../fonts/]
\setmathfont{PlusJakartaSans-Medium.ttf}[Path=../../../fonts/,range={up/num,up/{Latin,latin},\mathup}]
% (icomma retiré : décimales avec point en anglais)
% Caractères mathématiques Unicode absents des polices de texte (Plus Jakarta,
% Archivo Black) : rendus par leur équivalent mathématique, en texte comme en
% formule — sinon ils s'affichent en carré vide (ex. « Arithmétique dans ℤ »).
\usepackage{newunicodechar}
\newunicodechar{ℝ}{\ensuremath{\mathbb{R}}}
\newunicodechar{ℤ}{\ensuremath{\mathbb{Z}}}
\newunicodechar{ℂ}{\ensuremath{\mathbb{C}}}
\newunicodechar{ℕ}{\ensuremath{\mathbb{N}}}
\newunicodechar{ℚ}{\ensuremath{\mathbb{Q}}}
\newunicodechar{𝔻}{\ensuremath{\mathbb{D}}}
\newunicodechar{α}{\ensuremath{\alpha}}
\newunicodechar{β}{\ensuremath{\beta}}
\newunicodechar{γ}{\ensuremath{\gamma}}
\newunicodechar{δ}{\ensuremath{\delta}}
\newunicodechar{ε}{\ensuremath{\varepsilon}}
\newunicodechar{θ}{\ensuremath{\theta}}
\newunicodechar{λ}{\ensuremath{\lambda}}
\newunicodechar{μ}{\ensuremath{\mu}}
\newunicodechar{σ}{\ensuremath{\sigma}}
\newunicodechar{τ}{\ensuremath{\tau}}
\newunicodechar{φ}{\ensuremath{\varphi}}
\newunicodechar{ψ}{\ensuremath{\psi}}
\newunicodechar{ω}{\ensuremath{\omega}}
\newunicodechar{Σ}{\ensuremath{\Sigma}}
% majuscules produites par \MakeUppercase dans les titres (α-aminés -> Α)
\newunicodechar{Α}{\ensuremath{\alpha}}
\newunicodechar{Β}{\ensuremath{\beta}}
\newunicodechar{Γ}{\ensuremath{\Gamma}}
\newunicodechar{Δ}{\ensuremath{\Delta}}
\newunicodechar{Ε}{\ensuremath{\varepsilon}}
\newunicodechar{Θ}{\ensuremath{\Theta}}
\newunicodechar{Λ}{\ensuremath{\Lambda}}
\newunicodechar{Μ}{\ensuremath{\mu}}
\newunicodechar{Τ}{\ensuremath{\tau}}
\newunicodechar{Φ}{\ensuremath{\Phi}}
\newunicodechar{Ψ}{\ensuremath{\Psi}}
\newunicodechar{Ω}{\ensuremath{\Omega}}
\newunicodechar{↦}{\ensuremath{\mapsto}}
\newunicodechar{∈}{\ensuremath{\in}}
\newunicodechar{∉}{\ensuremath{\notin}}
\newunicodechar{∧}{\ensuremath{\wedge}}
\newunicodechar{⇔}{\ensuremath{\Leftrightarrow}}
\newunicodechar{⟺}{\ensuremath{\Longleftrightarrow}}
\newunicodechar{⇒}{\ensuremath{\Rightarrow}}
\newunicodechar{⟹}{\ensuremath{\Longrightarrow}}
\newunicodechar{∘}{\ensuremath{\circ}}
\newunicodechar{∪}{\ensuremath{\cup}}
\newunicodechar{∩}{\ensuremath{\cap}}
\newunicodechar{≡}{\ensuremath{\equiv}}
\newunicodechar{⊂}{\ensuremath{\subset}}
\newunicodechar{⊥}{\ensuremath{\perp}}
\newunicodechar{∃}{\ensuremath{\exists}}
\newunicodechar{∀}{\ensuremath{\forall}}
\newunicodechar{∛}{\ensuremath{\sqrt[3]{\,}}}
\newunicodechar{∜}{\ensuremath{\sqrt[4]{\,}}}
\newunicodechar{⃗}{\ensuremath{{}^{\to}}}
\newunicodechar{ⁿ}{\ifmmode^{n}\else\textsuperscript{\ensuremath{n}}\fi}
\newunicodechar{ˣ}{\ifmmode^{x}\else\textsuperscript{\ensuremath{x}}\fi}
\newunicodechar{ᵇ}{\ifmmode^{b}\else\textsuperscript{\ensuremath{b}}\fi}
\newunicodechar{ʳ}{\ifmmode^{r}\else\textsuperscript{\ensuremath{r}}\fi}
\newunicodechar{ᵘ}{\ifmmode^{u}\else\textsuperscript{\ensuremath{u}}\fi}
\newunicodechar{ʸ}{\ifmmode^{y}\else\textsuperscript{\ensuremath{y}}\fi}
\newunicodechar{ᵃ}{\ifmmode^{a}\else\textsuperscript{\ensuremath{a}}\fi}
\newunicodechar{ᵉ}{\ifmmode^{e}\else\textsuperscript{\ensuremath{e}}\fi}
\newunicodechar{ᵏ}{\ifmmode^{k}\else\textsuperscript{\ensuremath{k}}\fi}
\newunicodechar{ᵖ}{\ifmmode^{p}\else\textsuperscript{\ensuremath{p}}\fi}
\newunicodechar{ᴺ}{\ifmmode^{N}\else\textsuperscript{\ensuremath{N}}\fi}
\newunicodechar{ᶜ}{\ifmmode^{c}\else\textsuperscript{\ensuremath{c}}\fi}
\newunicodechar{ʰ}{\ifmmode^{h}\else\textsuperscript{\ensuremath{h}}\fi}
\newunicodechar{ᵠ}{\ifmmode^{q}\else\textsuperscript{\ensuremath{q}}\fi}
\newunicodechar{ₙ}{\ifmmode_{n}\else\textsubscript{\ensuremath{n}}\fi}
\newunicodechar{ₐ}{\ifmmode_{a}\else\textsubscript{\ensuremath{a}}\fi}
\newunicodechar{ᵢ}{\ifmmode_{i}\else\textsubscript{\ensuremath{i}}\fi}
\newunicodechar{ᵣ}{\ifmmode_{r}\else\textsubscript{\ensuremath{r}}\fi}
\newunicodechar{ₖ}{\ifmmode_{k}\else\textsubscript{\ensuremath{k}}\fi}
\newunicodechar{ᵦ}{\ifmmode_{\beta}\else\textsubscript{\ensuremath{\beta}}\fi}
\newunicodechar{ₓ}{\ifmmode_{x}\else\textsubscript{\ensuremath{x}}\fi}
\newfontfamily\popdisplay{ArchivoBlack-Regular.ttf}[Path=../../../fonts/]
\newfontfamily\poplogo{SpaceGrotesk-Bold.ttf}[Path=../../../fonts/]
\newfontfamily\popsemi{PlusJakartaSans-SemiBold.ttf}[Path=../../../fonts/]
\newfontfamily\popxbold{PlusJakartaSans-ExtraBold.ttf}[Path=../../../fonts/]
\newfontfamily\popxboldls{PlusJakartaSans-ExtraBold.ttf}[Path=../../../fonts/,LetterSpace=3.0]

\setlist[enumerate]{leftmargin=1.7em,itemsep=5pt,topsep=4pt}
\setlist[itemize]{leftmargin=1.7em,itemsep=4pt,topsep=4pt,label={\color{poporange}\textbullet}}
\setlength{\parindent}{0pt}
\setlength{\parskip}{5pt}
\renewcommand{\arraystretch}{1.25}

% pgfplots : style Pop par défaut
\pgfplotsset{
  every axis/.append style={axis line style={popink,line width=1pt},tick style={popink},
    grid style={popline,line width=0.5pt},label style={font=\small},tick label style={font=\footnotesize}},
  popcurve/.style={poporange,line width=1.8pt,no markers,smooth},
}
% Même style disponible dans un \draw TikZ simple (hors pgfplots)
\tikzset{popcurve/.style={poporange,line width=1.8pt}}
\tikzset{popgrid/.style={popline,very thin}}
\colorlet{popcurve}{poporange} % le nom du style est parfois utilisé comme couleur

% ============================================================
% Pill / squiggle
% ============================================================
\newcommand{\poppill}[3]{%
  \tikz[baseline=-0.55ex]\node[draw=popink,line width=1.2pt,rounded corners=2.6pt,%
    fill=#2,inner xsep=6.5pt,inner ysep=3pt]{%
    \popxboldls\fontsize{7}{7}\selectfont\color{#3}\MakeUppercase{#1}};%
}
\newcommand{\popsquiggle}[1]{%
  \tikz[baseline=-0.4ex]{
    \draw[#1,line width=2.6pt,line cap=round]
      plot [smooth] coordinates {(0,0.09) (0.34,0.01) (0.68,0.21) (1.02,0.01) (1.36,0.10)};
  }%
}
% Mot-clé surligné (marqueur orange sous le texte)
\newcommand{\cle}[1]{{\popxbold\color{popdark}#1}}

% ============================================================
% En-tête / pied
% ============================================================
\pagestyle{fancy}
\fancyhf{}
\renewcommand{\headrulewidth}{0pt}
\renewcommand{\footrulewidth}{0pt}
\lhead{\raisebox{-0.15ex}{\poplogo\fontsize{13}{13}\selectfont\color{popink}OnBuch\color{poporange}+}}
\rhead{\poppill{\DOCMATIERE\ \textbullet\ \DOCNIVEAU\ \textbullet\ Chapter \DOCLECON}{white}{popink}}
\lfoot{\popxbold\fontsize{7.6}{7.6}\selectfont\color{popmuted}\MakeUppercase{Signed course OnBuch+}\ \textbullet\ {\popsemi\DOCTITRE}}
\rfoot{\tikz[baseline=-0.6ex]\node[rounded corners=8.5pt,fill=popink,minimum height=17pt,minimum width=17pt,inner xsep=5pt,inner sep=0]{\popxbold\fontsize{8}{8}\selectfont\color{popcream}\thepage\,/\,\pageref*{LastPage}};}

% ============================================================
% Titres
% ============================================================
\newcommand{\chaptitre}[2]{%
  \begin{tcolorbox}[enhanced,colback=poporange,colframe=popink,boxrule=2.6pt,arc=11pt,
      drop shadow=popink,left=15pt,right=15pt,top=12pt,bottom=11pt,before skip=3pt,after skip=18pt]
    \poppill{#2}{popink}{popcream}\par\vspace{9pt}
    {\raggedright\hyphenpenalty=10000\exhyphenpenalty=10000\popdisplay\fontsize{22}{24}\selectfont\color{popcream}\MakeUppercase{#1}\par}
  \end{tcolorbox}
}
% Section numérotée : pastille orange avec le numéro + titre Archivo
\newcounter{popsec}
\newcommand{\coursec}[1]{%
  \stepcounter{popsec}\setcounter{popsub}{0}%
  \par\vspace{16pt}\noindent
  \begin{minipage}{\linewidth}
  \tikz[baseline=-0.7ex]\node[draw=popink,line width=1.6pt,fill=poporange,rounded corners=4pt,minimum size=22pt,inner sep=0]{\popdisplay\fontsize{12}{12}\selectfont\color{popcream}\thepopsec};%
  \hspace{8pt}{\popdisplay\fontsize{15}{17}\selectfont\color{popink}\MakeUppercase{#1}}\par
  \vspace{4pt}\noindent{\color{poporange}\rule{36pt}{3.6pt}}
  \end{minipage}\par\nopagebreak\vspace{8pt}
}
\newcounter{popsub}
\newcommand{\courssub}[1]{%
  \stepcounter{popsub}%
  \par\vspace{9pt}\noindent{\popxbold\fontsize{11.5}{13}\selectfont\color{popink}{\color{poporange}\thepopsec.\thepopsub}\hspace{6pt}#1}\par\nopagebreak\vspace{4pt}
}

% ============================================================
% Boîtes pédagogiques — même recette : fond blanc, bordure épaisse colorée,
% ombre dure encre, badge plein. Titre optionnel [#1].
% ============================================================
\tcbset{popbase/.style={breakable,enhanced,colback=white,boxrule=2.3pt,arc=9pt,
  shadow={3pt}{-3pt}{0pt}{popink},left=14pt,right=14pt,top=11pt,bottom=11pt,before skip=11pt,after skip=15pt,
  fontupper=\popsemi\fontsize{10.6}{14.5}\selectfont\color{popink},
  fontlower=\popsemi\fontsize{10.6}{14.5}\selectfont\color{popink}}}
% Coche dessinée (le glyphe ✓ n'existe pas dans les polices du document)
\newcommand{\popcheck}[1]{\tikz[baseline=-0.5ex]\draw[#1,line width=1.6pt,line cap=round,line join=round](-0.13,0.0)--(-0.04,-0.09)--(0.13,0.1);}
\providecommand{\coche}{\popcheck{popgreen}}
\providecommand{\resultat}[1]{\par\smallskip\noindent\fcolorbox{popgreen}{popgreenL}{\parbox{\dimexpr\linewidth-2\fboxsep-2\fboxrule\relax}{\textbf{Result.} #1}}\par\smallskip}
\newcommand{\popboxhead}[3]{\poppill{#1}{#2}{white}\ifx\relax#3\relax\else\hspace{6pt}{\popxbold\fontsize{10.6}{12}\selectfont\color{popink}#3}\fi\par\vspace{7pt}}

\newtcolorbox{aretenir}[1][]{popbase,colframe=popgreen,before upper={\popboxhead{Key points}{popgreen}{#1}}}
\newtcolorbox{exemplebox}[1][]{popbase,colframe=poporange,before upper={\popboxhead{Exemple}{poporange}{#1}}}
\newtcolorbox{definition}[1][]{popbase,colframe=popblue,before upper={\popboxhead{Definition}{popblue}{#1}}}
\newtcolorbox{propriete}[1][]{popbase,colframe=poppurple,before upper={\popboxhead{Property}{poppurple}{#1}}}
\newtcolorbox{methode}[1][]{popbase,colframe=popgold,before upper={\popboxhead{Method}{popgold}{#1}}}
\newtcolorbox{attention}[1][]{popbase,colframe=poppink,before upper={\popboxhead{Warning}{poppink}{#1}}}
\newtcolorbox{experience}[1][]{popbase,colframe=popblue,colback=popblueL,before upper={\popboxhead{Experiment}{popblue}{#1}}}
% « Le savais-tu ? » : carte violette pleine (contexte, culture, Cameroun)
\newtcolorbox{savaistu}[1][]{popbase,colback=poppurple,colframe=popink,
  fontupper=\popsemi\fontsize{10.4}{14.2}\selectfont\color{white},
  before upper={\poppill{Did you know?}{white}{poppurple}\ifx\relax#1\relax\else\hspace{6pt}{\popxbold\color{white}#1}\fi\par\vspace{7pt}}}
% Exercice résolu : énoncé en haut, \tcblower puis la solution sur fond crème
\newcounter{popexo}\newcounter{popexr}
\newtcolorbox{exoresolu}[1][]{popbase,colframe=poporange,colbacklower=popcream,
  segmentation style={popink,line width=1.2pt,dashed},
  before upper={\stepcounter{popexr}\popboxhead{Worked example \thepopexr}{poporange}{#1}},
  before lower={\poppill{Solution}{popink}{popcream}\par\vspace{6pt}}}
% Exercice (non résolu en place) — la correction est dans la partie Corrigés
\newtcolorbox{exercice}[1][]{popbase,colframe=popink,
  before upper={\stepcounter{popexo}\popboxhead{Exercise \thepopexo}{popink}{#1}}}
% Corrigé
\newtcolorbox{corrige}[1][]{popbase,colframe=popgreen,colback=popgreenL,
  before upper={\popboxhead{Solution}{popgreen}{#1}}}
% Objectifs / prérequis / bilan
\newtcolorbox{objectifs}{popbase,colframe=popink,colback=poporangeL,
  before upper={\popboxhead{Chapter objectives}{popink}{}}}
\newtcolorbox{prerequis}{popbase,colframe=popink,colback=popblueL,
  before upper={\popboxhead{Prerequisites}{popblue}{}}}
\newtcolorbox{bilan}{popbase,colframe=popink,colback=white,boxrule=2.8pt,
  before upper={\popboxhead{Fiche bilan}{poporange}{L'essentiel à connaître pour l'examen}}}
% Figure encadrée avec légende
\newtcolorbox{popfigure}[1][]{enhanced,breakable=false,colback=white,colframe=popline,boxrule=1.4pt,arc=8pt,
  left=10pt,right=10pt,top=10pt,bottom=8pt,before skip=10pt,after skip=12pt,halign=center,
  lower separated=false,halign lower=center,
  fontlower=\popsemi\fontsize{9}{11.5}\selectfont\color{popmuted},
  #1}
\newcommand{\legende}[1]{\par\vspace{6pt}{\popsemi\fontsize{9}{11.5}\selectfont\color{popmuted}#1}}

% ============================================================
% Signature OnBuch+
% ============================================================
\newcommand{\poplogoinline}[1]{{\poplogo\fontsize{#1}{#1}\selectfont\color{popink}OnBuch\color{poporange}+}}
\newcommand{\signatureonbuch}{%
  \begin{tcolorbox}[enhanced,colback=popink,colframe=popink,boxrule=2.2pt,arc=10pt,
      drop shadow=poporange,left=14pt,right=14pt,top=10pt,bottom=10pt,before skip=0pt,after skip=0pt]
    \tikz[baseline=-0.7ex]\node[circle,fill=popgreen,draw=popcream,line width=1.3pt,minimum size=22pt,inner sep=0]{\popcheck{white}};%
    \hspace{9pt}\begin{minipage}[c]{0.62\linewidth}
      {\popxbold\fontsize{12}{13}\selectfont\color{popcream}Signed\ \textbullet\ {\poplogo OnBuch\color{poporange}+}}\par\vspace{2pt}
      {\raggedright\popsemi\fontsize{8}{10.5}\selectfont\color{popline}\MakeUppercase{OnBuch+ teaching team}\\\MakeUppercase{Follows the official MINESEC syllabus}\par}
    \end{minipage}\hfill
    {\poplogo\fontsize{22}{22}\selectfont\color{popcream}OnBuch\color{poporange}+}
  \end{tcolorbox}%
}

% ============================================================
% Page de garde
% #1 titre · #2 matière · #3 niveau · #4 module · #5 n° leçon · #6 durée ·
% #7 description / objectifs (texte ou itemize)
% ============================================================
\newcommand{\pagegarde}[7]{%
  \thispagestyle{empty}
  \newgeometry{a4paper,margin=1.7cm,top=1.6cm,bottom=1.4cm}%
  \begin{tikzpicture}[remember picture,overlay]
    \fill[poporange!18] ([xshift=-3.2cm,yshift=-3.6cm]current page.north east) circle (4.6cm);
    \fill[poppurple!14] ([xshift=2.4cm,yshift=7.6cm]current page.south west) circle (3.8cm);
  \end{tikzpicture}%
  {\poplogo\fontsize{30}{30}\selectfont\color{popink}OnBuch\color{poporange}+}\hfill
  \raisebox{0.6ex}{\poppill{Signed course \textbullet\ Premium}{popink}{popcream}}\par
  \vspace{1pt}\popsquiggle{poppurple}\par
  \vspace{8pt}
  \begin{tcolorbox}[enhanced,colback=poporange,colframe=popink,boxrule=2.8pt,arc=13pt,
      drop shadow=popink,left=18pt,right=18pt,top=12pt,bottom=13pt,after skip=10pt]
    \poppill{#2 \textbullet\ #3}{popink}{popcream}\hspace{5pt}\poppill{#4}{white}{popink}\par\vspace{8pt}
    {\popdisplay\fontsize{14}{14}\selectfont\color{popink}CHAPTER #5}\par\vspace{4pt}
    {\raggedright\hyphenpenalty=10000\exhyphenpenalty=10000\popdisplay\fontsize{25}{27}\selectfont\color{popcream}\MakeUppercase{#1}\par}
    \vspace{8pt}
    {\popxbold\fontsize{9.5}{9.5}\selectfont\color{popink}\MakeUppercase{Suggested time: #6}}
  \end{tcolorbox}
  \begin{tcolorbox}[enhanced,colback=white,colframe=popink,boxrule=2.2pt,arc=10pt,
      drop shadow=popink,left=16pt,right=16pt,top=10pt,bottom=10pt,after skip=8pt]
    \poppill{In this chapter}{poppurple}{white}\par\vspace{5pt}
    \popsemi\fontsize{9.3}{12.6}\selectfont\color{popink}#7
  \end{tcolorbox}
  \noindent\begin{minipage}[t]{0.487\linewidth}
    \begin{tcolorbox}[enhanced,colback=popgreen,colframe=popink,boxrule=2.2pt,arc=10pt,
        drop shadow=popink,left=11pt,right=11pt,top=10pt,bottom=10pt,height=3.1cm,valign=center]
      \tcbox[colback=white,colframe=popink,boxrule=1.3pt,arc=5pt,boxsep=0pt,nobeforeafter,
        left=3pt,right=3pt,top=3pt,bottom=3pt,valign=center]{\qrcode[height=1.9cm]{https://play.google.com/store/apps/details?id=com.luvvix.onbuch&hl=en}}%
      \hspace{8pt}\begin{minipage}[c]{0.5\linewidth}
        {\popdisplay\fontsize{11}{12.5}\selectfont\color{white}\MakeUppercase{Download OnBuch}}\par\vspace{4pt}
        {\raggedright\popsemi\fontsize{8.4}{11}\selectfont\color{white}Free \textbullet\ Google Play\\AI tutor, past papers, results\par}
      \end{minipage}
    \end{tcolorbox}
  \end{minipage}\hfill
  \begin{minipage}[t]{0.487\linewidth}
    \begin{tcolorbox}[enhanced,colback=poppurple,colframe=popink,boxrule=2.2pt,arc=10pt,
        drop shadow=popink,left=11pt,right=11pt,top=10pt,bottom=10pt,height=3.1cm,valign=center]
      \tikz[baseline=-0.7ex]\node[circle,fill=white,draw=popink,line width=1.3pt,minimum size=1.6cm,inner sep=0]{\popdisplay\fontsize{24}{24}\selectfont\color{poppurple}+};%
      \hspace{8pt}\begin{minipage}[c]{0.6\linewidth}
        {\popdisplay\fontsize{11}{12.5}\selectfont\color{white}\MakeUppercase{Go OnBuch+}}\par\vspace{4pt}
        {\raggedright\popsemi\fontsize{8.4}{11}\selectfont\color{white}All signed courses, unlimited AI tutor, PDF sheets — OnBuch+ tab in the app\par}
      \end{minipage}
    \end{tcolorbox}
  \end{minipage}
  \vspace{8pt}
  \popcontenu
  \vfill
  \signatureonbuch
  \restoregeometry
  \newpage
}

% Rangée « Ce cours contient » (page de garde)
\newcommand{\poptile}[3]{% #1 couleur #2 gros texte #3 légende
  \begin{tcolorbox}[enhanced,colback=#1,colframe=popink,boxrule=2pt,arc=9pt,shadow={2.5pt}{-2.5pt}{0pt}{popink},
      left=6pt,right=6pt,top=9pt,bottom=9pt,halign=center,valign=center,height=1.75cm,width=\linewidth,nobeforeafter]
    {\popdisplay\fontsize{12.5}{13}\selectfont\color{white}\MakeUppercase{#2}}\par\vspace{3pt}
    {\centering\popsemi\fontsize{7.8}{9.5}\selectfont\color{white}#3\par}
  \end{tcolorbox}}
\newcommand{\popcontenu}{%
  \par\noindent{\popxboldls\fontsize{7.5}{7.5}\selectfont\color{popmuted}THIS SIGNED COURSE CONTAINS}\par\vspace{7pt}
  \noindent\begin{minipage}[t]{0.235\linewidth}\poptile{popblue}{Course}{complete, illustrated, step by step}\end{minipage}\hfill
  \begin{minipage}[t]{0.235\linewidth}\poptile{poporange}{Methods}{and worked examples}\end{minipage}\hfill
  \begin{minipage}[t]{0.235\linewidth}\poptile{poppink}{Exercises}{exam style + solutions}\end{minipage}\hfill
  \begin{minipage}[t]{0.235\linewidth}\poptile{popgreen}{Summary sheet}{the essentials to remember}\end{minipage}\par}

% Sommaire
\newtcolorbox{sommaire}{enhanced,colback=white,colframe=popink,boxrule=2.2pt,arc=10pt,
  drop shadow=popink,left=18pt,right=18pt,top=13pt,bottom=13pt,before skip=6pt,after skip=20pt,
  before upper={\poppill{Contents}{poppurple}{white}\par\vspace{9pt}\popsemi\fontsize{10.6}{14.5}\selectfont\color{popink}}}

% Signature de fin de cours — poussée tout en bas de la dernière page
\newcommand{\findecours}{%
  \par\vfill\nopagebreak
  \begin{center}\popsquiggle{poporange}\end{center}\vspace{4pt}
  \signatureonbuch
}
\AtBeginDocument{\def\llbracket{\mathopen{[\mkern-3.5mu[}}\def\rrbracket{\mathclose{]\mkern-3.5mu]}}} % intervalles d'entiers (glyphe ⟦ absent de la police)
% Glyphes absents de Fira Math : redéfinis à partir de glyphes disponibles
\AtBeginDocument{%
  \def\setminus{\mathbin{\backslash}}\let\smallsetminus\setminus
  \def\vdots{\mathord{\vbox{\baselineskip3.2pt\lineskiplimit0pt\kern4pt\hbox{.}\hbox{.}\hbox{.}}}}%
  \def\ddots{\mathinner{\mkern1mu\raise6.5pt\hbox{.}\mkern2mu\raise3.6pt\hbox{.}\mkern2mu\raise0.7pt\hbox{.}\mkern1mu}}%
  \def\triangle{\mathord{\scalebox{0.9}{$\symup{\Delta}$}}}%
  \def\nabla{\mathord{\raisebox{\depth}{\scalebox{1}[-1]{$\symup{\Delta}$}}}}%
  \def\checkmark{\popcheck{popdark}}%
  \def\star{\mathbin{\ast}}\def\bigstar{\mathord{\ast}}%
  \def\diamond{\mathbin{\scalebox{0.6}{$\Diamond$}}}%
  \def\bigcup{\mathop{\mathchoice{\raisebox{-0.3em}{\scalebox{1.7}{$\cup$}}}{\scalebox{1.25}{$\cup$}}{\cup}{\cup}}\displaylimits}%
  \def\bigcap{\mathop{\mathchoice{\raisebox{-0.3em}{\scalebox{1.7}{$\cap$}}}{\scalebox{1.25}{$\cap$}}{\cap}{\cap}}\displaylimits}%
  \def\lesseqqgtr{\mathrel{\lessgtr}}\def\bigoplus{\mathop{\oplus}\displaylimits}%
}
\providecommand{\ssi}{if and only if} % abréviation fréquente
\AtBeginDocument{\providecommand{\wideparen}{\overparen}} % arc de cercle

% Fonctions usuelles que les modèles utilisent sans que amsmath les fournisse
\providecommand{\arsinh}{\operatorname{arsinh}}\providecommand{\arcosh}{\operatorname{arcosh}}
\providecommand{\artanh}{\operatorname{artanh}}\providecommand{\arsech}{\operatorname{arsech}}
\providecommand{\arcsch}{\operatorname{arcsch}}\providecommand{\arcoth}{\operatorname{arcoth}}
\providecommand{\arcsinh}{\operatorname{arsinh}}\providecommand{\arccosh}{\operatorname{arcosh}}
\providecommand{\arctanh}{\operatorname{artanh}}\providecommand{\sech}{\operatorname{sech}}
\providecommand{\csch}{\operatorname{csch}}\providecommand{\arcsec}{\operatorname{arcsec}}
\providecommand{\arccsc}{\operatorname{arccsc}}\providecommand{\arccot}{\operatorname{arccot}}
\providecommand{\sgn}{\operatorname{sgn}}\providecommand{\lcm}{\operatorname{lcm}}
\providecommand{\Var}{\operatorname{Var}}\providecommand{\Cov}{\operatorname{Cov}}
\providecommand{\permil}{\ensuremath{{}^{0}\!/_{00}}}\providecommand{\textperthousand}{\ensuremath{{}^{0}\!/_{00}}}

\begin{document}

\pagegarde{Public Finance}{Economics}{Upper Sixth}{Upper Sixth — Economics}{5}{13 periods}{This chapter introduces public finance: the meaning and role of state financial activity, taxation and its types, the calculation of VAT in Cameroon, the structure and economic consequences of taxation, the State budget and programmed budgeting, the public debt, and the effectiveness of fiscal policy in Cameroon.
\vspace{4pt}
\begin{multicols}{2}\small
\begin{itemize}
\item By the end of this chapter I can define public finance and distinguish its main branches (revenue, expenditure, debt).
\item By the end of this chapter I can define a tax, state its characteristics and distinguish taxation from other public revenues.
\item By the end of this chapter I can classify taxes (direct/indirect; proportional, progressive, regressive) with Cameroonian examples.
\item By the end of this chapter I can calculate Value Added Tax (VAT) on a transaction in Cameroon.
\item By the end of this chapter I can describe the structure of taxation in Cameroon (tax base, rates, main taxes).
\item By the end of this chapter I can analyse the economic consequences of taxation on consumption, production, saving and income distribution.
\end{itemize}\end{multicols}}

\begin{sommaire}
\begin{multicols}{2}
\begin{enumerate}
\item Meaning and Scope of Public Finance
\item Taxation: Meaning, Characteristics and Principles
\item Types of Taxes
\item VAT in Cameroon and the Structure of Taxation (Practical Work)
\item Economic Consequences of Taxation
\item The State Budget and the Programmed Budget
\item The Public Debt and the Effectiveness of Fiscal Policy in Cameroon
\item Integration activity
\item Methods and worked examples
\item Exercises
\item Solutions to the exercises
\item Summary sheet
\end{enumerate}
\end{multicols}
\end{sommaire}

\begin{objectifs}
\begin{itemize}
\item By the end of this chapter I can define public finance and distinguish its main branches (revenue, expenditure, debt).
\item By the end of this chapter I can define a tax, state its characteristics and distinguish taxation from other public revenues.
\item By the end of this chapter I can classify taxes (direct/indirect; proportional, progressive, regressive) with Cameroonian examples.
\item By the end of this chapter I can calculate Value Added Tax (VAT) on a transaction in Cameroon.
\item By the end of this chapter I can describe the structure of taxation in Cameroon (tax base, rates, main taxes).
\item By the end of this chapter I can analyse the economic consequences of taxation on consumption, production, saving and income distribution.
\item By the end of this chapter I can explain the preparation, content and principles of the State budget, including the programmed budget.
\item By the end of this chapter I can explain the causes, forms and burden of the public (national) debt.
\item By the end of this chapter I can evaluate the effectiveness of fiscal policy in Cameroon.
\end{itemize}
\end{objectifs}

\begin{prerequis}
\begin{itemize}
\item The role of the State in the economy (Lower Sixth: economic agents and the public sector).
\item National income and GDP concepts (Lower Sixth: national accounting).
\item Demand, supply and price formation (Lower Sixth: market mechanisms).
\item Basic percentage and proportion calculations.
\item Macroeconomic objectives: growth, full employment, price stability, external balance.
\end{itemize}
\end{prerequis}

\newpage

\coursec{Meaning and Scope of Public Finance}

\courssub{Introduction: Where Does the Money Go?}

Every morning, a trader at Mokolo Market in Yaoundé pays market tolls before she can open her stall. Her daughter attends a government secondary school built with state money, and the tarred road she uses to transport her goods was financed partly by loans contracted by the government. Yet when VAT raised the price of her cooking oil, she wondered: \emph{where does all this money go, and who decides?} In recent years, Cameroon's annual state budget has exceeded 6,000 billion FCFA, while the public debt has stood at around 45\% of GDP. How does the State collect, spend and borrow money --- and with what consequences for citizens and businesses?

This chapter opens the doors of \cle{public finance} to answer these questions. In this first section, we define the field, distinguish it from private finance, identify its components and sources of revenue, and explain why the State intervenes in the economy through its finances.

\courssub{Definition of Public Finance}

When a household manages its money, it decides how much to earn, spend and save. The State does exactly the same thing, but on a national scale and with a fundamental difference: it can \emph{compel} individuals and firms to contribute through taxation. The study of these financial operations of the State is called public finance.

\begin{definition}[Public Finance]
\cle{Public finance} is the branch of economics that studies the \textbf{revenue}, \textbf{expenditure} and \textbf{borrowing} of the State and other public authorities (regions, councils, public enterprises), and the effects of these operations on the economy and on the welfare of citizens.
\end{definition}

Public finance therefore asks three great questions:
\begin{itemize}
\item \textbf{Where does public money come from?} (taxation, fees, state property, grants, borrowing);
\item \textbf{How is it spent?} (salaries of civil servants, roads, schools, hospitals, defence);
\item \textbf{With what effects?} (on production, employment, prices and the distribution of income).
\end{itemize}

\begin{attention}[Public Finance $\neq$ Public Funds]
Do not confuse \textbf{public finance} (the \emph{field of study} --- the science of the State's financial operations) with \textbf{public funds} (the \emph{money itself} belonging to the State). Saying ``public finance is mismanaged'' is imprecise; one should say ``public \emph{funds} are mismanaged''. In the GCE examination, this confusion costs marks in definition questions.
\end{attention}

\courssub{Public Finance versus Private Finance}

A useful way to understand public finance is to compare it with the finance of a household or a firm.

\begin{center}
\begin{tabularx}{\linewidth}{|l|X|X|}
\hline
\rowcolor{popblueL}
\textbf{Criterion} & \textbf{Private finance} & \textbf{Public finance} \\
\hline
Objective & Maximise private profit or satisfaction & Maximise social welfare and the general interest \\
\hline
Source of revenue & Sales, wages, gifts --- obtained \emph{voluntarily} & Mainly taxes, collected by \emph{coercion} (legal compulsion) \\
\hline
Constraint & Must balance income and spending; persistent deficit leads to bankruptcy & Can run deficits and borrow; can even influence the whole economy \\
\hline
Time horizon & Short or medium term & Long term (roads, dams, education of a whole generation) \\
\hline
Beneficiary & The individual or the firm & The community as a whole, including non-payers (public goods) \\
\hline
\end{tabularx}
\end{center}

The key word is \cle{coercion}: a trader cannot refuse to pay VAT simply because she dislikes it, whereas no customer can be forced to buy her tomatoes. This coercive power is what allows the State to finance goods --- security, justice, street lighting --- that no private firm would profitably provide.

\courssub{The Components of Public Finance}

Public finance has four interrelated components, which form a circuit:

\begin{definition}[Public Revenue and Public Expenditure]
\cle{Public revenue} is the totality of financial resources obtained by the State and other public authorities during a given period (usually one year): taxes, fees and dues, income from state property, grants and loans.

\cle{Public expenditure} is the totality of spending carried out by the State and other public authorities to fulfil their missions: payment of civil servants, purchase of goods and services, investment (roads, schools, hospitals), subsidies and debt servicing.
\end{definition}

The two other components are:
\begin{itemize}
\item \textbf{The budget}: the annual document that \emph{forecasts and authorises} revenue and expenditure. When expenditure exceeds revenue, there is a \cle{budget deficit}; the opposite case is a \cle{budget surplus}; when the two are equal, the budget is \emph{balanced}.
\item \textbf{The public debt}: the accumulation of past borrowing that the State must repay, with interest. Deficits are normally financed by borrowing, which feeds the debt.
\end{itemize}

\begin{popfigure}
\begin{center}
\begin{tikzpicture}[
box/.style={draw=popblue, thick, rounded corners, fill=popblueL, text width=3.1cm, align=center, minimum height=1.1cm, font=\small},
arr/.style={-{Stealth[length=3mm]}, very thick, popdark}]
\node[box, align=center] (rev) at (0,0){\textbf{Public revenue}\\ taxes, fees, grants, loans};
\node[box, align=center] (bud) at (5.6,0){\textbf{The Budget}\\ forecast \& authorisation};
\node[box, align=center] (exp) at (11.2,0){\textbf{Public expenditure}\\ salaries, investment, subsidies};
\node[box, fill=popgreenL, draw=popgreen, align=center] (eco) at (5.6,-3.2){\textbf{The Economy}\\ output, jobs, prices, equity};
\node[box, fill=white, draw=poporange, align=center] (debt) at (0,-3.2){\textbf{Public debt}\\ borrowing to cover deficits};
\draw[arr] (rev) -- (bud);
\draw[arr] (bud) -- (exp);
\draw[arr] (exp.south) |- (eco.east);
\draw[arr] (eco.west) -| (rev.south) node[pos=0.25, below, font=\scriptsize, align=center, text width=2.6cm]{taxes on income \& activity};
\draw[arr, dashed, poporange] (debt.north) -- (rev.south west) node[midway, left, font=\scriptsize]{loans};
\draw[arr, dashed, poporange] (eco.south) .. controls (2,-4.6) .. (debt.south) node[midway, below, font=\scriptsize, align=center, text width=3.4cm]{deficit $\Rightarrow$ new borrowing};
\end{tikzpicture}
\end{center}
\legende{The circuit of public finance: revenue is authorised by the budget, spent on the economy, which in turn generates the tax base --- while deficits are financed by borrowing.}
\end{popfigure}

\courssub{Sources of Public Revenue}

The State draws its resources from five main sources:

\begin{enumerate}
\item \textbf{Taxation}: compulsory levies on income, profits, property and consumption (income tax, company tax, VAT, customs duties). This is by far the largest source of revenue in Cameroon.
\item \textbf{Fees and dues}: payments for specific services rendered by the State --- market tolls, court fees, passport fees, university registration fees. Unlike taxes, they are paid only by users of the service, in exchange for that service.
\item \textbf{Income from state property}: rents, dividends from state-owned or state-participated enterprises, and royalties from natural resources (for example oil revenue channelled through the national oil company).
\item \textbf{Grants and donations}: non-repayable transfers from foreign governments and international organisations, often tied to specific projects (health, roads).
\item \textbf{Borrowing}: issuing treasury bills and bonds, or contracting loans from banks, foreign states and institutions. Borrowing is revenue \emph{today} but creates a repayment obligation \emph{tomorrow}.
\end{enumerate}

\begin{attention}[Tax or Fee?]
A \textbf{tax} is paid by all who fall within its legal scope, whether or not they receive a particular service in return; a \textbf{fee} is paid only by the user of an identified public service. Examination questions often test this distinction: a passport fee is a fee, but income tax is a tax even though the taxpayer also uses roads and schools.
\end{attention}

\begin{popfigure}
\begin{center}
\begin{tikzpicture}
\begin{axis}[
ybar, bar width=0.9cm, width=11cm, height=6.2cm,
ymin=0, ymax=100,
ylabel={Share of total revenue (\%)},
symbolic x coords={Tax revenue, Non-tax revenue},
xtick=data,
nodes near coords, nodes near coords align={vertical},
every node near coord/.append style={font=\small},
ymajorgrids, grid style={popline, dashed},
]
\addplot[fill=popblue, draw=popdark] coordinates {(Tax revenue,82) (Non-tax revenue,18)};
\end{axis}
\end{tikzpicture}
\end{center}
\legende{Indicative composition of Cameroon's public revenue in a recent year: taxation (domestic taxes plus customs duties) provides the overwhelming bulk of resources, while non-tax revenue (fees, dividends, grants, oil royalties) plays a smaller role. Values are illustrative orders of magnitude.}
\end{popfigure}

\begin{exemplebox}[A Council's Revenue in the North-West Region]
A small urban council collects market tolls and business licences (fees and dues), receives its share of national taxes transferred by the central government, obtains a grant from a development partner to rehabilitate a water scheme, and contracts a small bank loan to buy a grader. All five sources appear --- but only the loan must be repaid.
\end{exemplebox}

\courssub{The Objectives of Public Finance: Musgrave's Three Functions}

Why does the State collect and spend at all? The economist Richard Musgrave summarised the answer in three fundamental functions.

\begin{propriete}[Musgrave's Three Functions of the State]
According to \textbf{Musgrave}, the State uses public finance to perform three functions:
\begin{enumerate}
\item \textbf{Allocation of resources}: providing public goods (defence, justice, roads) and correcting market failures (externalities, monopolies) so that resources are used efficiently;
\item \textbf{Distribution (redistribution) of income}: reducing inequalities through progressive taxation and transfers (subsidies, bursaries, social assistance);
\item \textbf{Stabilisation of the economy}: using the budget to fight unemployment and inflation --- spending more (or taxing less) in a recession, and the reverse in a boom.
\end{enumerate}
This classification is a standard reference at the GCE Advanced Level; it should be known with a short explanation of each function.
\end{propriete}

\begin{exemplebox}[The Three Functions in One Road Project]
The construction of the Yaoundé--Nsimalen motorway illustrates all three functions at once: it \emph{allocates} resources to a public good no private firm would finance alone; it \emph{redistributes} by connecting poorer rural producers to urban markets; and by employing thousands of workers during its construction, it contributes to \emph{stabilisation} by sustaining aggregate demand.
\end{exemplebox}

\courssub{Public Finance in Developing Countries: The Case of Cameroon}

In developing countries, public finance carries an additional mission: \textbf{financing development}. Private capital is scarce, so the State must directly finance:
\begin{itemize}
\item \textbf{Infrastructure}: roads, ports, dams (such as hydroelectric projects on the Sanaga river), which lower production costs for the whole economy;
\item \textbf{Education and health}: building schools and hospitals and paying teachers and nurses --- investments in \cle{human capital} whose returns appear only after many years;
\item \textbf{Support to key sectors}: subsidies and credit to agriculture (cocoa, cotton, food crops) that sustain rural incomes and exports.
\end{itemize}

Cameroon faces the classic dilemma of developing countries: \emph{development needs are enormous, but the tax base is narrow}. A large informal sector escapes taxation, so the State relies heavily on VAT, customs duties and oil-related revenue, and resorts to borrowing --- which explains why debt sustainability has become a central policy concern.

\begin{savaistu}[Did You Know?]
Cameroon belongs to the CEMAC zone, where member states have adopted convergence criteria --- including limits on budget deficits and public debt --- to protect the stability of the common currency, the CFA franc. This is why debates about the size of Cameroon's deficit are not merely national: they concern the whole monetary union.
\end{savaistu}

\begin{exoresolu}[Identifying the Components and Functions of Public Finance]
\textbf{Statement.} The Government of Cameroon (i) raises the rate of company tax on large firms; (ii) builds 500 new classrooms in rural areas; (iii) pays interest on bonds issued five years earlier; (iv) grants bursaries to students from low-income families. For each operation, state (a) the component of public finance concerned and (b) the Musgrave function it mainly serves.
\tcblower
\textbf{Solution.}
\begin{itemize}
\item (i) \emph{Public revenue} (taxation). Function: \textbf{redistribution} (and stabilisation), since profitable firms contribute more, freeing resources for social spending.
\item (ii) \emph{Public expenditure} (investment). Function: \textbf{allocation} --- providing a public service (education) that markets under-provide in rural areas.
\item (iii) \emph{Public debt servicing} (a charge arising from past borrowing). It is an expenditure, but it serves no Musgrave function directly; it is the \emph{cost of earlier deficits}.
\item (iv) \emph{Public expenditure} (transfer). Function: \textbf{redistribution} of income towards poorer households.
\end{itemize}
\textbf{Examination tip:} always answer in two steps --- component first, function second --- and justify with one sentence, as above.
\end{exoresolu}

\begin{aretenir}[Key Points of the Section]
\begin{itemize}
\item Public finance studies the \textbf{revenue, expenditure and borrowing} of the State and public authorities; it must not be confused with public funds.
\item Unlike private finance, public finance pursues the \textbf{general interest} and relies on \textbf{coercion} (compulsory taxation).
\item Its four components --- revenue, expenditure, budget, debt --- form a circuit linking the State to the economy.
\item Revenue comes from taxation, fees and dues, state property, grants and borrowing; in Cameroon, taxation dominates.
\item Musgrave's three functions --- \textbf{allocation, distribution, stabilisation} --- explain why the State intervenes; in developing countries, public finance additionally finances \textbf{development} (infrastructure, education, health).
\end{itemize}
\end{aretenir}

\begin{exercice}[Check Your Understanding]
\begin{enumerate}
\item Define public finance and state two differences between public and private finance.
\item Classify the following as tax, fee, income from property, grant or borrowing: a passport fee; a dividend paid by a state enterprise; a loan from the World Bank; VAT on cooking oil; a donation from a foreign government.
\item A government cuts taxes during a recession to stimulate demand. Which of Musgrave's functions is being used? Justify.
\item Explain why the large informal sector makes the financing of development difficult in Cameroon.
\end{enumerate}
\end{exercice}

\begin{corrige}[Exercise 1]
\begin{enumerate}
\item Public finance is the branch of economics dealing with the revenue, expenditure and borrowing of the State and public authorities. Differences (any two): objective (general interest vs private profit); coercion (compulsory taxes vs voluntary transactions); ability to run deficits; longer time horizon.
\item Passport fee: \emph{fee}; dividend: \emph{income from state property}; World Bank loan: \emph{borrowing}; VAT: \emph{tax}; donation: \emph{grant}.
\item \textbf{Stabilisation}: the tax cut raises households' disposable income and firms' after-tax profits, sustaining consumption, investment and aggregate demand during the downturn.
\item Informal activities are not declared and therefore escape taxation; the tax base is narrow, so revenue is insufficient for infrastructure, education and health, forcing the State to rely on a few taxes (VAT, customs) and on borrowing, which increases the public debt.
\end{enumerate}
\end{corrige}

\coursec{Taxation: Meaning, Characteristics and Principles}

In the previous section, we saw that public finance is the branch of economics that studies how the State raises resources (public revenue) and how it spends them (public expenditure) in order to achieve its economic and social objectives. Among all the sources of public revenue — taxes, fees, fines, borrowing, aid — one dominates by far in every modern economy, and Cameroon is no exception: \cle{taxation}. When the State of Cameroon builds the Yaoundé--Douala motorway, pays teachers in Government High Schools, or equips district hospitals, most of the money comes from taxes paid by households and businesses. But what exactly is a tax? What distinguishes it from the price you pay for a passport or from a fine for a traffic offence? And what makes a tax system ``good''? These are the questions of this section.

\courssub{Definition of Taxation and of a Tax}

Imagine a trader at Mokolo market in Yaoundé who sells tomatoes. Every year, she must pay an amount of money to the State, calculated according to rules fixed by law. She receives nothing specific in exchange: the State does not deliver tomatoes, a stall, or a personal service to her in return for that payment. Yet if she refuses to pay, she can be prosecuted and punished. This compulsory, unrequited payment is precisely what economists call a \cle{tax}.

\begin{definition}[Taxation and tax]
\cle{Taxation} is the system by which the State compulsorily collects levies (taxes) from individuals and businesses in order to finance public expenditure and to influence economic activity.

A \cle{tax} (French: \emph{impôt}) is a \textbf{compulsory}, \textbf{unrequited} (without direct counterpart) \textbf{pecuniary levy} imposed by the State (or another public authority) on persons and businesses, by virtue of its sovereign power and according to rules established by law.
\end{definition}

Taxation therefore serves two great purposes:
\begin{itemize}
\item a \textbf{fiscal (budgetary) purpose}: to raise revenue to finance public expenditure — this is the primary and historical purpose;
\item an \textbf{extra-fiscal (economic and social) purpose}: to influence behaviour and the economy, for example taxing cigarettes heavily to discourage smoking, or granting tax reductions to companies that invest in rural areas.
\end{itemize}

\courssub{Essential Characteristics of a Tax}

The definition above contains three essential characteristics. A levy is a genuine tax only if all three are present.

\begin{enumerate}
\item \textbf{Compulsory (obligatory) character.} A tax is imposed by the sovereign authority of the State; the taxpayer cannot choose whether or not to pay. Refusal to pay is sanctioned by penalties (surcharges, seizure of property, prosecution). This distinguishes a tax from a voluntary contribution or a donation.
\item \textbf{Absence of direct counterpart (unrequited character).} The taxpayer does not receive, in exchange for the tax, an individualised and proportionate service from the State. The trader who pays tax cannot demand that the State repair the road in front of her shop \emph{because} she paid. Taxes finance public services for the community as a whole, not for each contributor individually.
\item \textbf{Pecuniary character and legal basis.} A tax is paid \textbf{in money} (not in goods or labour) and is \textbf{established by law}: in Cameroon, no tax can be created, modified or collected without a legal text — principally the annual Finance Law and the General Tax Code. This is the principle of \emph{legality of taxation}: ``no taxation without law''.
\end{enumerate}

\begin{attention}[A tax has no direct counterpart]
The most frequent examination mistake is to confuse a tax with the \textbf{price of a public service}. When you pay for a passport, a hospital consultation in a public hospital, or a water connection, you receive a \emph{specific, individual service} in exchange: these are \textbf{fees}, not taxes. A tax finances the community generally; a fee pays for a service rendered to you personally. Never write in an essay that ``the taxpayer receives services in exchange for his tax'' as if it were a purchase.
\end{attention}

\courssub{Tax, Fee and Fine: Three Different Levies}

The State collects many compulsory payments, but they do not all have the same nature. It is essential to distinguish three of them.

\begin{definition}[Fee and fine]
A \cle{fee} (French: \emph{redevance}) is a payment made to a public authority \textbf{in exchange for a service rendered} to the person who pays (e.g.\ passport fees, university registration fees, television licence fee).

A \cle{fine} (French: \emph{amende}) is a pecuniary \textbf{punishment} imposed for the violation of a law or regulation (e.g.\ a fine for dangerous driving or for illegal dumping of waste).
\end{definition}

\begin{center}
\begin{tabularx}{\linewidth}{|l|X|X|X|}
\hline
\rowcolor{popblueL} \textbf{Criterion} & \textbf{Tax (impôt)} & \textbf{Fee (redevance)} & \textbf{Fine (amende)} \\
\hline
Purpose & Finance public expenditure & Pay for a service rendered & Punish an offence \\
\hline
Counterpart & None (unrequited) & Direct, individual service & None \\
\hline
Trigger & Legal obligation (income, sales, property\dots) & Voluntary use of the service & Violation of the law \\
\hline
Sanction if unpaid & Penalties, prosecution & Service simply not provided & Aggravated penalties \\
\hline
Example in Cameroon & Company tax, VAT, income tax & Passport fee, court registry fee & Traffic offence fine \\
\hline
\end{tabularx}
\end{center}

\begin{popfigure}
\begin{tikzpicture}[font=\small, node distance=0.9cm,
box/.style={draw=popblue, thick, rounded corners, fill=popblueL, text width=3.4cm, align=center, minimum height=1.1cm},
crit/.style={draw=popmuted, fill=white, text width=2.6cm, align=center, minimum height=0.9cm}]
\node[box, align=center] (tax){\textbf{TAX}\\ compulsory, no counterpart};
\node[box, right=1.1cm of tax, align=center] (fee){\textbf{FEE}\\ payment for a service};
\node[box, right=1.1cm of fee, align=center] (fine){\textbf{FINE}\\ punishment of an offence};
\node[crit, below=of tax] (t1) {Finances the budget};
\node[crit, below=of fee] (f1) {Covers cost of service};
\node[crit, below=of fine] (s1) {Deters illegal behaviour};
\draw[->, thick, popdark] (tax) -- (t1);
\draw[->, thick, popdark] (fee) -- (f1);
\draw[->, thick, popdark] (fine) -- (s1);
\end{tikzpicture}
\legende{Tax, fee and fine: three compulsory levies with different economic natures.}
\end{popfigure}

\begin{methode}[How to identify whether a levy is a tax]
To decide whether a given levy is a genuine tax, apply the three characteristics one after the other:
\begin{enumerate}
\item \textbf{Is it compulsory?} If the person can avoid the payment simply by not using a service, it is probably a fee, not a tax.
\item \textbf{Is there a direct, individual counterpart?} If the payer receives a specific service proportionate to the payment, it is a fee. If the payment punishes an offence, it is a fine. A tax has \emph{no} direct counterpart.
\item \textbf{Is it paid in money and established by law?} If yes, and the first two answers point to a tax, conclude: the levy is a \textbf{tax}.
\end{enumerate}
\emph{Example of application:} the annual payment for a television licence is compulsory and fixed by law, but it is linked to the use of the public broadcasting service — many authors classify it as a fee. The company tax, paid on profits with no service in return, is a true tax.
\end{methode}

\courssub{The Principles of a Good Tax System}

Not every way of raising taxes is acceptable. Since the birth of modern economics, thinkers have asked: what makes a tax system \emph{good}? The most famous answer was given by Adam Smith in \emph{The Wealth of Nations} (1776), in the form of four ``canons'' (maxims) of taxation.

\begin{propriete}[Adam Smith's four canons of taxation]
According to Adam Smith, a good tax system must respect four principles (the statement of these canons is examinable; no formal proof is required, but each must be explained):
\begin{enumerate}
\item \textbf{Equity (canon of equality):} taxpayers should contribute according to their \textbf{ability to pay} — those who earn more should pay more. Equity has two dimensions: \emph{horizontal equity} (persons in identical situations pay identical taxes) and \emph{vertical equity} (persons with greater means pay proportionately more).
\item \textbf{Certainty:} the tax each person must pay should be \textbf{clear, fixed and known in advance} (amount, time, manner of payment), not arbitrary. Certainty protects taxpayers from abuse by tax officials.
\item \textbf{Convenience:} taxes should be levied at the \textbf{time and in the manner most convenient} for the taxpayer — for example, income tax withheld monthly from salaries rather than demanded in one lump sum.
\item \textbf{Economy:} the \textbf{cost of collecting} a tax should be as small as possible relative to the revenue it yields; a tax that costs nearly as much to collect as it brings in is a bad tax.
\end{enumerate}
\end{propriete}

Modern economists have added two further principles to Smith's four canons:

\begin{itemize}
\item \textbf{Productivity (yield):} a good tax should bring in \textbf{sufficient revenue} for the State with a broad base and reasonable rates. Very high rates on a narrow base may yield little, because they encourage evasion and discourage the taxed activity.
\item \textbf{Neutrality:} taxation should, as far as possible, \textbf{not distort economic decisions} — choices between working and leisure, saving and consuming, investing in one sector or another should not be twisted by the tax system. In practice, perfect neutrality is impossible (every tax affects behaviour), so governments seek an acceptable compromise between neutrality and the extra-fiscal goals they deliberately pursue (e.g.\ discouraging tobacco).
\end{itemize}

\begin{exemplebox}[Applying the canons to a Cameroonian situation]
Suppose the government wishes to tax small informal traders in Douala.
\begin{itemize}
\item \emph{Equity:} a flat amount identical for a large wholesaler and a small retailer would violate vertical equity; a tax proportional to turnover is fairer.
\item \emph{Certainty:} if market collectors demand arbitrary amounts that change from week to week, the canon of certainty is violated — this is a frequent complaint in practice.
\item \emph{Convenience:} allowing payment by mobile money at the end of the market day respects convenience.
\item \emph{Economy:} if collecting the tax requires an army of agents whose salaries absorb most of the revenue, the canon of economy fails.
\end{itemize}
This simple grid is an excellent tool for the evaluation paragraph of a GCE essay on taxation.
\end{exemplebox}

\courssub{Tax Evasion, Tax Avoidance and Tax Fraud}

Because taxes are compulsory and unrequited, some taxpayers try to escape them. Three behaviours must be carefully distinguished.

\begin{definition}[Tax evasion and tax avoidance]
\cle{Tax evasion} (in its strict, illegal form also called \cle{tax fraud}) consists in \textbf{illegally} escaping tax by violating the law: concealing income, keeping false accounts, failing to declare sales, smuggling goods across the border.

\cle{Tax avoidance} consists in \textbf{legally} reducing one's tax burden by exploiting loopholes, exemptions or favourable provisions of the law — for example, locating an investment in a zone where the law grants a tax holiday, or arranging one's affairs to fall into a lower tax bracket.
\end{definition}

The boundary is legality: avoidance stays within the letter of the law (even if against its spirit), while evasion and fraud break the law and expose the taxpayer to penalties, fines and imprisonment. In everyday language, ``tax evasion'' is often used loosely to cover both ideas; in an examination, always state the distinction precisely.

The \textbf{consequences for public revenue} are serious, especially in developing countries:
\begin{itemize}
\item a \textbf{reduction in budget revenue}, limiting the State's capacity to build roads, schools and hospitals;
\item \textbf{unfair competition}: honest businesses that pay their taxes compete against evaders who do not;
\item a heavier burden on the \textbf{narrow formal sector} (civil servants, registered companies), which cannot easily escape taxation, pushing rates up and discouraging formality;
\item a weakening of the \textbf{civic tax culture}: when evasion is widespread and unpunished, honest taxpayers feel cheated and compliance falls further.
\end{itemize}

\begin{savaistu}[The informal sector and tax revenue in Cameroon]
A large share of economic activity in Cameroon — as in most African economies — takes place in the \textbf{informal sector}: small trade, artisanal production, transport. These activities are difficult to identify and tax, which narrows the tax base considerably. This is one reason why the Cameroonian State relies heavily on taxes that are easier to collect, such as VAT on formal sales and taxes on foreign trade collected at the port of Douala, and why broadening the tax base is a permanent objective of tax policy.
\end{savaistu}

\courssub{The Organisation of Taxation in Cameroon}

In Cameroon, taxation is governed by the \cle{General Tax Code} (\emph{Code Général des Impôts}), which assembles the rules on the different taxes, their bases, rates and procedures, together with the annual \textbf{Finance Law} which fixes the budget and may create or modify taxes. The administration responsible for assessing and collecting taxes is the \cle{Directorate General of Taxation} (DGI — \emph{Direction Générale des Impôts}), under the Ministry of Finance, while customs duties are collected by the Directorate General of Customs.

Every taxpayer — individual or business — is identified by a unique \cle{taxpayer identification number} (\emph{numéro d'identifiant unique}), which appears on a taxpayer card and is required for most administrative and commercial operations (opening a business bank account, bidding for public contracts, importing goods). This identification system is a key weapon against evasion, because it allows the administration to trace a taxpayer's operations.

\begin{popfigure}
\begin{tikzpicture}[font=\small, node distance=0.8cm,
stage/.style={draw=popgreen, thick, rounded corners, fill=popgreenL, text width=4.6cm, align=center, minimum height=1cm}]
\node[stage, align=center] (law){\textbf{1. Tax law}\\ Finance Law \& General Tax Code};
\node[stage, below=of law, align=center] (admin){\textbf{2. Administration}\\ Directorate General of Taxation (DGI)};
\node[stage, below=of admin, align=center] (ident){\textbf{3. Identification}\\ Taxpayer number assigned};
\node[stage, below=of ident, align=center] (decl){\textbf{4. Declaration \& assessment}\\ Taxpayer declares; DGI assesses};
\node[stage, below=of decl, align=center] (coll){\textbf{5. Collection}\\ Payment, controls, penalties};
\node[stage, below=of coll, align=center] (budget){\textbf{6. Public revenue}\\ Finances the State budget};
\draw[->, thick, popdark] (law) -- (admin);
\draw[->, thick, popdark] (admin) -- (ident);
\draw[->, thick, popdark] (ident) -- (decl);
\draw[->, thick, popdark] (decl) -- (coll);
\draw[->, thick, popdark] (coll) -- (budget);
\end{tikzpicture}
\legende{From tax law to tax collection in Cameroon.}
\end{popfigure}

\begin{exoresolu}[Tax, fee or fine?]
\emph{Statement.} Classify each of the following payments collected in Cameroon as a tax, a fee or a fine, justifying your answer with the characteristics studied: (a) the company tax paid by a brewery on its profits; (b) the sum paid to obtain a national identity card; (c) the sum paid by a lorry driver sanctioned for overloading; (d) VAT paid by consumers on purchases in a supermarket.
\tcblower
\emph{Solution.}
\begin{itemize}
\item \textbf{(a) Company tax:} a \textbf{tax}. It is compulsory, established by the General Tax Code, paid in money, and the brewery receives no individual service in exchange.
\item \textbf{(b) Identity card payment:} a \textbf{fee}. The payment is made in exchange for a specific, individual service rendered by the administration (production and delivery of the document). There is a direct counterpart.
\item \textbf{(c) Overloading sanction:} a \textbf{fine}. It is a pecuniary punishment imposed for the violation of a regulation; its purpose is to punish and deter, not to finance ordinary expenditure or pay for a service.
\item \textbf{(d) VAT:} a \textbf{tax}. It is compulsory, fixed by law, paid in money on consumption, with no direct counterpart for the consumer. It is an \emph{indirect} tax, collected by the seller and transferred to the State (studied in detail in the next section).
\end{itemize}
\end{exoresolu}

\begin{exercice}[Checking understanding]
\begin{enumerate}
\item Define a tax and state its three essential characteristics.
\item Explain, with an example, the difference between tax avoidance and tax evasion, and give two consequences of widespread evasion for the Cameroonian economy.
\item A municipality collects a fixed sum from every shopkeeper to finance the daily cleaning of the market, a service each shopkeeper directly benefits from. Is this a tax? Justify using the identification method.
\item State Adam Smith's four canons and use them to evaluate a proposal to tax mobile phone top-ups.
\end{enumerate}
\end{exercice}

\begin{corrige}[Exercise: elements of answers]
\begin{enumerate}
\item A tax is a compulsory, unrequited pecuniary levy imposed by the State by law. Characteristics: compulsory; no direct counterpart; paid in money and established by law.
\item Avoidance is legal (using loopholes, exemptions); evasion/fraud is illegal (concealment, false declarations). Consequences: loss of budget revenue; unfair competition and heavier burden on the formal sector; erosion of tax culture.
\item No: there is a direct, individual counterpart (the cleaning service), so it is a \textbf{fee}, even though it is compulsory.
\item Equity, certainty, convenience, economy; then evaluate: e.g.\ the tax is certain and convenient (collected at purchase), economical (operators collect it), but its equity is debatable if it hits poor users proportionately harder.
\end{enumerate}
\end{corrige}

\begin{aretenir}[Key points]
\begin{itemize}
\item A \textbf{tax} is a compulsory, unrequited, pecuniary levy established by law; taxation is the system of such levies.
\item A \textbf{fee} pays for a service rendered; a \textbf{fine} punishes an offence; only the tax has no direct counterpart.
\item A good tax system respects Adam Smith's canons — \textbf{equity, certainty, convenience, economy} — plus the modern principles of \textbf{productivity} and \textbf{neutrality}.
\item \textbf{Avoidance} is legal, \textbf{evasion/fraud} is illegal; both reduce public revenue, but only fraud is punishable.
\item In Cameroon, taxation is governed by the \textbf{General Tax Code} and administered by the \textbf{DGI}, with each taxpayer identified by a unique \textbf{taxpayer number}.
\end{itemize}
\end{aretenir}

Having defined taxation and its principles, we must now examine the different \emph{forms} that taxes take — direct and indirect taxes, proportional, progressive and regressive taxation — which is the object of the next section.

\coursec{Types of Taxes}

In the previous section we defined taxation as a compulsory levy imposed by the State on persons and businesses, and we examined its characteristics (compulsoriness, absence of direct counterpart, generality) and the principles of a good tax system (equity, certainty, convenience, economy, productivity). But ``tax'' is not a single instrument. When a trader in Douala pays tax on her profits, when a consumer in Yaound\'e pays VAT on a bag of rice, and when an importer pays customs duty on a container at the port, three very different taxes are at work. The purpose of this section is to classify taxes, to understand how each type works, and to evaluate its advantages and disadvantages --- a classic GCE essay question.

\courssub{Direct Taxes and Indirect Taxes}

\textbf{The basic distinction: impact versus incidence.}\par
The oldest and most fundamental classification of taxes rests on the distinction between the \cle{impact} of a tax (the person who legally pays it to the Treasury) and its \cle{incidence} (the person who ultimately bears the economic burden). If the person who pays the tax to the Treasury is the same person who finally bears its burden, the tax is \cle{direct}. If the person who hands the money to the State can pass the burden on to someone else (usually by raising prices), the tax is \cle{indirect}.

\begin{definition}[Direct tax]
A \cle{direct tax} is a tax levied directly on the \textbf{income} or \textbf{wealth} of an identifiable person (individual or company), who is expected to bear the burden personally and cannot normally shift it to someone else. The impact and incidence fall on the same agent. Examples: personal income tax, company tax, property (land) tax.
\end{definition}

\begin{definition}[Indirect tax]
An \cle{indirect tax} is a tax levied on \textbf{expenditure, consumption or transactions} (goods and services, imports). It is collected by an intermediary (seller, importer) who can \textbf{shift (transfer) the burden} to the final consumer through higher prices. The impact is on the intermediary; the incidence is on the consumer. Examples: VAT, excise duties, customs duties.
\end{definition}

\textbf{The main direct taxes in Cameroon.}\par
\begin{itemize}
\item \textbf{Personal income tax (PIT / IRPP):} levied on the earnings of individuals --- salaries, wages, business profits of sole traders. In Cameroon, wages are subject to the \emph{Imp\^ot sur le Revenu des Personnes Physiques} (IRPP), withheld at source by the employer according to progressive brackets.
\item \textbf{Company tax (corporate income tax):} levied on the profits of companies. In Cameroon, limited companies such as SONARA, ENEO or the mobile telephone operators pay company tax on their taxable profits at a proportional rate (currently 30\% + 10\% council tax).
\item \textbf{Property (real estate) tax:} levied on the ownership of land and buildings, usually collected by local councils (decentralised authorities).
\end{itemize}

\textbf{The main indirect taxes in Cameroon.}\par
\begin{itemize}
\item \textbf{Value Added Tax (VAT):} a general tax on consumption, charged at each stage of production and distribution but borne by the final consumer. It is studied in detail in Practical Work 21 of this chapter.
\item \textbf{Excise duties:} selective taxes on particular goods such as alcoholic drinks, tobacco, cigarettes and certain luxury products. They raise revenue and also discourage the consumption of harmful products.
\item \textbf{Customs (import) duties:} taxes on goods entering the country, collected at the port or border. Within CEMAC, a common external tariff (CET) applies to imports from outside the community.
\end{itemize}

\begin{attention}[VAT is an indirect tax -- common trap]
A very common mistake is to say that VAT is ``paid by the shopkeeper, therefore it is a direct tax on the shopkeeper''. This is wrong. The seller merely \emph{collects} VAT on behalf of the State and adds it to the price; the \textbf{final consumer bears the burden}. A tax is classified according to who \emph{bears} it (incidence), not according to who \emph{hands over the money} (impact). VAT is therefore an \cle{indirect tax}.
\end{attention}

\begin{popfigure}
\begin{center}
\begin{tikzpicture}[grow=right, level distance=3.4cm,
  level 1/.style={sibling distance=4.6cm},
  level 2/.style={sibling distance=1.5cm},
  edge from parent/.style={draw=popline, thick},
  every node/.style={align=center, font=\small}]
\node[draw=popdark, fill=popblueL, rounded corners, font=\bfseries] {Taxes}
  child { node[draw=poporange, rounded corners, align=center]{Indirect taxes\\ \emph{(on expenditure)}}
    child { node[draw=popline, rounded corners, fill=white] {Customs duties} }
    child { node[draw=popline, rounded corners, fill=white] {Excise duties} }
    child { node[draw=popline, rounded corners, fill=white] {VAT} }
  }
  child { node[draw=popblue, rounded corners, align=center]{Direct taxes\\ \emph{(on income \& wealth)}}
    child { node[draw=popline, rounded corners, fill=white] {Property tax} }
    child { node[draw=popline, rounded corners, fill=white] {Company tax} }
    child { node[draw=popline, rounded corners, fill=white] {Personal income tax (IRPP)} }
  };
\end{tikzpicture}
\end{center}
\legende{Classification tree of taxes: direct versus indirect, with the main sub-types.}
\end{popfigure}

\textbf{Advantages and disadvantages.}\par
\begin{center}
\begin{tabularx}{\linewidth}{|l|X|X|}
\hline
\rowcolor{popblueL} & \textbf{Advantages} & \textbf{Disadvantages} \\
\hline
\textbf{Direct taxes} & Fair (can be made progressive and adjusted to ability to pay); certain and economical to collect once taxpayers are identified; reduce income inequality; civic awareness (taxpayer feels the burden). & Easy to evade (informal sector, false declarations); high rates may discourage work, saving and investment; require a large, honest tax administration; unpopular. \\
\hline
\textbf{Indirect taxes} & Difficult to evade (built into prices); cheap and convenient to collect at point of sale/entry; wide base, so very productive; can discourage harmful consumption (tobacco, alcohol); less visible to taxpayer. & Regressive: the poor spend a larger share of their income, so they bear a relatively heavier burden; inflationary (raise prices); can distort consumption choices; taxpayer unaware of amount paid. \\
\hline
\end{tabularx}
\end{center}

\begin{savaistu}[Why developing countries rely on indirect taxes]
In Cameroon, as in most African countries, a large part of economic activity is informal and many incomes escape declaration. Direct taxes are therefore hard to collect, while VAT and customs duties are collected automatically at the shop or the port. This is why indirect taxes provide the bulk of tax revenue in Cameroon, whereas rich countries rely more on direct taxes.
\end{savaistu}

\courssub{Proportional, Progressive and Regressive Taxation}

A second classification concerns the relationship between the \cle{tax rate} and the \cle{tax base} (the income, wealth or value on which the tax is calculated).

\begin{definition}[Proportional tax]
A \cle{proportional tax} applies a \textbf{constant rate} to the tax base, whatever its size. Tax paid $=$ constant rate $\times$ base. The average tax rate (ATR) equals the marginal tax rate (MTR) at all levels. Example: a flat tax of $10\%$ on all incomes; VAT at a single rate is proportional to the value of the purchase; company tax in Cameroon is proportional.
\end{definition}

\begin{definition}[Progressive tax]
A \cle{progressive tax} is a tax whose \textbf{rate increases as the base increases}: higher incomes are taxed at higher rates. This is achieved through \textbf{brackets} (tranches): each slice of income is taxed at its own, rising rate. The marginal tax rate exceeds the average tax rate (MTR $>$ ATR). Example: personal income tax (IRPP) in Cameroon.
\end{definition}

\begin{definition}[Regressive (degressive) tax]
A \cle{regressive tax} is a tax whose \textbf{effective rate falls as the base rises}: the rich pay a smaller \emph{proportion} of their income than the poor, even if they pay more in absolute terms. The marginal tax rate is below the average tax rate (MTR $<$ ATR). Indirect taxes on necessities are regressive in effect, because the poor devote a larger share of their income to consumption.
\end{definition}

\begin{methode}[Computing average and marginal tax rates with brackets]
\begin{enumerate}
\item Identify the brackets and their corresponding rates.
\item For a given income, apply each bracket's rate \textbf{only to the part of income falling inside that bracket} (slice by slice).
\item Sum the tax from each slice to get \textbf{total tax paid}.
\item \cle{Average tax rate (ATR)} $= \dfrac{\text{total tax paid}}{\text{total income}} \times 100$.
\item \cle{Marginal tax rate (MTR)} $=$ the rate of the bracket in which the \emph{last franc earned} falls.
\item Compare ATR across income levels: constant ATR $\Rightarrow$ proportional; rising ATR $\Rightarrow$ progressive; falling ATR $\Rightarrow$ regressive.
\end{enumerate}
\end{methode}

\begin{exoresolu}[Progressive income tax with brackets]
Suppose a country taxes income as follows: $0$--$100\,000$ FCFA at $0\%$; $100\,001$--$300\,000$ FCFA at $10\%$; $300\,001$--$500\,000$ FCFA at $20\%$. Compute the tax, the average rate and the marginal rate for Mr Eyang, who earns $400\,000$ FCFA.
\tcblower
\textbf{Solution.} Apply each rate to its slice only:
\begin{align*}
\text{Slice 1: } & 100\,000 \times 0\% = 0 \\
\text{Slice 2: } & (300\,000 - 100\,000) \times 10\% = 20\,000 \\
\text{Slice 3: } & (400\,000 - 300\,000) \times 20\% = 20\,000 \\
\text{Total tax} &= 0 + 20\,000 + 20\,000 = 40\,000 \text{ FCFA}
\end{align*}
Average rate: $\mathrm{ATR} = \dfrac{40\,000}{400\,000}\times 100 = 10\%$. \par
Marginal rate: the last franc earned falls in the $20\%$ bracket, so $\mathrm{MTR} = 20\%$. \par
\textbf{Remark:} the marginal rate ($20\%$) is higher than the average rate ($10\%$) --- a typical feature of progressive taxation. A common error is to apply $20\%$ to the whole $400\,000$ FCFA, which would give $80\,000$ FCFA: brackets apply \emph{slice by slice}, never to total income.
\end{exoresolu}

\begin{popfigure}
\begin{center}
\begin{tikzpicture}
\begin{axis}[
  width=11cm, height=7.5cm,
  xlabel={Income (thousands of FCFA)},
  ylabel={Tax paid (thousands of FCFA)},
  xmin=0, xmax=500, ymin=0, ymax=120,
  xtick={0,100,200,300,400,500},
  ytick={0,20,40,60,80,100,120},
  grid=both, grid style={popline!40},
  legend pos=north west,
  legend style={font=\small},
  axis line style={popdark},
]
\addplot[very thick, popblue, domain=0:500] {0.2*x};
\addlegendentry{Proportional (20\%)}
\addplot[very thick, poporange, domain=0:500, samples=100] {0.0004*x^2};
\addlegendentry{Progressive}
\addplot[very thick, popgreen, domain=0:500, samples=100] {4.5*sqrt(max(0,x))};
\addlegendentry{Regressive}
\end{axis}
\end{tikzpicture}
\end{center}
\legende{Tax paid against income. Proportional: straight line (constant rate). Progressive: curve steepening (rate rises with income). Regressive: curve flattening (effective rate falls as income rises). All three meet at (500,100) for comparison.}
\end{popfigure}

Reading the graph: for the proportional tax, tax paid grows in a straight line --- the rate never changes. For the progressive tax, the curve bends upwards: each additional franc of income is taxed more heavily than the previous one. For the regressive tax, the curve flattens: tax still rises in absolute terms, but more and more slowly relative to income, so the \emph{effective} rate falls.

\courssub{Other Classifications of Taxes}

\textbf{Personal taxes versus real taxes.}\par
A \cle{personal tax} takes into account the situation of the taxpayer (family size, total income, ability to pay): the progressive income tax is personal. A \cle{real tax} is attached to the \emph{thing} (the object) regardless of who owns it: property tax on a building or customs duty on a container is the same whoever the owner is.

\textbf{Central versus local taxes.}\par
\cle{Central (State) taxes} are levied by the national government (income tax, company tax, VAT, customs duties). \cle{Local taxes} are levied by decentralised authorities --- in Cameroon, the regional and local authorities (councils) collect taxes such as property tax, market dues and business licence fees (\emph{patente}) to finance local services like markets, roads and waste collection.

\textbf{Specific versus ad valorem duties.}\par
A \cle{specific duty} is a fixed amount per physical unit: for example, $X$ FCFA per litre of beer or per packet of cigarettes. An \cle{ad valorem duty} is a percentage of the \emph{value} of the good: for example, a customs duty of $20\%$ of the value of an imported vehicle. VAT is an ad valorem tax; most excise duties combine specific and ad valorem elements.

\begin{exemplebox}[Cameroonian examples of each type]
\begin{itemize}
\item \textbf{Direct, progressive, personal, central:} the IRPP on salaries, withheld monthly by employers according to rising brackets.
\item \textbf{Direct, proportional, central:} company tax on the profits of firms such as ENEO or the breweries.
\item \textbf{Direct, real, local:} property tax on buildings, collected by councils.
\item \textbf{Indirect, proportional, ad valorem:} VAT charged on goods and services sold in shops and supermarkets.
\item \textbf{Indirect, specific/ad valorem:} excise duties on beer, spirits and cigarettes.
\item \textbf{Indirect, ad valorem:} customs duties on imports through the port of Douala under the CEMAC common external tariff.
\end{itemize}
\end{exemplebox}

\begin{attention}[Exam tip: always illustrate with a Cameroonian example]
In the GCE, a definition without an example earns only half the marks. Whenever you define a type of tax (direct, indirect, proportional, progressive, regressive, specific, ad valorem), immediately give a Cameroonian illustration: IRPP brackets for progressive income tax, VAT for indirect taxation, excise duties on tobacco for specific duties, customs duties at the port of Douala for ad valorem duties. This habit transforms an average essay into an excellent one.
\end{attention}

\begin{exercice}[Classify and compute]
\begin{enumerate}
\item Classify each of the following as direct or indirect, and justify: (a) company tax paid by a brewery; (b) VAT on a telephone recharge; (c) property tax on a house in Buea; (d) customs duty on imported rice.
\item Miss Ngo earns $600\,000$ FCFA per year. Using the brackets of the worked example above ($0\%$ up to $100\,000$; $10\%$ from $100\,001$ to $300\,000$; $20\%$ from $300\,001$ to $500\,000$; assume $30\%$ above $500\,000$), compute her total tax, her average rate and her marginal rate.
\item Explain, with one example each, the difference between a specific duty and an ad valorem duty.
\end{enumerate}
\end{exercice}

\begin{corrige}[Exercise: Classify and compute]
\textbf{1.} (a) Direct: levied on the company's profit, borne by the company. (b) Indirect: collected by the seller but borne by the consumer. (c) Direct: levied on the owner's wealth, cannot be shifted. (d) Indirect: the importer shifts it into the selling price of rice. \par
\textbf{2.} Slices: $0 + (200\,000 \times 10\%) + (200\,000 \times 20\%) + (100\,000 \times 30\%) = 0 + 20\,000 + 40\,000 + 30\,000 = 90\,000$ FCFA. Average rate: $90\,000/600\,000 = 15\%$. Marginal rate: $30\%$ (the bracket of the last franc earned). \par
\textbf{3.} A specific duty is a fixed sum per unit (e.g.\ a fixed amount of FCFA per packet of cigarettes), independent of price; an ad valorem duty is a percentage of value (e.g.\ $20\%$ of the value of an imported car), so the tax rises automatically with the price of the good.
\end{corrige}

\begin{aretenir}[Key points]
\begin{itemize}
\item \textbf{Direct taxes} fall on income and wealth of identifiable persons (PIT, company tax, property tax); \textbf{indirect taxes} fall on expenditure and can be shifted to consumers (VAT, excise duties, customs duties).
\item Classification depends on who \emph{bears} the tax (incidence), not who collects it (impact): VAT is indirect.
\item \textbf{Proportional}: constant rate (ATR = MTR); \textbf{progressive}: rate rises with the base (MTR $>$ ATR); \textbf{regressive}: effective rate falls as the base rises (MTR $<$ ATR).
\item Average rate $=$ tax $\div$ income; marginal rate $=$ rate on the last franc earned. With brackets, tax is computed slice by slice.
\item Other distinctions: personal vs real, central vs local, specific vs ad valorem.
\item Direct taxes are fairer but easier to evade; indirect taxes are productive and hard to evade but regressive and inflationary.
\end{itemize}
\end{aretenir}

\coursec{VAT in Cameroon and the Structure of Taxation (Practical Work)}

\courssub{Definition and Mechanism of VAT}

\begin{definition}[Value Added Tax (VAT)]
The \cle{Value Added Tax (VAT)} is a \cle{multi-stage indirect tax} levied on the \cle{value added} at each stage of the production and distribution chain. Unlike a single-stage sales tax, VAT is collected fractionally: each taxable person (producer, wholesaler, retailer) charges VAT on their sales (\cle{output VAT}) and deducts the VAT they have paid on their purchases (\cle{input VAT}). Only the difference is remitted to the State. The final consumer bears the full tax burden because they cannot deduct any input VAT.
\end{definition}

\begin{definition}[Value Added]
At a given stage, \cle{value added} is the difference between the \cle{sales revenue (excluding VAT)} of the firm and the \cle{cost of intermediate goods and services purchased from other firms (excluding VAT)}. It represents the new value created by the firm's own factors of production (labour, capital, entrepreneurship).
\[
\text{Value Added} = \text{Sales (HT)} - \text{Purchases of intermediate goods (HT)}
\]
\end{definition}

\begin{savaistu}[Why VAT replaced the old turnover tax in Cameroon]
Before 1999, Cameroon applied a \cle{turnover tax} levied on the \emph{total} sales value at each stage, causing \cle{cascading} (tax charged on tax already embedded in the price). This distorted production decisions and encouraged artificial vertical integration. VAT, introduced by the Finance Law No. 98/004 of 14 April 1998 and effective from 1 January 1999, eliminates cascading because tax is charged only on the \emph{value added} at each stage. It has since become the backbone of Cameroon's domestic indirect tax revenue.
\end{savaistu}

\courssub{VAT Rate in Cameroon and Exempt/Zero-Rated Goods}

\begin{aretenir}[VAT Rate in Cameroon]
\begin{itemize}
    \item \textbf{Standard rate:} \cle{19.25\%}, composed of 17.5\% VAT proper plus 10\% council surcharges on the VAT (\emph{centimes additionnels communaux}), i.e. $17.5\% \times 1.10 = 19.25\%$.
    \item \textbf{Zero-rated (exonérés avec droit à déduction):} exports and international transport. The seller charges 0\% but keeps the right to deduct input VAT.
    \item \textbf{Exempt (exonérés):} a list fixed by the Finance Law, which typically includes certain basic necessities, medical and educational services, financial services and insurance, and some agricultural inputs. The exact list changes with each Finance Law; in the examination, use the list given in the question.
\end{itemize}
\end{aretenir}

\begin{attention}[Zero-rated vs. Exempt — a crucial distinction for exams]
\begin{itemize}
    \item \textbf{Zero-rated (with right to deduct):} the firm charges 0\% VAT on sales \textbf{BUT} can still deduct input VAT on purchases. The State may even refund the resulting credit. This \textbf{encourages} the sector (e.g. exporters).
    \item \textbf{Exempt without right to deduct:} the firm charges no VAT on sales \textbf{AND} cannot deduct input VAT. The VAT paid on inputs becomes a \textbf{cost} for the firm, usually passed on in higher prices.
\end{itemize}
\end{attention}

\courssub{Calculation of VAT: From HT to TTC and Reverse}

\begin{methode}[Three-Step VAT Calculation (HT $\to$ VAT $\to$ TTC)]
Let $r = 0.1925$ be the VAT rate.
\begin{enumerate}
    \item \textbf{Start from the VAT-exclusive price (Prix HT).}
    \item \textbf{Compute the VAT amount:} $\text{VAT} = \text{HT} \times r$.
    \item \textbf{Compute the VAT-inclusive price (Prix TTC):} $\text{TTC} = \text{HT} + \text{VAT} = \text{HT} \times (1+r)$.
\end{enumerate}
\end{methode}

\begin{methode}[Reverse Calculation (TTC $\to$ HT $\to$ VAT)]
Given a TTC price:
\begin{enumerate}
    \item \textbf{Compute HT:} $\text{HT} = \dfrac{\text{TTC}}{1+r}$.
    \item \textbf{Compute VAT:} $\text{VAT} = \text{TTC} - \text{HT} = \text{TTC} \times \dfrac{r}{1+r}$.
\end{enumerate}
\end{methode}

\begin{exoresolu}[VAT Calculation on a Cameroonian Invoice]
A retailer in Douala sells a bag of locally milled rice at a VAT-exclusive price of \SI{25000}{FCFA}.
\begin{enumerate}
    \item Calculate the VAT amount.
    \item Calculate the TTC price the consumer pays.
    \item If the price tag shows \SI{29812.5}{FCFA} TTC, verify the HT price.
\end{enumerate}
\tcblower
\textbf{Solution:}
Rate $r = 19.25\% = 0.1925$.
\begin{enumerate}
    \item $\text{VAT} = 25\,000 \times 0.1925 = \SI{4812.5}{FCFA}$.
    \item $\text{TTC} = 25\,000 + 4\,812.5 = \SI{29812.5}{FCFA}$.
    \item Reverse: $\text{HT} = \dfrac{29\,812.5}{1.1925} = \SI{25000}{FCFA}$. Correct.
\end{enumerate}
\end{exoresolu}

\begin{attention}[Classic Mistake: Applying the Rate Directly to TTC]
\textbf{Wrong:} $\text{VAT} = \text{TTC} \times 0.1925$.
\textbf{Why it is wrong:} the tax base is the HT price, not the TTC price. The VAT fraction contained in a TTC price is $\frac{r}{1+r} = \frac{0.1925}{1.1925} \approx 16.14\%$, not 19.25\%.
\textbf{Example:} if TTC = \SI{119.25}{FCFA}, the wrong VAT is $119.25 \times 0.1925 = 22.95$ FCFA (too high). The correct VAT is $119.25 \times \frac{0.1925}{1.1925} = 19.25$ FCFA, and HT = 100 FCFA. Check: $100 \times 0.1925 = 19.25$. Always extract HT first!
\end{attention}

\courssub{VAT Deduction Mechanism: Output VAT minus Input VAT}

\begin{propriete}[VAT Payable to the State]
At each reporting period (monthly in Cameroon), a VAT-registered taxpayer computes:
\[
\text{VAT Payable} = \text{Output VAT (collected on sales)} - \text{Input VAT (deductible on purchases)}
\]
\begin{itemize}
    \item If positive: the firm pays the difference to the tax administration (DGI).
    \item If negative: the firm carries forward a \cle{VAT credit} (crédit de TVA) to offset future liabilities. Refunds are possible in specific cases fixed by law, notably for exporters (subject to conditions).
\end{itemize}
\end{propriete}

\begin{aretenir}[Conditions for Input VAT Deduction in Cameroon]
Input VAT is deductible only if:
\begin{enumerate}
    \item The purchase is used for \textbf{taxable operations} (not exempt without right to deduct).
    \item The supplier is \textbf{VAT-registered} and issues a valid \cle{invoice} mentioning the taxpayer identification number (NIU), the HT amount, the VAT rate and the VAT amount.
    \item The invoice is properly recorded in the buyer's \textbf{purchase journal}.
    \item The goods or services are not \textbf{excluded from deduction by law} (e.g. private passenger vehicles and related fuel, certain luxury or personal expenses).
\end{enumerate}
\end{aretenir}

\courssub{Practical Work 21: VAT Calculation Along a Production-Distribution Chain}

We illustrate the fractional collection of VAT along a simplified chain: \textbf{Producer $\to$ Wholesaler $\to$ Retailer $\to$ Final Consumer}. Assume the standard rate $r = 19.25\%$. All amounts are in FCFA.

\begin{popfigure}
\begin{center}
\begin{tabular}{|c|c|c|c|c|c|c|}
\hline
\rowcolor{popblueL}
\textbf{Stage} & \textbf{Sales HT} & \textbf{Output VAT} & \textbf{Purchases HT} & \textbf{Input VAT} & \textbf{VAT Payable} & \textbf{Value Added} \\
\hline
Producer & 100\,000 & 19\,250 & 0 & 0 & 19\,250 & 100\,000 \\
\hline
Wholesaler & 150\,000 & 28\,875 & 100\,000 & 19\,250 & 9\,625 & 50\,000 \\
\hline
Retailer & 200\,000 & 38\,500 & 150\,000 & 28\,875 & 9\,625 & 50\,000 \\
\hline
\textbf{Total} & & \textbf{86\,625} & & \textbf{48\,125} & \textbf{38\,500} & \textbf{200\,000} \\
\hline
\end{tabular}
\end{center}
\legende{VAT calculation along a three-stage chain (Cameroon, rate 19.25\%)}
\end{popfigure}

\begin{experience}[Guided Demonstration: Filling the Table]
\begin{enumerate}
    \item \textbf{Producer:} assume no deductible purchases. Sales HT = 100\,000. Output VAT = $100\,000 \times 0.1925 = 19\,250$. Input VAT = 0. VAT Payable = 19\,250. Value Added = 100\,000.
    \item \textbf{Wholesaler:} buys from the Producer at 100\,000 HT + 19\,250 VAT. Sales HT = 150\,000. Output VAT = $150\,000 \times 0.1925 = 28\,875$. Input VAT = 19\,250 (shown on the Producer's invoice). VAT Payable = $28\,875 - 19\,250 = 9\,625$. Value Added = $150\,000 - 100\,000 = 50\,000$. Check: $50\,000 \times 0.1925 = 9\,625$. \checkmark
    \item \textbf{Retailer:} buys from the Wholesaler at 150\,000 HT + 28\,875 VAT. Sales HT = 200\,000. Output VAT = $200\,000 \times 0.1925 = 38\,500$. Input VAT = 28\,875. VAT Payable = $38\,500 - 28\,875 = 9\,625$. Value Added = $200\,000 - 150\,000 = 50\,000$. Check: $50\,000 \times 0.1925 = 9\,625$. \checkmark
    \item \textbf{Consumer:} pays TTC = $200\,000 + 38\,500 = 238\,500$ and bears the full VAT burden of 38\,500.
    \item \textbf{State revenue:} sum of VAT Payable = $19\,250 + 9\,625 + 9\,625 = 38\,500$, exactly equal to the VAT on final consumption.
\end{enumerate}
\end{experience}

\begin{aretenir}[Key Properties of the VAT Chain]
\begin{itemize}
    \item \textbf{Neutrality:} the total VAT collected by the State (38\,500) equals the VAT on the final sale to the consumer. Intermediate firms are mere \cle{tax collectors}; VAT does not enter their costs (when fully deductible).
    \item \textbf{Value added = tax base at each stage:} VAT Payable = Value Added $\times$ Rate. This proves that VAT is genuinely a tax on value added.
    \item \textbf{Invoice trail:} the mechanism relies on the \cle{invoice trail}. If the Wholesaler fails to issue a valid invoice, the Retailer cannot deduct input VAT and the chain is broken.
\end{itemize}
\end{aretenir}

\courssub{The Structure of Taxation: Tax Base, Assessment, Collection}

\begin{definition}[Tax Base (Assiette)]
The \cle{tax base} is the \textbf{quantitative measure} of the taxable matter (income, profit, turnover, value added, property value, etc.) to which the tax rate is applied. It must be clearly defined by law to avoid arbitrariness.
\end{definition}

\begin{definition}[Assessment (Liquidation)]
\cle{Assessment} is the administrative operation of \textbf{determining the tax liability} of a specific taxpayer: identifying the base, applying the rate, calculating the amount due, and issuing the \cle{tax notice (avis d'imposition)}.
\end{definition}

\begin{definition}[Collection (Recouvrement)]
\cle{Collection} is the \textbf{actual recovery} of the assessed tax amount from the taxpayer into the Public Treasury. It involves payment deadlines, enforcement procedures in case of default, and the accounting for revenue.
\end{definition}

\begin{aretenir}[The Three Phases in Cameroon]
\begin{center}
\begin{tabular}{|l|l|l|}
\hline
\rowcolor{popblueL}
\textbf{Phase} & \textbf{Responsible Body} & \textbf{Key Document} \\
\hline
1. Base determination & DGI (Direction Générale des Impôts) & Tax return (déclaration) \\
\hline
2. Assessment (liquidation) & DGI / Taxation centre & Avis d'imposition \\
\hline
3. Collection (recouvrement) & Public Treasury (Trésor Public) & Receipt (quittance) \\
\hline
\end{tabular}
\end{center}
\end{aretenir}

\courssub{Practical Work 22: Structure of Taxation in Cameroon}

Cameroon's tax system is divided into \cle{direct taxes} (borne by the taxpayer on income or wealth) and \cle{indirect taxes} (borne by the consumer on expenditure). The share of each in total tax revenue reflects the country's economic structure: a large informal sector, a narrow formal wage base, and the importance of consumption and foreign trade.

\begin{popfigure}
\begin{center}
\begin{tikzpicture}
\begin{axis}[
    ybar,
    bar width=12pt,
    width=\linewidth,
    height=8cm,
    enlarge x limits=0.12,
    ylabel={Share in total tax revenue (\%)},
    ymin=0, ymax=40,
    xtick=data,
    xticklabels={VAT, Company tax (IS), Personal income tax (IRPP), Excise duties, Customs duties, Other direct, Other indirect},
    x tick label style={rotate=45, anchor=east, font=\small},
    ymajorgrids=true,
    grid style=dashed,
    nodes near coords,
    nodes near coords align={vertical},
    every node near coord/.append style={font=\small, popdark},
]
\addplot[fill=popblue] coordinates {(1,32) (2,22) (3,12) (4,10) (5,8) (6,8) (7,8)};
\end{axis}
\end{tikzpicture}
\end{center}
\legende{Indicative structure of tax revenue in Cameroon (orders of magnitude for illustration)}
\end{popfigure}

\begin{savaistu}[Reading the Chart — Cameroon's Tax Structure]
\begin{itemize}
    \item \textbf{VAT} is the \textbf{single largest tax} (roughly one third of tax revenue), reflecting the importance of consumption taxation and formal-sector compliance.
    \item \textbf{Company tax (Impôt sur les Sociétés, IS)} comes second, driven by oil and refining (SONARA), banking, telecommunications and agro-industry (e.g. CDC, SOCAPALM, HEVECAM).
    \item \textbf{Personal income tax (IRPP)} yields relatively little because of the large informal sector, the exemption threshold and the narrow formal wage base.
    \item \textbf{Excise duties} fall on alcohol, tobacco, petroleum products and certain luxury goods.
    \item \textbf{Customs duties} have declined in relative importance with regional integration (CEMAC Common External Tariff) and trade liberalisation (AfCFTA).
    \item \textbf{Other direct taxes:} property tax, business licence (patente), etc. \textbf{Other indirect taxes:} registration and stamp duties.
\end{itemize}
\emph{Note:} the exact shares vary from year to year; the chart gives orders of magnitude, not official statistics.
\end{savaistu}

\begin{exoresolu}[Interpreting a Tax Revenue Table]
A simplified tax revenue table for Cameroon (in billion FCFA):

\begin{center}
\begin{tabular}{|l|r|r|}
\hline
\rowcolor{popblueL}
\textbf{Tax} & \textbf{Amount} & \textbf{\% of total} \\
\hline
VAT & 1\,200 & 34.3\% \\
\hline
Company tax (IS) & 800 & 22.9\% \\
\hline
Personal income tax (IRPP) & 400 & 11.4\% \\
\hline
Excise duties & 350 & 10.0\% \\
\hline
Customs duties & 250 & 7.1\% \\
\hline
Other direct & 250 & 7.1\% \\
\hline
Other indirect & 250 & 7.1\% \\
\hline
\textbf{Total} & \textbf{3\,500} & \textbf{100\%} \\
\hline
\end{tabular}
\end{center}

\textbf{Questions:}
\begin{enumerate}
    \item Calculate the total tax revenue.
    \item What percentage of total revenue comes from indirect taxes?
    \item If the government wants to reduce reliance on customs duties (trade liberalisation), which tax should it strengthen? Justify.
\end{enumerate}
\tcblower
\textbf{Solution:}
\begin{enumerate}
    \item Total = $1\,200 + 800 + 400 + 350 + 250 + 250 + 250 = 3\,500$ billion FCFA.
    \item Indirect taxes = VAT + Excise + Customs + Other indirect = $1\,200 + 350 + 250 + 250 = 2\,050$. Share = $\dfrac{2\,050}{3\,500} \times 100 \approx 58.6\%$.
    \item Strengthen \textbf{VAT} (broaden the base, improve compliance) or \textbf{excise duties} on goods with inelastic demand (fuel, tobacco, alcohol). VAT is preferred because it is broad-based, neutral between firms, and grows automatically with the formal economy. Excise duties are easier to administer but rest on a narrow base.
\end{enumerate}
\end{exoresolu}

\courssub{Reading a Cameroonian Price Tag or Tax Notice}

\begin{experience}[Case Study: Decoding a Cameroonian Price Tag]
A price tag in a supermarket in Yaoundé reads:
\begin{center}
\textbf{PRIX TTC : 11\,925 FCFA} \quad \textbf{TVA (19.25\%) INCLUSE}
\end{center}
\textbf{Task:} determine the HT price and the VAT amount.
\textbf{Step-by-step:}
\begin{enumerate}
    \item Identify TTC = 11\,925 FCFA and rate $r = 19.25\%$.
    \item HT = TTC / (1+r) = $11\,925 / 1.1925 = \textbf{10\,000 FCFA}$.
    \item VAT = TTC $-$ HT = $11\,925 - 10\,000 = \textbf{1\,925 FCFA}$.
    \item Check: $10\,000 \times 0.1925 = 1\,925$. \checkmark
\end{enumerate}
\textbf{Note:} a tag marked ``TTC'' or ``TVA comprise'' already includes VAT; a tag marked ``HT'' means VAT is added at the till. Always check that the rate used is the current legal rate.
\end{experience}

\begin{experience}[Case Study: Reading a Simplified Income Tax Notice (IRPP)]
An employee receives a tax notice based on a \emph{simplified illustrative} progressive scale (the real scale is fixed by the Finance Law and includes council surcharges; always use the scale given in the examination question):
\begin{itemize}
    \item \textbf{Net taxable income (RNI):} 3\,600\,000 FCFA (after allowances on gross salary).
    \item \textbf{Illustrative brackets and rates:}
    \begin{center}
    \begin{tabular}{|l|c|}
    \hline
    \rowcolor{popblueL}
    \textbf{Bracket (FCFA)} & \textbf{Rate} \\
    \hline
    0 -- 1\,000\,000 & 0\% \\
    \hline
    1\,000\,001 -- 2\,000\,000 & 10\% \\
    \hline
    2\,000\,001 -- 3\,000\,000 & 15\% \\
    \hline
    3\,000\,001 -- 4\,000\,000 & 20\% \\
    \hline
    Above 4\,000\,000 & 25\% \\
    \hline
    \end{tabular}
    \end{center}
    \item \textbf{Tax computation (liquidation):}
    \begin{align*}
        \text{Bracket 1: } & 1\,000\,000 \times 0\% = 0 \\
        \text{Bracket 2: } & 1\,000\,000 \times 10\% = 100\,000 \\
        \text{Bracket 3: } & 1\,000\,000 \times 15\% = 150\,000 \\
        \text{Bracket 4: } & 600\,000 \times 20\% = 120\,000 \\
        \text{Total tax} & = 370\,000 \text{ FCFA}
    \end{align*}
    \item \textbf{Pay-as-you-earn withholding:} 370\,000 FCFA already deducted monthly by the employer.
    \item \textbf{Balance to pay / refund:} 0 FCFA.
\end{itemize}
\textbf{Lesson:} the notice displays the three phases studied above: the \textbf{base} (RNI), the \textbf{assessment} (application of the scale), and the \textbf{collection} (withholding at source). The taxpayer's job is to verify the base and the computation.
\end{experience}

\begin{attention}[Common Exam Traps on VAT and Tax Structure]
\begin{enumerate}
    \item \textbf{Confusing the VAT rate with the VAT fraction in a TTC price.} Always use $\frac{r}{1+r}$ to extract VAT from TTC.
    \item \textbf{Forgetting that exemption without right to deduct breaks the deduction chain.} A landlord letting exempt residential housing cannot recover the VAT paid on repairs or electricity.
    \item \textbf{Mixing up base, assessment and collection.} Base = \emph{what} is taxed (e.g. turnover); assessment = \emph{calculation} of the amount due; collection = \emph{actually getting the money} into the Treasury.
    \item \textbf{Assuming customs duties are still the main revenue.} With the CEMAC Common External Tariff and AfCFTA, their relative share has fallen; VAT and company tax now dominate.
    \item \textbf{Ignoring the invoice requirement.} No valid invoice $\Rightarrow$ no input VAT deduction $\Rightarrow$ a higher effective tax cost for the buyer.
\end{enumerate}
\end{attention}

\begin{exercice}[Practice Questions]
\begin{enumerate}
    \item A furniture maker in Bafoussam buys wood for 500\,000 FCFA HT and pays 96\,250 FCFA VAT. He sells finished furniture for 1\,200\,000 FCFA HT. Calculate his output VAT, input VAT, VAT payable, and value added.
    \item Explain why a zero-rated export firm (e.g. a cocoa bean exporter) benefits from VAT refunds, while a bank (exempt financial services) does not.
    \item Using the tax structure chart, argue whether Cameroon's tax system relies more on \emph{direct} or on \emph{indirect} taxation. What are the equity implications?
    \item A price tag shows 50\,000 FCFA TTC. A student calculates VAT as $50\,000 \times 19.25\% = 9\,625$ FCFA. Explain the error and give the correct VAT and HT.
\end{enumerate}
\end{exercice}

\begin{corrige}[Exercise 1]
\begin{enumerate}
    \item Purchases HT = 500\,000; input VAT = $500\,000 \times 0.1925 = 96\,250$ (given). Sales HT = 1\,200\,000; output VAT = $1\,200\,000 \times 0.1925 = 231\,000$. VAT payable = $231\,000 - 96\,250 = 134\,750$. Value added = $1\,200\,000 - 500\,000 = 700\,000$. Check: $700\,000 \times 0.1925 = 134\,750$. \checkmark
    \item Zero-rated exporter: charges 0\% VAT on exports but deducts input VAT on all inputs (fertiliser, transport, bags). Since output VAT = 0 and input VAT $>$ 0, a VAT credit arises, which may be refunded. Bank (exempt): charges no VAT on its services but cannot deduct input VAT on rent, electricity, IT or security; the VAT on inputs becomes a cost, passed on to customers through higher bank charges.
    \item Indirect taxes (VAT + excise + customs + other indirect) $\approx$ 58--60\%; direct taxes (IS + IRPP + other direct) $\approx$ 40--42\%. The system is \textbf{indirect-heavy}. Equity implication: indirect taxes are \textbf{regressive} (the poor spend a higher share of their income on consumption). Heavy reliance on VAT and excise duties burdens low-income households unless mitigated by exempting or zero-rating basic goods (basic foodstuffs, essential medicines).
    \item Error: the rate was applied to the TTC price instead of the HT base. Correct: HT = $50\,000 / 1.1925 \approx 41\,928.7$ FCFA; VAT = $50\,000 - 41\,928.7 \approx 8\,071.3$ FCFA (equivalently, VAT = $50\,000 \times \frac{0.1925}{1.1925}$).
\end{enumerate}
\end{corrige}

\coursec{Economic Consequences of Taxation}

In the previous section we examined how taxes are structured in Cameroon, with a practical focus on Value Added Tax. We now turn to the \emph{economic impact} of those taxes. A tax is not merely a transfer of money from the private sector to the State; it changes the incentives that guide the decisions of households and firms. Understanding these consequences is essential for evaluating whether a tax system is efficient, equitable, and sustainable — a classic GCE essay question.

\courssub{Effects on Consumption and Household Disposable Income}

The most immediate effect of a direct tax (such as personal income tax) is a reduction in the \cle{disposable income} of households.

\begin{definition}[Disposable Income]
The income available to a household for spending and saving after payment of direct taxes and receipt of transfer payments (pensions, family allowances, subsidies).
\[
Y_d = Y - T_d + Tr
\]
where \(Y\) is gross income, \(T_d\) direct taxes, and \(Tr\) transfers.
\end{definition}

When disposable income falls, consumption typically falls as well, but by a smaller proportion because households also reduce saving. The reduction in consumption depends on the \cle{marginal propensity to consume} (MPC). For a tax increase \(\Delta T\), the initial drop in consumption is \(\text{MPC} \times \Delta T\). In Cameroon, where a large share of the population lives near the subsistence level, the MPC is high; thus direct taxes on low-income earners can significantly depress aggregate demand. This is one reason why the personal income tax exempts the lowest income band.

\courssub{Effects on Production, Investment and the Cost of Labour (Tax Wedge)}

Taxes on firms (corporate tax, payroll taxes, social security contributions) and on factors of production alter production costs and investment decisions.

\begin{definition}[Tax Wedge]
The difference between the total labour cost paid by the employer and the net wage received by the employee, expressed as a percentage of total labour cost. It measures the distortion introduced by taxes and social contributions on the labour market.
\[
\text{Tax Wedge} = \frac{\text{Total Labour Cost} - \text{Net Wage}}{\text{Total Labour Cost}} \times 100\%
\]
\end{definition}

In Cameroon, the employer and the employee both contribute to the \emph{Caisse Nationale de Prévoyance Sociale} (CNPS), and employers pay additional payroll levies. A high tax wedge can:
\begin{itemize}
    \item Discourage formal hiring, pushing employment into the informal sector.
    \item Reduce the return on capital, lowering investment.
    \item Encourage capital-intensive rather than labour-intensive techniques, which is costly in a labour-abundant economy.
\end{itemize}

\begin{exemplebox}[Effect of Payroll Charges on a Cameroonian Firm (Illustrative Figures)]
A firm in Douala employs 50 workers. The gross wage per worker is \SI{200000}{FCFA} per month. For illustration, assume the employer pays social contributions and payroll levies totalling 20\% of the gross wage, while the employee pays 10\% in social contributions and income tax.
\begin{itemize}
    \item Total labour cost per worker \(= 200000 + 20\% \times 200000 = 200000 + 40000 = \SI{240000}{FCFA}\).
    \item Net wage \(= 200000 - 10\% \times 200000 = \SI{180000}{FCFA}\).
    \item Tax wedge \(= \dfrac{240000 - 180000}{240000} \times 100\% = 25\%\).
\end{itemize}
If the government raises the employer's levies from 20\% to 25\%, the total cost rises to \SI{250000}{FCFA}, the wedge widens to 28\%, and the firm may hire fewer workers, delay replacement of departing staff, or automate.
\end{exemplebox}

\courssub{Effects on Saving and Capital Formation}

Taxes on interest income, dividends, and capital gains reduce the net return on saving. This can discourage household saving and, by reducing the pool of loanable funds, raise interest rates and crowd out private investment. Conversely, tax incentives (for example, tax advantages granted to approved enterprises under Cameroon's investment legislation) aim to stimulate capital formation.

\begin{propriete}[Taxation and the Cost of Capital]
The \cle{user cost of capital} is the real cost of using one unit of capital for one period; in its simplest form \(c = r + \delta\), where \(r\) is the real interest rate and \(\delta\) the depreciation rate. A corporate tax at rate \(\tau\) raises the pre-tax return a project must earn to be worthwhile: the firm must earn enough to pay the tax and still leave investors their required net return. Hence, other things equal, a higher \(\tau\) raises the effective cost of capital and reduces the desired capital stock — unless the tax system offsets this through deductions, depreciation allowances or tax credits. (The formal proof is not examinable; the reasoning and conclusion are.)
\end{propriete}

\courssub{Redistributive Effects: Progressive Taxation and Inequality}

A \cle{progressive tax} takes a larger percentage of income as income rises. By reducing the disposable income of the rich proportionally more than that of the poor, it narrows the post-tax income distribution.

\begin{propriete}[Progressive Taxation Reduces Income Inequality]
If the tax function \(T(Y)\) is progressive (the average tax rate \(T(Y)/Y\) rises with \(Y\)), then the distribution of post-tax income \(Y - T(Y)\) is more equal than that of pre-tax income \(Y\), provided the tax does not cause re-ranking of individuals.
\end{propriete}

In Cameroon, the personal income tax (IRPP) is progressive, with rates rising across successive income bands. However, the large informal sector limits its redistributive reach, since many incomes escape the tax net. Indirect taxes (VAT, excises) are often \cle{regressive} because the poor spend a larger share of their income on taxed goods.

\begin{aretenir}[Redistribution in Practice]
\begin{itemize}
    \item Direct progressive taxes (IRPP) reduce inequality.
    \item Indirect taxes (VAT) tend to be regressive unless essential goods are zero-rated or exempted.
    \item In Cameroon, basic necessities such as unprocessed foodstuffs, medicines and educational services benefit from VAT exemptions or reduced treatment to protect the poor.
\end{itemize}
\end{aretenir}

\courssub{Tax Incidence: Legal vs. Economic Incidence}

One of the most important concepts in public finance is that the entity legally obliged to pay a tax may not be the one who ultimately bears its burden.

\begin{definition}[Tax Incidence]
\begin{itemize}
    \item \textbf{Legal (statutory) incidence:} the person or firm on whom the tax is legally imposed (who writes the cheque to the tax authority).
    \item \textbf{Economic incidence:} the final distribution of the tax burden after all market adjustments (price changes, quantity changes).
\end{itemize}
\end{definition}

\begin{definition}[Tax Shifting]
The process by which the legal taxpayer passes the tax burden to others through changes in prices.
\begin{itemize}
    \item \textbf{Forward shifting:} the producer raises the price paid by consumers.
    \item \textbf{Backward shifting:} the producer lowers the price paid to factor suppliers (e.g., lower wages, lower rent).
\end{itemize}
\end{definition}

\begin{attention}[Who Bears the Tax?]
\textbf{Never assume the legal taxpayer bears the burden.} The economic incidence depends on the \emph{price elasticities of supply and demand}. The more inelastic side of the market bears a larger share of the tax, because it cannot easily adjust the quantity it buys or sells when the price changes.
\end{attention}

\courssub{Analysing Tax Incidence with Supply and Demand}

We use a standard supply-demand diagram to show how an indirect tax (per unit) is shared between buyers and sellers.

\begin{methode}[Analysing the Effect of a Per-Unit Tax with a Supply-Demand Diagram]
\begin{enumerate}
    \item Draw the initial supply curve \(S_0\) and demand curve \(D\). Mark equilibrium \(E_0\) with price \(P_0\) and quantity \(Q_0\).
    \item Shift the supply curve vertically upward by the amount of the tax \(t\) to get \(S_1\). (The vertical distance between \(S_0\) and \(S_1\) equals \(t\) at every quantity.)
    \item Find the new equilibrium \(E_1\) where \(S_1\) intersects \(D\). The new price paid by buyers is \(P_b\); the price received by sellers (after paying the tax) is \(P_s = P_b - t\).
    \item The tax burden on buyers \(= P_b - P_0\). The burden on sellers \(= P_0 - P_s\).
    \item Government revenue \(= t \times Q_1\) (the rectangle between \(P_b\) and \(P_s\) up to \(Q_1\)).
    \item The deadweight loss (excess burden) is the triangle between \(S_0\), \(D\) and \(Q_1\): mutually beneficial trades between \(Q_1\) and \(Q_0\) no longer take place.
\end{enumerate}
\end{methode}

\begin{popfigure}
\begin{tikzpicture}
\begin{axis}[
    width=12cm, height=8cm,
    axis lines=left,
    xlabel={Quantity ($Q$)}, ylabel={Price ($P$)},
    xmin=0, xmax=13, ymin=0, ymax=15,
    xtick={4,6}, xticklabels={$Q_1$, $Q_0$},
    ytick={6,8,10}, yticklabels={$P_s$, $P_0$, $P_b$},
    clip=false,
    domain=0:13,
    samples=100,
]
% Demand curve: P = 14 - Q
\addplot[thick, popblue, domain=0:13] {14 - x} node[right, pos=0.85] {$D$};
% Supply curve before tax: P = 2 + Q
\addplot[thick, popgreen, domain=0:12] {2 + x} node[right, pos=0.85] {$S_0$};
% Supply curve after tax (shift up by 4): P = 6 + Q
\addplot[thick, poporange, dashed, domain=0:8] {6 + x} node[right, pos=0.85] {$S_1 = S_0 + t$};
% Tax wedge vertical line at Q1=4
\draw[poppurple, thick, <->] (axis cs:4,6) -- (axis cs:4,10) node[midway, right] {$t$};
% Equilibrium points
\fill[popblue] (axis cs:6,8) circle (2pt) node[above right] {$E_0$};
\fill[poporange] (axis cs:4,10) circle (2pt) node[above left] {$E_1$};
\fill[popgreen] (axis cs:4,6) circle (2pt) node[below left] {$E_s$};
% Dashed lines to axes
\draw[dashed, thin] (axis cs:0,8) -- (axis cs:6,8);
\draw[dashed, thin] (axis cs:0,10) -- (axis cs:4,10);
\draw[dashed, thin] (axis cs:0,6) -- (axis cs:4,6);
\draw[dashed, thin] (axis cs:6,0) -- (axis cs:6,8);
\draw[dashed, thin] (axis cs:4,0) -- (axis cs:4,10);
% Deadweight loss triangle
\fill[poppink, opacity=0.35] (axis cs:4,6) -- (axis cs:6,8) -- (axis cs:4,10) -- cycle;
% Government revenue rectangle
\fill[popblue, opacity=0.15] (axis cs:0,6) rectangle (axis cs:4,10);
\end{axis}
\end{tikzpicture}
\legende{Incidence of a per-unit tax $t$. The tax shifts supply from $S_0$ to $S_1$. Buyers pay $P_b$, sellers receive $P_s = P_b - t$. The burden is shared: buyers bear $P_b-P_0$, sellers bear $P_0-P_s$. Government revenue is the blue rectangle; the deadweight loss is the pink triangle.}
\end{popfigure}

\begin{exemplebox}[Incidence Calculation]
Suppose the demand and supply for a beverage in Yaoundé are:
\[
Q_d = 100 - 5P, \qquad Q_s = 5P - 20
\]
where \(P\) is in hundreds of FCFA per litre and \(Q\) in thousands of litres.
\begin{itemize}
    \item \textbf{Initial equilibrium:} \(100 - 5P = 5P - 20 \Rightarrow 10P = 120 \Rightarrow P_0 = 12\) (i.e., \SI{1200}{FCFA}), and \(Q_0 = 100 - 5(12) = 40\).
    \item \textbf{Tax:} the government imposes \(t = 4\) (i.e., \SI{400}{FCFA}) per litre on sellers. Sellers now receive \(P_s = P_b - 4\), so:
    \[
    Q_s = 5(P_b - 4) - 20 = 5P_b - 40.
    \]
    \item \textbf{New equilibrium:} \(100 - 5P_b = 5P_b - 40 \Rightarrow 10P_b = 140 \Rightarrow P_b = 14\) (\SI{1400}{FCFA}); \(P_s = 10\) (\SI{1000}{FCFA}); \(Q_1 = 30\).
    \item \textbf{Burdens:} buyers bear \(14 - 12 = 2\) (\SI{200}{FCFA}); sellers bear \(12 - 10 = 2\) (\SI{200}{FCFA}). The burden is shared equally because the two curves have slopes of equal magnitude (equal elasticities at the initial equilibrium).
    \item \textbf{Revenue:} \(R = t \times Q_1 = 4 \times 30 = 120\) (i.e., \SI{400}{FCFA} \(\times\) 30,000 litres \(=\) 12 million FCFA).
    \item \textbf{Deadweight loss:} \(DWL = \tfrac{1}{2} \times t \times (Q_0 - Q_1) = 0.5 \times 4 \times (40 - 30) = 20\) (i.e., 2 million FCFA).
\end{itemize}
\end{exemplebox}

\courssub{The Laffer Curve: Tax Rates and Tax Revenue (GCE Extension)}

There is a theoretical relationship between the tax rate and total tax revenue. At a 0\% rate, revenue is zero. At a 100\% rate, revenue is also zero, because there is no incentive to work, produce or declare income. Somewhere in between lies a revenue-maximising rate.

\begin{definition}[Laffer Curve]
A curve showing the relationship between the tax rate and tax revenue, illustrating that beyond a certain tax rate, further increases in the rate lead to a decrease in revenue due to the disincentive effects on the tax base (labour supply, investment, compliance).
\end{definition}

\begin{popfigure}
\begin{tikzpicture}
\begin{axis}[
    width=10cm, height=7cm,
    axis lines=middle,
    xlabel={Tax Rate ($t$)}, ylabel={Tax Revenue ($R$)},
    xmin=0, xmax=1.05, ymin=0, ymax=1.1,
    xtick={0.5}, xticklabels={$t^*$},
    ytick={1}, yticklabels={$R_{max}$},
    domain=0:1,
    samples=100,
    clip=false,
]
% Laffer curve: inverted parabola R = 4t(1-t)
\addplot[thick, poppurple, domain=0:1] {4*x*(1-x)} node[right, pos=0.8] {$R(t)$};
% Maximum point
\fill[poppurple] (axis cs:0.5,1) circle (2pt) node[above right] {$E$};
% Dashed lines
\draw[dashed, thin] (axis cs:0.5,0) -- (axis cs:0.5,1);
\draw[dashed, thin] (axis cs:0,1) -- (axis cs:0.5,1);
% Labels
\node[align=center] at (axis cs:0.22,0.55) {Revenue\\rising};
\node[align=center] at (axis cs:0.78,0.55) {Revenue\\falling};
\end{axis}
\end{tikzpicture}
\legende{The Laffer curve. Revenue rises with the tax rate up to $t^*$, then falls. The peak $t^*$ is the revenue-maximising rate. Note: this is a theoretical construct; the exact shape and the value of $t^*$ are unknown and debated.}
\end{popfigure}

\begin{savaistu}[The Laffer Curve in Policy Debates]
The Laffer curve became famous in the 1980s (the supply-side policies of Reagan in the United States and Thatcher in the United Kingdom), where it was used to argue that cutting very high marginal tax rates could \emph{increase} revenue by boosting effort, growth and compliance. In Cameroon, the standard corporate income tax rate is 33\% (30\% plus the council surtax), having been lowered from 38.5\% by earlier reforms, partly to broaden the tax base and reduce evasion. Most economists consider that Cameroon lies on the rising (left-hand) side of the curve for most taxes: moderate rate increases still raise revenue.
\end{savaistu}

\courssub{Negative and Positive Effects of Taxation: A Summary}

\begin{center}
\begin{tabularx}{\linewidth}{|l|X|X|}
\hline
\rowcolor{popblueL}
\textbf{Dimension} & \textbf{Negative Effects (Costs)} & \textbf{Positive Effects (Benefits)} \\
\hline
\textbf{Allocation} & Distorts relative prices $\rightarrow$ deadweight loss (excess burden). & Corrects externalities (taxes on pollution, tobacco, alcohol). \\
\hline
\textbf{Distribution} & Regressive taxes worsen inequality. & Progressive taxes and targeted transfers reduce inequality. \\
\hline
\textbf{Stabilisation} & Pro-cyclical tax hikes in recession deepen the downturn. & Automatic stabilisers (progressive income tax) dampen cycles. \\
\hline
\textbf{Growth} & High rates discourage effort, saving, investment, entrepreneurship. & Finances infrastructure, education, health $\rightarrow$ raises productivity. \\
\hline
\textbf{Administration} & Compliance costs, evasion, corruption. & Funds the state apparatus that enforces contracts and property rights. \\
\hline
\end{tabularx}
\end{center}

\begin{exoresolu}[Analysing the Economic Impact of a New Excise Tax]
\textbf{Statement:} The government of Cameroon proposes a new excise tax of \SI{100}{FCFA} per litre on locally produced fruit juices to finance a rural electrification programme. The current market price is \SI{800}{FCFA} per litre. Demand is relatively elastic (many substitutes: water, soft drinks, imports), while supply is relatively inelastic in the short run (fixed plant capacity).
\begin{enumerate}
    \item Using a diagram, show the incidence of the tax.
    \item Who bears the larger burden? Explain.
    \item Discuss two potential negative effects and one positive effect of this tax.
\end{enumerate}
\tcblower
\textbf{Detailed Solution:}
\begin{enumerate}
    \item \textbf{Diagram:} Draw a steep supply curve (inelastic) and a flat demand curve (elastic). Shift supply vertically upward by \SI{100}{FCFA}. The new consumer price \(P_b\) rises by only a small amount (say \SI{20}{FCFA}), while the net producer price \(P_s\) falls by a large amount (about \SI{80}{FCFA}). The quantity traded falls slightly.
    \item \textbf{Incidence:} because demand is more elastic than supply, \textbf{producers (sellers) bear the larger burden}. They cannot easily pass the tax on to consumers, who would switch to substitutes. Here the legal incidence (on producers) and the economic incidence (mostly on producers) coincide.
    \item \textbf{Negative effects:}
        \begin{itemize}
            \item \textbf{Deadweight loss:} the tax reduces the quantity traded below the efficient level, creating a welfare loss (the triangle on the diagram).
            \item \textbf{Discouragement of local industry:} local juice producers receive lower net prices, reducing profitability and potentially leading to layoffs or factory closures — which contradicts the policy of promoting local processing ("Made in Cameroon").
        \end{itemize}
        \textbf{Positive effect:}
        \begin{itemize}
            \item \textbf{Revenue for public goods:} the tax finances rural electrification, which can raise productivity and living standards in rural areas, potentially benefiting the very communities where the fruit is grown.
        \end{itemize}
\end{enumerate}
\end{exoresolu}

\courssub{Worked Exercise: Tax Incidence with Linear Functions}

\begin{exercice}[Tax Incidence Calculation]
The market for palm oil in a Cameroonian region is described by:
\[
Q_d = 2000 - 100P, \qquad Q_s = 200P - 400
\]
where \(Q\) is in tonnes per month and \(P\) is in thousands of FCFA per tonne.
The government imposes a specific tax of \SI{2000}{FCFA} per tonne (i.e., \(t = 2\) in our units) on producers.
\begin{enumerate}
    \item Find the initial equilibrium price and quantity.
    \item Find the new equilibrium price paid by buyers (\(P_b\)), the price received by sellers (\(P_s\)), and the quantity.
    \item Calculate the tax burden on buyers and on sellers, and interpret using elasticities.
    \item Calculate the government tax revenue and the deadweight loss.
\end{enumerate}
\end{exercice}

\begin{corrige}[Exercise: Tax Incidence Calculation]
\begin{enumerate}
    \item \textbf{Initial equilibrium:} set \(Q_d = Q_s\):
    \[
    2000 - 100P = 200P - 400 \Rightarrow 300P = 2400 \Rightarrow P_0 = 8 \ (\SI{8000}{FCFA}/\text{tonne}).
    \]
    \(Q_0 = 2000 - 100(8) = 1200\) tonnes.
    \item \textbf{New supply with tax:} producers receive \(P_s = P_b - 2\), so
    \[
    Q_s = 200(P_b - 2) - 400 = 200P_b - 800.
    \]
    New equilibrium: \(2000 - 100P_b = 200P_b - 800 \Rightarrow 300P_b = 2800 \Rightarrow P_b = \dfrac{28}{3} \approx 9.33\) (\(\approx \SI{9333}{FCFA}\)).
    \[
    P_s = P_b - 2 = \dfrac{22}{3} \approx 7.33 \ (\approx \SI{7333}{FCFA}); \qquad Q_1 = 2000 - 100 \times \dfrac{28}{3} = \dfrac{3200}{3} \approx 1066.7 \text{ tonnes}.
    \]
    \item \textbf{Burdens:}
    \begin{itemize}
        \item Buyers' burden \(= P_b - P_0 = \dfrac{28}{3} - 8 = \dfrac{4}{3} \approx 1.33\) (\(\approx \SI{1333}{FCFA}\)).
        \item Sellers' burden \(= P_0 - P_s = 8 - \dfrac{22}{3} = \dfrac{2}{3} \approx 0.67\) (\(\approx \SI{667}{FCFA}\)).
    \end{itemize}
    Buyers bear two-thirds of the tax and sellers one-third. \textbf{Interpretation:} at the initial equilibrium, \(\eta_d = -100 \times \dfrac{8}{1200} \approx -0.67\) and \(\eta_s = 200 \times \dfrac{8}{1200} \approx 1.33\). Demand is more inelastic than supply, so buyers bear the larger share — confirming the general rule.
    \item \textbf{Revenue:} \(R = t \times Q_1 = 2 \times \dfrac{3200}{3} = \dfrac{6400}{3} \approx 2133.3\) (thousand FCFA \(\times\) tonnes, i.e., about 2.13 billion FCFA per month).
    \textbf{Deadweight loss:}
    \[
    DWL = \frac{1}{2} \times t \times (Q_0 - Q_1) = 0.5 \times 2 \times \left(1200 - \frac{3200}{3}\right) = \frac{400}{3} \approx 133.3 \text{ (about 133 million FCFA)}.
    \]
\end{enumerate}
\end{corrige}

\courssub{Key Points to Remember}

\begin{aretenir}[Summary: Economic Consequences of Taxation]
\begin{itemize}
    \item \textbf{Disposable income:} \(Y_d = Y - T_d + Tr\). Taxes reduce \(Y_d\) and thus consumption (by \(\text{MPC} \times \Delta T\)).
    \item \textbf{Tax wedge:} drives a wedge between the labour cost to the employer and the net wage to the employee; a high wedge $\rightarrow$ informal employment, less hiring.
    \item \textbf{Saving/Investment:} taxes on capital income raise the user cost of capital and reduce capital formation; incentives can offset this.
    \item \textbf{Redistribution:} progressive direct taxes reduce inequality; indirect taxes are often regressive unless basic goods are exempted.
    \item \textbf{Incidence:} legal incidence \(\neq\) economic incidence. The burden falls mainly on the more inelastic side of the market. Use the supply-demand diagram to analyse it.
    \item \textbf{Laffer curve:} revenue rises then falls as the tax rate increases. A useful conceptual tool (GCE extension), but the peak rate is unknown.
    \item \textbf{Net effects:} taxes cause a deadweight loss (allocation cost) but finance public goods, redistribution and stabilisation. Good tax design minimises distortions for a given revenue.
\end{itemize}
\end{aretenir}

\begin{attention}[Common GCE Traps]
\begin{itemize}
    \item \textbf{Confusing legal and economic incidence.} Always state: "The legal incidence is on X, but the economic incidence falls mainly on Y because..."
    \item \textbf{Forgetting the deadweight loss.} A tax always creates a DWL (except when demand or supply is perfectly inelastic). Draw the triangle!
    \item \textbf{Mis-stating the Laffer curve.} It does \emph{not} say that tax cuts always raise revenue. It says revenue is zero at 0\% and at 100\%, with a maximum in between. Most economies are on the rising side.
    \item \textbf{Ignoring the difference between average and marginal tax rates.} The marginal rate affects incentives (the next hour of work, the next franc of investment); the average rate measures the overall burden.
    \item \textbf{Sign errors in incidence calculations.} With a tax on sellers, replace \(P\) by \(P_b - t\) in the \emph{supply} function only; the demand function keeps \(P_b\).
\end{itemize}
\end{attention}

\coursec{The State Budget and the Programmed Budget}

In the previous section, we analysed how taxation — the State's primary revenue tool — affects economic behaviour, resource allocation and income distribution. We saw that taxes reduce disposable income, create deadweight losses, and can be used to correct market failures. But collecting revenue is only half the story. The State must also \emph{plan, authorise, execute and control} the spending of that revenue. This is the realm of \textbf{Public Budgeting}. A budget is not merely a spreadsheet; it is the legal instrument that translates political priorities into economic action. In Cameroon, as in all modern States, the budget is the heart of fiscal policy. Let us now dissect its anatomy, its governing principles, its cycle, and the modern shift toward \emph{programme budgeting}.

\courssub{Meaning and Legal Nature of the State Budget}

\begin{definition}[The State Budget]
The \textbf{State Budget} (or \textbf{National Budget}) is the \textbf{annual legal act} (an Act of Parliament) which \textbf{forecasts} and \textbf{authorises} the State's revenue and expenditure for a given financial year. It is the quantitative expression of the government's economic and social policy.
\end{definition}

\begin{aretenir}[Key Characteristics]
\begin{itemize}
\item \textbf{Annual:} It covers a specific financial year (in Cameroon: 1\textsuperscript{st} January – 31\textsuperscript{st} December).
\item \textbf{Legal:} It is voted as a \textbf{Finance Law} (\emph{Loi de Finances}) by Parliament. Without this law, the Executive cannot legally collect taxes or spend public funds.
\item \textbf{Authorisatory:} It sets \emph{limits} (appropriations) for expenditure and \emph{estimates} for revenue. The Executive cannot exceed expenditure ceilings without a supplementary finance law (\emph{Loi de Finances Rectificative}).
\item \textbf{Forecasting:} Revenue figures are estimates; expenditure figures are legal authorisations (upper bounds).
\end{itemize}
\end{aretenir}

\begin{definition}[The Finance Law (\emph{Loi de Finances})]
In Cameroon, the Finance Law is the legal instrument that enacts the budget. It comprises two main parts:
\begin{itemize}
    \item \textbf{The Revenue Section:} Provisions on tax rates, bases, customs duties, non-tax revenue, and borrowing authorisations.
    \item \textbf{The Expenditure Section:} Appropriations (\emph{dotations}) allocated to each Ministry/Institution, broken down by \emph{recurrent} and \emph{capital} expenditure.
\end{itemize}
\end{definition}

\courssub{Classical Budgetary Principles}

These principles, rooted in 19\textsuperscript{th}-century French and British constitutional practice, guarantee transparency, parliamentary control, and sound financial management. They are enshrined in Cameroon's Constitution (Article 27) and the \emph{Loi n\textsuperscript{o} 2018/012 du 11 juillet 2018} (Financial Regime of the State).

\begin{propriete}[The Five Classical Budgetary Principles]
\begin{enumerate}
    \item \textbf{Annuality (\emph{Annualité}):} The budget is voted for \textbf{one year only} (Jan 1 – Dec 31). It prevents the Executive from committing funds indefinitely without parliamentary renewal. \textit{Exception:} Multi-annual investment programmes exist, but annual appropriations are voted yearly.
    \item \textbf{Unity (\emph{Unité}):} All State revenue and expenditure must appear in \textbf{a single document} (the General Budget). No "off-budget" accounts or special treasury accounts (except those explicitly authorised by law, e.g., Autonomous Sinking Fund \emph{CAA}). This ensures a global view of public finances.
    \item \textbf{Universality (\emph{Universalité}):} \textbf{All revenue finances all expenditure} (no earmarking / \emph{affectation spéciale}). Tax revenue goes into a common pot ("Treasury Single Account") and funds the general budget. \textit{Exception:} Certain specific taxes may be earmarked by law (e.g., Petroleum royalties to producing regions, VAT share to communes).
    \item \textbf{Equilibrium / Balance (\emph{Équilibre}):} \textbf{Estimated Revenue = Authorised Expenditure.} A balanced budget is the ideal. In practice, deficits are frequent and legal, but they must be \emph{financed} (see Section 5).
    \item \textbf{Sincerity (\emph{Sincérité}):} Revenue estimates must be \textbf{realistic} (neither over- nor under-estimated) and expenditure must reflect \textbf{true costs}. "Creative accounting" (hiding deficits, inflating revenue) violates this principle.
\end{enumerate}
\end{propriete}

\begin{attention}[Exam Trap: Principles vs. Reality]
\textbf{Do not confuse the \emph{principle} with the \emph{practice}.}
\begin{itemize}
    \item \textbf{Equilibrium} is a principle, but \textbf{deficits are legal} if financed by borrowing (Article 27 of Constitution). A "balanced budget" in the exam means \emph{Revenue = Expenditure} in the \emph{voted} Finance Law.
    \item \textbf{Unity} is often violated by \emph{Comptes d'Affectation Spéciale} (Special Treasury Accounts) and extra-budgetary funds. Know the exceptions!
    \item \textbf{Annuality} does not forbid multi-year \emph{planning} (MTEF — Medium-Term Expenditure Framework), only multi-year \emph{appropriation} without annual vote.
\end{itemize}
\end{attention}

\courssub{The Budgetary Process in Cameroon: Four Stages}

The budget follows a strict annual cycle involving the Executive (Government) and the Legislature (Parliament), with ex-post control by the Audit Bench (\emph{Chambre des Comptes}) of the Supreme Court.

\begin{popfigure}
\centering
\begin{tikzpicture}[
    node distance=1.8cm and 2.5cm,
    stage/.style={rectangle, draw=popblue, thick, fill=popblueL, rounded corners, minimum width=3.2cm, minimum height=1.4cm, text width=3cm, align=center, font=\small\bfseries},
    arrow/.style={-Stealth, thick, popdark},
    detail/.style={rectangle, draw=popgreen, fill=popgreenL, rounded corners, minimum width=3cm, minimum height=1cm, text width=2.8cm, align=center, font=\footnotesize, opacity=0.9}
]
% Nodes
\node[stage, align=center] (prep){1. PREPARATION\\ (Executive)};
\node[stage, right=of prep, align=center] (vote){2. VOTE\\ (Parliament)};
\node[stage, right=of vote, align=center] (exec){3. EXECUTION\\ (Executive)};
\node[stage, right=of exec, align=center] (ctrl){4. CONTROL\\ (Audit Bench / Parliament)};

% Details under Preparation
\node[detail, below=0.8cm of prep, align=center] (p1){Circular from Minfi\\ (May)};
\node[detail, below=0.8cm of p1, align=center] (p2){Ministries submit\\ estimates};
\node[detail, below=0.8cm of p2, align=center] (p3){Arbitration / MTEF\\ alignment};
\node[detail, below=0.8cm of p3, align=center] (p4){Draft Finance Law\\ (Oct/Nov)};

% Details under Vote
\node[detail, below=0.8cm of vote, align=center] (v1){Tabling at National\\ Assembly (Nov)};
\node[detail, below=0.8cm of v1, align=center] (v2){Committee scrutiny\\ (Finance Commission)};
\node[detail, below=0.8cm of v2] (v3) {Plenary debate \& vote};
\node[detail, below=0.8cm of v3] (v4) {Senate review};
\node[detail, below=0.8cm of v4, align=center] (v5){Promulgation by\\ President (Dec 31)}; 

% Details under Execution
\node[detail, below=0.8cm of exec] (e1) {Cash management plan};
\node[detail, below=0.8cm of e1] (e2) {Commitment (\emph{Engagement})};
\node[detail, below=0.8cm of e2] (e3) {Verification (\emph{Visa})};
\node[detail, below=0.8cm of e3] (e4) {Payment order (\emph{Ordonnancement})};
\node[detail, below=0.8cm of e4] (e5) {Payment (\emph{Paiement})};
\node[detail, below=0.8cm of e5, align=center] (e6){Supplementary Budget\\ if needed};

% Details under Control
\node[detail, below=0.8cm of ctrl, align=center] (c1){Administrative control\\ (Minfi, Contrôle Financier)};
\node[detail, below=0.8cm of c1, align=center] (c2){Parliamentary control\\ (Finance Laws review)};
\node[detail, below=0.8cm of c2] (c3) {Audit Bench (\emph{Chambre des Comptes})};
\node[detail, below=0.8cm of c3, align=center] (c4){Certification of\\ Settlement Act (\emph{Loi de Règlement})};

% Arrows main flow
\draw[arrow] (prep) -- (vote);
\draw[arrow] (vote) -- (exec);
\draw[arrow] (exec) -- (ctrl);
% Feedback loop
\draw[arrow, poporange, bend right=30] (ctrl.south west) to node[midway, below, font=\footnotesize\itshape, poporange] {Feedback for next year} (prep.south west);

% Connect details (optional visual lines)
\foreach \x in {p1,p2,p3,p4} \draw[popline, thin] (prep.south) -- (\x.north);
\foreach \x in {v1,v2,v3,v4,v5} \draw[popline, thin] (vote.south) -- (\x.north);
\foreach \x in {e1,e2,e3,e4,e5,e6} \draw[popline, thin] (exec.south) -- (\x.north);
\foreach \x in {c1,c2,c3,c4} \draw[popline, thin] (ctrl.south) -- (\x.north);
\end{tikzpicture}
\legende{The Four Stages of the Budgetary Process in Cameroon}
\end{popfigure}

\begin{methode}[The Expenditure Chain (\emph{La Chaîne de la Dépense}) — Examinable Steps]
During \textbf{Execution (Stage 3)}, every payment follows a legal chain. Missing a step makes the expenditure irregular.
\begin{enumerate}
    \item \textbf{Engagement (Commitment):} The Minister (\emph{Ordonnateur}) recognises the debt (e.g., signs a contract). \emph{Creates the obligation.}
    \item \textbf{Visa / Liquidation (Verification):} The Financial Controller (\emph{Contrôleur Financier}) checks: \textbf{Service fait} (service rendered), \textbf{Disponibilité budgétaire} (budget line has funds), \textbf{Régularité} (legal procurement). \emph{Authorises payment.}
    \item \textbf{Ordonnancement (Payment Order):} The Minister issues the payment order (\emph{Ordre de paiement}) to the Public Accountant.
    \item \textbf{Paiement (Payment):} The Public Accountant (\emph{Comptable Public} — Paymaster) pays the beneficiary. \emph{Extinguishes the debt.}
\end{enumerate}
\textbf{Separation of functions:} \emph{Ordonnateur} (authorises/commits) $\neq$ \emph{Comptable} (pays/handles cash). This prevents fraud.
\end{methode}

\courssub{Structure of the Budget: Revenue and Expenditure}

The Finance Law presents two columns: \textbf{Revenue (Recettes)} and \textbf{Expenditure (Dépenses)}.

\begin{definition}[Budget Revenue (Recettes Budgétaires)]
\begin{itemize}
    \item \textbf{Tax Revenue (\emph{Recettes fiscales}):} Direct taxes (Income tax, Company tax), Indirect taxes (VAT, Excise duties, Customs duties). \emph{Main source (approx. 60--70\%)}.
    \item \textbf{Non-Tax Revenue (\emph{Recettes non fiscales}):} Domain revenue (rents, dividends from State enterprises like SONARA, SNH), Administrative fees, Fines.
    \item \textbf{Capital Revenue (\emph{Recettes en capital}):} Privatisation proceeds, Loan repayments received, Grants (\emph{Dons}) from donors.
    \item \textbf{Borrowing (\emph{Emprunts}):} Treasury bills/bonds (domestic), Project loans / Programme loans (external). \emph{Not "revenue" in economic sense, but financing.}
\end{itemize}
\end{definition}

\begin{definition}[Budget Expenditure (Dépenses Budgétaires)]
Classified by \textbf{Economic Nature} (most important for analysis) and \textbf{Administrative Nature} (Ministry).
\begin{itemize}
    \item \textbf{Recurrent Expenditure (\emph{Dépenses de fonctionnement / Courantes}):} Consumption spending. Wages (\emph{Personnel}), Goods \& Services (\emph{Fonctionnement}), Transfers (\emph{Transferts} — subsidies, pensions, debt interest), Subsidies to public enterprises (SONARA, CAMWATER, etc.).
    \item \textbf{Capital / Investment Expenditure (\emph{Dépenses d'investissement / En capital}):} Spending that creates assets. Infrastructure (roads, dams, schools), Equipment, Acquisition of shares, Capital transfers.
\end{itemize}
\end{definition}

\begin{aretenir}[Cameroon Context: The "Wage Bill" Burden]
In Cameroon's recent budgets (2020–2024), \textbf{Recurrent expenditure} typically absorbs \textbf{70--80\%} of total expenditure, driven by the \textbf{wage bill} (\emph{Masse salariale}) and \textbf{debt service} (\emph{Service de la dette}). This leaves a narrow fiscal space for \textbf{Capital expenditure} (often < 20--25\%), constraining growth-enhancing investment. The IMF-supported programmes often target a reduction in the wage bill-to-revenue ratio.
\end{aretenir}

\begin{popfigure}
\centering
\begin{tikzpicture}
\begin{axis}[
    ybar stacked,
    bar width=1.8cm,
    width=0.9\linewidth,
    height=10cm,
    enlarge x limits=0.3,
    ylabel={Amount (Billions $\mathrm{CFAF}$)},
    ylabel style={font=\small},
    xtick={1,2},
    xticklabels={Revenue, Expenditure},
    xtick style={font=\small},
    ytick style={font=\small},
    ymin=0, ymax=6500,
    legend style={at={(0.5,-0.15)}, anchor=north, legend columns=3, font=\footnotesize, /tikz/every even column/.append style={column sep=0.5cm}},
    grid=major, grid style={popline},
    axis lines=left,
    axis line style={popdark},
    tick style={popdark},
]
% Revenue bars (2023/2024 approx figures for illustration)
\addplot[fill=popblue] coordinates {(1, 2800)}; % Tax
\addplot[fill=popblue!60!popgreen] coordinates {(1, 400)}; % Non-tax
\addplot[fill=poporange] coordinates {(1, 600)}; % Grants/Loans (Financing)
\addplot[fill=poppurple] coordinates {(1, 200)}; % Capital revenue/Privatization

% Expenditure bars
\addplot[fill=poppink] coordinates {(2, 2200)}; % Recurrent (Wages+Goods+Transfers+Interest)
\addplot[fill=poporange] coordinates {(2, 1200)}; % Capital/Investment

\legend{Tax Revenue, Non-Tax Revenue, Grants \& Loans (Financing), Capital Revenue, Recurrent Expenditure, Capital Expenditure}
\end{axis}
\end{tikzpicture}
\legende{Simplified Structure of Cameroon's Budget (Illustrative 2023/24 Figures in Billions $\mathrm{CFAF}$). Note: Deficit financed by Grants/Loans.}
\end{popfigure}

\courssub{Budget Balance: Surplus, Deficit and Financing}

\begin{definition}[Budget Balance]
The \textbf{Budget Balance} ($B$) is the difference between \textbf{Total Revenue ($R$)} and \textbf{Total Expenditure ($E$)} (excluding borrowing):
\[
B = R - E
\]
\begin{itemize}
    \item \textbf{Budget Surplus ($B > 0$):} Revenue exceeds expenditure. Rare for developing States. Allows debt repayment or reserve accumulation.
    \item \textbf{Balanced Budget ($B = 0$):} Revenue equals expenditure. The classical ideal.
    \item \textbf{Budget Deficit ($B < 0$):} Expenditure exceeds revenue. The norm for Cameroon and most countries.
\end{itemize}
\end{definition}

\begin{definition}[Financing the Deficit (\emph{Financement du Déficit})]
A deficit \textbf{must} be financed. The \textbf{Financing Requirement} = Deficit + Principal Repayments (Amortisation). Sources:
\begin{enumerate}
    \item \textbf{Domestic Borrowing:} Treasury Bills (\emph{Bons du Trésor} — short term), Treasury Bonds (\emph{Obligations du Trésor} — medium/long term) issued on CEMAC market (BVMAC). \emph{Crowding out risk.}
    \item \textbf{External Borrowing:}
        \begin{itemize}
            \item \textbf{Project Loans:} Tied to specific infrastructure (e.g., Kribi Port, Memve'ele Dam). Concessional (low interest, long grace) or Non-concessional.
            \item \textbf{Programme Loans / Budget Support:} Untied, for reform programmes (IMF, World Bank, AfDB, AFD).
        \end{itemize}
    \item \textbf{Monetary Financing (\emph{Financement monétaire / Planche à billets}):} Central Bank (BEAC) directly credits Treasury account. \textbf{Strictly limited by CEMAC Treaty (Statut de la BEAC):} No direct advances; only via market operations. Excess leads to inflation and loss of reserves.
    \item \textbf{Use of Reserves / Deposits:} Drawing down Treasury deposits at BEAC or commercial banks.
    \item \textbf{Privatisation / Asset Sales:} One-off capital revenue.
\end{enumerate}
\end{definition}

\begin{exoresolu}[Budget Balance Calculation]
\textbf{Statement:} The 2024 Finance Law of a hypothetical State \emph{X} forecasts: Tax Revenue = 3,200; Non-Tax Revenue = 450; Grants = 300; Recurrent Expenditure = 3,100; Capital Expenditure = 1,200. Principal debt repayment (amortisation) = 400.
\begin{enumerate}
    \item Calculate the \textbf{Budget Balance} (Overall Balance).
    \item Calculate the \textbf{Financing Requirement} (\emph{Besoin de Financement}).
    \item If the State issues Treasury Bills for 500 and receives a Programme Loan of 300, what is the remaining gap?
\end{enumerate}
\tcblower
\textbf{Solution:}
\begin{enumerate}
    \item \textbf{Total Revenue ($R$)} = Tax + Non-Tax + Grants = 3,200 + 450 + 300 = \textbf{3,950}.\\
    \textbf{Total Expenditure ($E$)} = Recurrent + Capital = 3,100 + 1,200 = \textbf{4,300}.\\
    \textbf{Budget Balance ($B$)} = $R - E$ = 3,950 - 4,300 = \textbf{-350 (Deficit of 350)}.
    \item \textbf{Financing Requirement} = Deficit + Amortisation = 350 + 400 = \textbf{750}.\\
    (This is the total new borrowing + drawdowns needed).
    \item \textbf{Financing secured} = T-Bills (500) + Programme Loan (300) = 800.\\
    \textbf{Remaining Gap} = 750 - 800 = \textbf{+50 (Surplus financing / Reserve build-up)}.\\
    \textit{Note: In practice, the Finance Law authorises a borrowing ceiling. If financing > requirement, the excess stays as Treasury deposits.}
\end{enumerate}
\end{exoresolu}

\courssub{The Programmed Budget (Budget-Programme / Performance Budgeting)}

Since the 2000s, Cameroon has undertaken a major budgetary reform (supported by IMF/World Bank) to move from a \textbf{Means-Based Budget} (\emph{Budget de Moyens}) to a \textbf{Programme Budget} (\emph{Budget-Programme}).

\begin{definition}[Programme Budget (\emph{Budget-Programme})]
A \textbf{Programme Budget} is a budget presentation where appropriations are allocated not just to \emph{input categories} (wages, fuel, buildings) nor merely to \emph{administrative units} (Ministries), but to \textbf{Programmes} linked to \textbf{Objectives}, \textbf{Outputs} (\emph{Produits}) and \textbf{Performance Indicators} (\emph{Indicateurs de performance}).
\end{definition}

\begin{center}
\begin{tabularx}{\linewidth}{|l|X|X|}
\hline
\rowcolor{popblueL}
\textbf{Dimension} & \textbf{Classical (Means) Budget} & \textbf{Programme Budget} \\
\hline
\textbf{Classification} & Economic nature (Wages, Investment) + Administrative (Ministry) & \textbf{Programmes $\to$ Actions $\to$ Activities} \\
\hline
\textbf{Focus} & \textbf{Inputs} (What we buy: litres of fuel, number of civil servants) & \textbf{Results} (What we achieve: km of road built, vaccination rate, exam pass rate) \\
\hline
\textbf{Question asked} & "How much did Ministry X spend?" & "What did Programme Y achieve with the money?" \\
\hline
\textbf{Accountability} & Financial regularity (Did they follow procurement rules?) & \textbf{Performance} (Did they deliver outcomes efficiently?) \\
\hline
\textbf{Management} & Compliance / Input control & \textbf{Results-Based Management (RBM) / \emph{Gestion Axée sur les Résultats (GAR)}} \\
\hline
\end{tabularx}
\end{center}

\begin{propriete}[Structure of a Programme Budget (Cameroon Format)]
\begin{enumerate}
    \item \textbf{Policy Area (\emph{Secteur de politique publique}):} e.g., Education, Health, Transport.
    \item \textbf{Programme:} A coherent set of actions for a major objective. \textit{Ex: "Programme 142: Access to Quality Basic Education".}
    \item \textbf{Action (\emph{Action}):} A homogeneous set of activities. \textit{Ex: "Action 1421: Construction and Equipment of Classrooms".}
    \item \textbf{Performance Indicators:}
        \begin{itemize}
            \item \textbf{Output Indicators (\emph{Indicateurs de produits}):} Quantity/Quality of goods/services delivered. \textit{Ex: "Number of classrooms built", "Pupil/Classroom ratio".}
            \item \textbf{Outcome/Impact Indicators (\emph{Indicateurs d'effets/impact}):} Changes in target population. \textit{Ex: "Primary completion rate", "Literacy rate".}
        \end{itemize}
    \item \textbf{Targets (\emph{Cibles}):} Quantified goals for the year (and medium term).
    \item \textbf{Costing:} Total cost of the Programme = Sum of costs of Actions (using standard costs).
\end{enumerate}
\end{propriete}

\begin{aretenir}[Advantages of Programme Budgeting (Essay Points)]
\begin{itemize}
    \item \textbf{Efficiency / Allocative Efficiency:} Resources shift to programmes with highest impact/cost ratio.
    \item \textbf{Transparency \& Accountability:} Parliament and citizens can judge "value for money", not just legality.
    \item \textbf{Strategic Alignment:} Links budget to \emph{National Development Strategy (SND30)} and \emph{MTEF (Cadre de Dépenses à Moyen Terme)}.
    \item \textbf{Incentive for Managers:} Programme managers (\emph{Gestionnaires de programmes}) have flexibility to reallocate inputs \emph{within} the programme to hit targets (virement).
    \item \textbf{Better Monitoring/Evaluation:} Ex-post evaluation (\emph{Évaluation des politiques publiques}) becomes possible.
\end{itemize}
\end{aretenir}

\begin{attention}[Limits and Challenges in Cameroon]
\begin{itemize}
    \item \textbf{Data Quality:} Reliable baseline data and monitoring systems (statistics) are often weak.
    \item \textbf{Cultural Shift:} Civil servants used to "compliance" (following rules) resist "performance" (taking risks for results).
    \item \textbf{Rigidities:} Wage bill (70\%+ of recurrent) is largely \emph{incompressible} and not easily "programmable" by results.
    \item \textbf{Complexity:} Defining good indicators (SMART: Specific, Measurable, Achievable, Relevant, Time-bound) is technically hard.
    \item \textbf{Partial Implementation:} Cameroon's budget is still largely presented by Ministry/Economic Nature; the "Programme" annex exists but does not yet fully drive appropriations.
\end{itemize}
\end{attention}

\courssub{The State Budget of Cameroon: Recent Orientations (SND30 Era)}

Cameroon's budget is framed by the \textbf{National Development Strategy 2020–2030 (SND30)}, which aims at structural transformation and import substitution. The annual Finance Law operationalises the \textbf{MTEF (Cadre de Dépenses à Moyen Terme — CDMT)}.

\begin{savaistu}[Cameroon Budget Key Figures (Illustrative Trends 2021–2024)]
\begin{itemize}
    \item \textbf{Total Budget Size:} Grew from $\sim$4,900 Billion (2021) to $\sim$5,800–6,000 Billion $\mathrm{CFAF}$ (2024).
    \item \textbf{Revenue Mobilisation:} Tax pressure (\emph{Pression fiscale}) remains low ($\sim$11--12\% of GDP), far below the 20\%+ target. Non-oil revenue effort is priority (VAT, Digital economy taxation).
    \item \textbf{Deficit:} Overall deficit (commitment basis) targeted around 2--3\% of GDP (CEMAC convergence criterion $\le$ 3\%), but often wider on cash basis due to arrears (\emph{arriérés}).
    \item \textbf{Debt Service:} Rising fast (Eurobond 2015/2021 repayments, BEAC statutory advances repayment). Crowds out investment.
    \item \textbf{Priority Sectors (SND30):} Agriculture (import substitution: rice, maize, fish), Energy (hydro/solar), Transport (roads, ports), Digital, Human Capital (Health/Education).
    \item \textbf{Decentralisation:} Increasing share of \emph{Dotation Générale de Décentralisation (DGD)} and \emph{Dotation d'Investissement} transferred to Communes/Regions (Law 2019/024).
\end{itemize}
\end{savaistu}

\begin{experience}[Case Study: Reading a Finance Law Table (Simplified)]
\textbf{Extract from a hypothetical Finance Law (Billions $\mathrm{CFAF}$):}

\begin{center}
\begin{tabularx}{\linewidth}{|l|r|r|>{\centering\arraybackslash}X|}
\hline
\rowcolor{popblueL}
\textbf{Programme / Action} & \textbf{Recurrent} & \textbf{Capital} & \textbf{Key Performance Indicator (Target 2024)} \\
\hline
\textbf{PROG 142: Basic Education} & \textbf{450} & \textbf{85} & \\
\quad Action 1421: Classroom Construction & 5 & 60 & Classrooms built: 1,200 \\
\quad Action 1422: Textbook Distribution & 20 & 0 & Pupil/Textbook ratio: 1:1 \\
\quad Action 1423: Teacher Recruitment/Training & 400 & 5 & Qualified teacher rate: 85\% \\
\quad Action 1424: School Feeding (Cantines) & 25 & 20 & Beneficiary pupils: 500,000 \\
\hline
\textbf{PROG 210: Road Infrastructure} & \textbf{30} & \textbf{420} & \\
\quad Action 2101: National Roads Maintenance & 20 & 150 & Km maintained: 5,000 \\
\quad Action 2102: New Construction (Bamenda-Mamfe) & 0 & 200 & Km paved: 80 \\
\quad Action 2103: Rural Roads & 10 & 70 & Km opened/rehabilitated: 1,500 \\
\hline
\end{tabularx}
\end{center}

\textbf{Analysis Questions (Exam Style):}
\begin{enumerate}
    \item Identify the \textbf{recurrent/capital split} for each programme. \textit{(Education: 84\% Recurrent / 16\% Capital; Roads: 7\% Recurrent / 93\% Capital).}
    \item Why is Education recurrent-heavy? \textit{(Wage bill for teachers dominates).}
    \item Why is Roads capital-heavy? \textit{(Asset creation).}
    \item If the wage bill rises by 10\%, which programme suffers most? \textit{(Education — less room for textbooks/training).}
    \item How does \textbf{Programme Budgeting} help the Minister of Finance arbitrate? \textit{(He sees \emph{cost per classroom} or \emph{cost per km} and can compare efficiency across years/regions).}
\end{enumerate}
\end{experience}

\courssub{Synthesis and Exam Focus}

\begin{aretenir}[Checklist for GCE Advanced Level]
\begin{itemize}
\item \textbf{Define:} Budget, Finance Law, Programme Budget, Budget Balance (Surplus/Deficit), Financing Requirement.
\item \textbf{List \& Explain:} 5 Classical Principles (Annuality, Unity, Universality, Equilibrium, Sincerity) — give Cameroon exceptions.
\item \textbf{Describe:} 4 Stages of Budgetary Process (Preparation $\to$ Vote $\to$ Execution $\to$ Control). Know the actors (Minfi, Parliament, Audit Bench, Financial Controller, Paymaster).
\item \textbf{Distinguish:} Recurrent vs Capital Expenditure; Tax vs Non-Tax Revenue; Domestic vs External Financing.
\item \textbf{Calculate:} Budget Balance, Financing Requirement, Financing Gap (Worked Example style).
\item \textbf{Evaluate:} Advantages/Limits of Programme Budgeting (RBM). Why is Cameroon's capital expenditure share low? (Wage bill, Debt service).
\item \textbf{Contextualise:} Mention SND30, MTEF (CDMT), CEMAC Convergence Criteria (Deficit $\le$ 3\% GDP, Debt $\le$ 70\% GDP, Inflation $\le$ 3\%, Non-accumulation of arrears).
\end{itemize}
\end{aretenir}

\begin{exercice}[Practice Questions]
\begin{enumerate}
    \item \textbf{(Essay)} "The principle of budgetary equilibrium is more theoretical than practical in Cameroon." Discuss. (Structure: Define principle $\to$ Show legal possibility of deficit $\to$ Explain financing methods $\to$ Link to debt sustainability/CEMAC criteria $\to$ Conclude).
    \item \textbf{(Data Response)} Given a table of Revenue/Expenditure for 3 years, calculate: Overall Balance, Primary Balance (Excl. Interest), Financing Requirement. Comment on trend.
    \item \textbf{(Short Answer)} Explain the role of the \emph{Contrôleur Financier} in the expenditure chain. Why is the separation \emph{Ordonnateur/Comptable} crucial?
    \item \textbf{(Application)} A Ministry wants to build a hospital. Trace the steps from \emph{Budget Preparation} to \emph{Payment of Contractor}, naming the responsible actor at each stage.
    \item \textbf{(Evaluation)} "Programme budgeting solves the problem of low capital execution rates in Cameroon." Critically assess.
\end{enumerate}
\end{exercice}

\begin{corrige}[Exercise 1 — Outline]
\textbf{Introduction:} Define Budgetary Equilibrium ($R=E$). State it is a constitutional principle (Art 27).
\textbf{Body 1 (Theory):} Classical view: Deficit = bad management. Balanced budget = sound finance.
\textbf{Body 2 (Practice):} Keynesian view / Modern reality: Deficit = tool of fiscal policy (stimulus). Cameroon: Structural deficits due to low revenue, rigid recurrent spending, investment needs.
\textbf{Body 3 (Legal Framework):} Deficit legal if financed by borrowing (Art 27). CEMAC convergence criterion: Deficit $\le$ 3\% GDP.
\textbf{Body 4 (Cameroon Evidence):} Recent deficits (2-4\% GDP). Financing: Eurobonds, BEAC advances (statutory), Treasury Bills, Project loans. Risk: Debt service crowding out investment.
\textbf{Body 5 (Limits):} "Equilibrium" in voted law $\neq$ "Equilibrium" in execution (arrears, revenue shortfalls). Cash basis vs Commitment basis.
\textbf{Conclusion:} Principle remains a benchmark (discipline), but flexible application is necessary for development. Credibility depends on \emph{quality} of deficit (investment vs consumption) and debt sustainability.
\end{corrige}

\begin{methode}[How to Answer "Distinguish between Recurrent and Capital Expenditure" (4-6 marks)]
\begin{enumerate}
    \item \textbf{Definition:} Recurrent = consumption, repeats yearly, does not create asset. Capital = investment, creates durable asset, one-off/large.
    \item \textbf{Examples (Cameroon):} Recurrent: Salaries, Fuel, Rent, Subsidies to SONARA, Debt Interest. Capital: Building Kribi Port, Buying MRI machines, Constructing classrooms, Equity in SONAMINES.
    \item \textbf{Economic Impact:} Recurrent $\to$ Maintains current welfare/administration. Capital $\to$ Expands productive capacity (AS shift), future growth.
    \item \textbf{Budget Treatment:} Recurrent financed by recurrent revenue (principle). Capital can be financed by borrowing (Golden Rule).
    \item \textbf{Indicator:} Capital/Total Expenditure ratio (Quality of budget).
\end{enumerate}
\end{methode}

\coursec{The Public Debt and the Effectiveness of Fiscal Policy in Cameroon}

In the previous section we studied the State budget, the annual document that authorises public revenue and expenditure, and the programmed budget which plans expenditure over several years. We saw that when expenditure exceeds revenue, the budget shows a \cle{deficit}. But a deficit must be financed somehow: the State borrows. Year after year, these borrowings accumulate into what is called the \cle{public debt}. This final section studies the public debt --- its meaning, forms, measurement and consequences --- and then evaluates the effectiveness of \cle{fiscal policy}, the use of the budget to steer the economy, in Cameroon.

\courssub{Meaning and Causes of the Public Debt}

\textbf{Intuition.} A trader in Mokolo market who spends more than she earns must borrow from a \emph{tontine} or a microfinance institution. If she borrows again the following year without repaying the first loan, her total debt grows. The State behaves in exactly the same way: each budget deficit adds to the stock of accumulated borrowing.

\begin{definition}[Public debt and debt service]
The \cle{public debt} (or \cle{national debt}) is the \textbf{total of outstanding loans contracted by the State and public authorities} (central government, local councils, and public enterprises whose debt is guaranteed by the State) that have not yet been repaid. It is a \textbf{stock} measured at a given date.

The \cle{debt service} is the amount the State must pay each year on the debt: it is the sum of \textbf{interest payments} plus \textbf{repayment of principal} (amortisation) falling due in that year.
\end{definition}

\textbf{Causes of public debt.} The State borrows for several reasons:
\begin{itemize}
\item \textbf{Budget deficits:} when current revenue (taxes, oil receipts) is insufficient to cover expenditure, the gap is financed by borrowing.
\item \textbf{Financing of investment:} large infrastructure projects (dams, ports, roads such as the Yaound\'e--Douala motorway) cost more than one year's revenue, so the State borrows to spread the cost over time.
\item \textbf{External shocks:} a fall in world prices of oil, cocoa or coffee reduces export and fiscal revenue, forcing the State to borrow to maintain expenditure.
\item \textbf{Exchange-rate effects:} when external debt is denominated in foreign currency (US dollars, euros), an appreciation of that currency against the CFA franc automatically increases the debt burden measured in FCFA, even without any new borrowing.
\end{itemize}

\begin{attention}[Deficit versus debt: do not confuse them!]
The \textbf{budget deficit} is a \emph{flow}: it is the excess of expenditure over revenue during \textbf{one year}. The \textbf{public debt} is a \emph{stock}: it is the \textbf{accumulation of all past deficits} (minus repayments) at a given date. A country can reduce its deficit while its debt continues to rise, as long as the deficit remains positive. In examinations, candidates who write ``the debt is the annual deficit'' lose marks.
\end{attention}

\courssub{Forms of the Public Debt}

\begin{definition}[Internal and external debt]
The \cle{internal debt} is the part of the public debt owed to lenders \textbf{inside the country or the monetary zone}: domestic banks, the central bank (BEAC), and treasury bills and bonds bought by residents of the CEMAC zone.

The \cle{external debt} is the part owed to \textbf{foreign lenders}. It is classified by type of creditor:
\begin{itemize}
\item \textbf{Multilateral debt:} owed to international institutions (World Bank, IMF, African Development Bank);
\item \textbf{Bilateral debt:} owed to foreign governments (e.g.\ Paris Club members, China);
\item \textbf{Commercial debt:} owed to private foreign banks and bondholders (e.g.\ the Eurobonds issued by Cameroon).
\end{itemize}
\end{definition}

Debt is also classified by \textbf{maturity}: \textbf{short-term} (less than one year, e.g.\ treasury bills), \textbf{medium-term} (one to five years) and \textbf{long-term} (more than five years, e.g.\ bonds and infrastructure loans). Long-term debt is generally less risky because it gives the State time to generate returns from the investments financed.

\courssub{Measuring the Debt}

The absolute size of the debt in FCFA means little on its own; economists compare it with the size of the economy and with the State's capacity to pay.

\begin{propriete}[Debt indicators]
The \cle{debt-to-GDP ratio} measures the weight of the debt relative to national output:
\[ \text{Debt ratio} = \frac{\text{Public debt}}{\text{GDP}} \times 100 \]
The \cle{debt service ratio} compares annual debt service with government revenue or with export earnings; it measures how much of the State's resources are absorbed by past borrowing. Within the \textbf{CEMAC convergence criteria}, member States (including Cameroon) commit to keeping public debt \textbf{at or below 70\% of GDP}.
\end{propriete}

\begin{exemplebox}[Calculating a debt ratio]
Suppose in a given year Cameroon's public debt stands at 12\,000 billion FCFA and GDP at 25\,000 billion FCFA. Then:
\[ \text{Debt ratio} = \frac{12\,000}{25\,000} \times 100 = 48\% \]
The ratio of 48\% is below the CEMAC ceiling of 70\%, so the convergence criterion is respected --- although the \emph{trend} of the ratio and the cost of \emph{debt service} must also be watched.
\end{exemplebox}

\begin{popfigure}
\begin{center}
\begin{tikzpicture}
\begin{axis}[
  width=0.85\linewidth, height=6.2cm,
  xlabel={Year}, ylabel={Public debt (\% of GDP)},
  xmin=2000, xmax=2022, ymin=0, ymax=90,
  xtick={2000,2004,2008,2012,2016,2020},
  ytick={0,20,40,60,80},
  grid=both, grid style={popline, dashed},
  legend style={at={(0.02,0.97)}, anchor=north west, font=\small},
]
\addplot[very thick, popblue, mark=*, mark size=1.2pt] coordinates {
(2000,80) (2002,62) (2004,50) (2006,18) (2008,12) (2010,14) (2012,15) (2014,20) (2016,30) (2018,37) (2020,44) (2022,45)
};
\addlegendentry{Cameroon public debt (indicative trend)}
\addplot[thick, poporange, dashed] coordinates {(2000,70) (2022,70)};
\addlegendentry{CEMAC ceiling: 70\%}
\end{axis}
\end{tikzpicture}
\legende{Evolution of Cameroon's public debt as a percentage of GDP (indicative trend): very high before the 2006 HIPC debt relief, then rising again with infrastructure borrowing.}
\end{center}
\end{popfigure}

\courssub{Consequences of the Public Debt}

\textbf{Negative consequences:}
\begin{itemize}
\item \textbf{Crowding out of public expenditure:} debt service absorbs budgetary resources that could have financed health, education or investment. When debt service rises, less is left for classrooms and hospitals.
\item \textbf{Crowding out of private investment:} heavy borrowing on the domestic market absorbs savings and pushes interest rates up, making credit scarcer and dearer for private firms.
\item \textbf{Burden on future generations:} today's borrowing must be repaid by tomorrow's taxpayers; excessive debt mortgages the future.
\item \textbf{Risk of over-indebtedness:} beyond a certain level, lenders demand higher interest rates or refuse to lend, and the country may lose access to finance or face a debt crisis.
\end{itemize}

\textbf{Possible benefits:} borrowing is not bad in itself. If loans finance \textbf{productive investment} (roads, energy, ports) that raises future output and tax revenue, the debt is self-financing: growth makes repayment easier. The problem arises when borrowing finances current consumption or poorly chosen projects.

\textbf{Debt management and relief.} Over-indebted countries can negotiate a \cle{rescheduling} (postponement of repayments) or obtain \cle{debt relief} (partial or total cancellation). Under the \textbf{HIPC initiative} (Heavily Indebted Poor Countries), Cameroon, after implementing agreed reforms, reached the \textbf{completion point in 2006} and obtained substantial cancellation of its external debt --- visible in the figure above as the sharp fall of the ratio around 2006.

\begin{savaistu}[Cameroon and the HIPC initiative]
The debt relief obtained in 2006 freed several hundred billion FCFA of budgetary resources, part of which was redirected towards poverty-reduction expenditure. However, heavy borrowing for infrastructure since the 2010s has pushed the debt ratio up again, showing that debt relief is not a permanent cure: sound management must follow.
\end{savaistu}

\courssub{Fiscal Policy: Definition and Instruments}

\begin{definition}[Fiscal policy]
\cle{Fiscal policy} is the deliberate use of the State budget --- \textbf{taxation} and \textbf{public expenditure} --- by the government to influence the level of economic activity, employment, prices and growth. Its two instruments are:
\begin{itemize}
\item \textbf{Public expenditure} (investment, wages, transfers, subsidies);
\item \textbf{Taxation} (rates, bases, exemptions).
\end{itemize}
An \cle{expansionary fiscal policy} raises expenditure and/or cuts taxes to stimulate aggregate demand during a slowdown. A \cle{restrictive fiscal policy} cuts expenditure and/or raises taxes to cool down demand and fight inflation or reduce a deficit.
\end{definition}

\begin{popfigure}
\begin{center}
\begin{tikzpicture}[
  box/.style={draw=popblue, very thick, rounded corners, fill=popblueL, text width=3.4cm, align=center, minimum height=1.3cm, font=\small},
  arr/.style={-{Stealth[length=3mm]}, very thick, popdark}
]
\node[box, align=center] (a) at (0,0){\textbf{Fiscal impulse}\\ tax cut or higher public spending};
\node[box, align=center] (b) at (5.4,0){\textbf{Aggregate demand}\\ $AD = C+I+G+(X-M)$ rises};
\node[box, align=center] (c) at (10.8,0){\textbf{Output and employment}\\ firms produce more, hire more};
\node[box, fill=popgreenL, draw=popgreen, align=center] (d) at (5.4,-2.6){\textbf{Multiplier effect}\\ initial injection amplified through respending};
\draw[arr] (a) -- (b);
\draw[arr] (b) -- (c);
\draw[arr] (b) -- (d);
\draw[arr] (d) -- (c);
\end{tikzpicture}
\legende{The fiscal policy transmission chain: a budgetary stimulus raises aggregate demand, which raises output and employment through the multiplier.}
\end{center}
\end{popfigure}

\courssub{Further Studies 6: Effectiveness of Fiscal Policy in Cameroon}

\begin{methode}[Evaluating fiscal policy effectiveness]
For each aspect of the evaluation, follow four steps:
\begin{enumerate}
\item \textbf{State the objective} (e.g.\ raise growth, mobilise revenue, reduce poverty);
\item \textbf{Identify the instrument} used (tax reform, public investment programme);
\item \textbf{Assess the result} (what was actually achieved);
\item \textbf{State the limit} (why the result fell short of the objective).
\end{enumerate}
\end{methode}

\textbf{Achievements of fiscal policy in Cameroon:}
\begin{itemize}
\item \textbf{Revenue mobilisation:} modernisation of tax administration (larger taxpayer offices, computerisation, taxation of mobile money and telephone services) has increased non-oil revenue.
\item \textbf{Infrastructure:} expansionary budgets have financed dams, roads, ports and stadiums, supporting growth and employment.
\item \textbf{Counter-cyclical action:} during external shocks (fall in oil prices, security and health crises), budget deficits sustained demand and protected economic activity.
\end{itemize}

\textbf{Limits of fiscal policy in Cameroon:}
\begin{itemize}
\item \textbf{Narrow tax base:} a large informal sector escapes taxation, so tax rates weigh heavily on the few formal firms and workers.
\item \textbf{Tax evasion and fraud:} smuggling, under-invoicing and corruption leak revenue away from the Treasury.
\item \textbf{Debt burden:} rising debt service crowds out social and investment expenditure, reducing the room for manoeuvre.
\item \textbf{Oil dependence:} revenue fluctuates with world oil prices, making fiscal policy unstable and sometimes pro-cyclical.
\item \textbf{Weak multiplier:} part of public spending leaks into imports, limiting the impact on domestic output.
\end{itemize}

\begin{exoresolu}[Essay practice: fiscal policy in Cameroon]
\textbf{Statement.} ``Fiscal policy has been an effective instrument of economic management in Cameroon.'' Discuss.
\tcblower
\textbf{Detailed solution (plan and key arguments).}
\textbf{Introduction:} define fiscal policy (use of taxation and public expenditure to influence activity); announce a balanced discussion.
\textbf{Part 1 --- Achievements:} improved revenue mobilisation through administrative modernisation; major infrastructure financed by expansionary budgets; deficits used to cushion external shocks; debt relief in 2006 (HIPC) freed resources for poverty reduction.
\textbf{Part 2 --- Limitations:} narrow tax base and large informal sector; evasion and fraud; heavy and rising debt service; dependence on volatile oil revenue; expenditure leakages into imports weaken the multiplier.
\textbf{Part 3 --- Suggestions:} widen the tax base by formalising small businesses; strengthen tax administration and fight corruption; prioritise productive investment; diversify the economy to reduce oil dependence; keep debt within the CEMAC ceiling of 70\% of GDP.
\textbf{Conclusion:} fiscal policy in Cameroon has achieved real but partial success; its effectiveness depends on broadening the tax base and improving the quality of expenditure.
\end{exoresolu}

\begin{attention}[Exam tip]
For essay questions on Cameroon's fiscal policy, never give a one-sided answer. Structure your essay as: \textbf{achievements} $\rightarrow$ \textbf{limitations} $\rightarrow$ \textbf{suggestions}, with a definition of fiscal policy in the introduction and a balanced conclusion. Use concrete Cameroonian references (HIPC 2006, CEMAC 70\% criterion, VAT, informal sector) to score application marks.
\end{attention}

\begin{exercice}[Consolidation exercise]
\begin{enumerate}
\item Distinguish clearly, with one example each, between the budget deficit and the public debt.
\item A country has a public debt of 9\,000 billion FCFA and a GDP of 15\,000 billion FCFA. Calculate the debt ratio and state whether the CEMAC criterion is respected.
\item Explain two ways in which a heavy debt service can reduce the effectiveness of fiscal policy in Cameroon.
\end{enumerate}
\end{exercice}

\begin{corrige}[Exercise 1]
\begin{enumerate}
\item The budget deficit is a \emph{flow}: the excess of expenditure over revenue in one year (e.g.\ a deficit of 500 billion FCFA in the 2024 budget). The public debt is a \emph{stock}: the accumulation of all past borrowing not yet repaid (e.g.\ 12\,000 billion FCFA outstanding at end-2024).
\item Debt ratio $= \dfrac{9\,000}{15\,000}\times 100 = 60\%$. Since $60\% < 70\%$, the CEMAC convergence criterion is respected.
\item (i) Debt service absorbs budgetary resources, crowding out investment and social spending, so the State cannot finance expansionary programmes. (ii) High debt reduces credibility with lenders, raising interest costs and limiting the State's ability to borrow counter-cyclically during shocks.
\end{enumerate}
\end{corrige}

\begin{aretenir}[Key points of the section]
\begin{itemize}
\item The \cle{public debt} is the stock of all outstanding State borrowing; \cle{debt service} is the annual interest plus principal repayment.
\item Debt is \textbf{internal} (domestic lenders) or \textbf{external} (multilateral, bilateral, commercial creditors), and short, medium or long term.
\item The key indicator is the \textbf{debt-to-GDP ratio}; the CEMAC ceiling is \textbf{70\% of GDP}.
\item Debt can finance productive investment, but excessive debt crowds out expenditure and burdens future generations; Cameroon obtained major relief under \textbf{HIPC in 2006}.
\item \cle{Fiscal policy} uses taxation and expenditure to steer the economy; in Cameroon it shows real achievements (revenue, infrastructure) but serious limits (narrow base, evasion, informality, oil dependence, debt burden).
\end{itemize}
\end{aretenir}

\newpage

\coursec{Integration activity}

\courssub{Aims of the activity}

This integration activity closes the chapter on \cle{Public Finance} by asking you to act, in groups, as the executive of a Cameroonian municipal council. By the end of the activity you should be able to:

\begin{itemize}
\item \textbf{Classify} the revenues of a public authority by their nature (direct taxes, indirect taxes, fees, grants, borrowing) and justify each classification;
\item \textbf{Compute} the \cle{Value Added Tax} (VAT) payable by a local trading company from purchase and sales invoices, applying the Cameroonian rate of $19.25\%$;
\item \textbf{Construct} a balanced mini-budget distinguishing \cle{recurrent expenditure} from \cle{investment (capital) expenditure}, in line with the logic of the State budget and the programmed budget;
\item \textbf{Evaluate} a borrowing decision by computing simple debt-service indicators and weighing the advantages and limits of the \cle{public debt};
\item \textbf{Defend} your financial plan orally before the class, and \textbf{link} your decisions to the \cle{effectiveness of fiscal policy} in Cameroon.
\end{itemize}

\begin{aretenir}[Skills mobilised]
Knowledge (types of taxes, budget structure, public debt), know-how (VAT calculation, budget arithmetic, ratios) and soft skills (teamwork, argumentation, public speaking). All the notions of Lessons 94 to 100 are called upon.
\end{aretenir}

\courssub{Situation}

You are the newly elected executive of the \textbf{Municipal Council of Nkongsamba-Nord}, a fictional municipality of the Littoral Region, on the road between Douala and Bafoussam. The municipality has a population of $45\,000$ inhabitants, of whom about $6\,000$ are registered economic operators (shopkeepers, transporters, artisans, market traders). The council must prepare its budget for the coming financial year and present it to the municipal council for adoption.

The finance committee has gathered the following data.

\textbf{A. Expected own revenues of the council}

\begin{center}
\begin{tabularx}{\linewidth}{|l|X|r|}
\hline
\rowcolor{popblueL} \textbf{Code} & \textbf{Revenue item} & \textbf{Amount (FCFA)} \\
\hline
R1 & Market stall fees and daily market tolls & $48\,000\,000$ \\
\hline
R2 & Business licences (\emph{patentes}) paid by local traders & $30\,000\,000$ \\
\hline
R3 & Property tax collected on built land in the town & $22\,000\,000$ \\
\hline
R4 & Parking and cattle-transit fees & $10\,000\,000$ \\
\hline
R5 & Annual grant transferred by the State (decentralisation funds) & $60\,000\,000$ \\
\hline
\end{tabularx}
\end{center}

\textbf{B. The company ``Nkongsamba Trading SARL''}

The largest VAT-registered firm in the municipality is Nkongsamba Trading SARL, which buys foodstuffs from farmers and wholesalers and resells them to retailers. For the quarter under review, its accountant presents the following summary of invoices (all amounts are \emph{exclusive} of VAT; the VAT rate in Cameroon is $19.25\%$):

\begin{center}
\begin{tabularx}{\linewidth}{|X|r|}
\hline
\rowcolor{popblueL} \textbf{Operation} & \textbf{Amount excl. VAT (FCFA)} \\
\hline
Purchases of goods (VAT invoiced by suppliers) & $40\,000\,000$ \\
\hline
Purchase of a delivery tricycle (VAT invoiced) & $4\,000\,000$ \\
\hline
Sales of goods (VAT charged to customers) & $70\,000\,000$ \\
\hline
\end{tabularx}
\end{center}

The council wants to understand how much VAT this company will actually pay over to the tax administration, because the municipality benefits indirectly from a buoyant, tax-compliant local economy.

\textbf{C. Expenditure needs identified by the council services}

\begin{center}
\begin{tabularx}{\linewidth}{|l|X|r|}
\hline
\rowcolor{popblueL} \textbf{Code} & \textbf{Expenditure item} & \textbf{Amount (FCFA)} \\
\hline
E1 & Salaries of council staff and allowances of councillors & $55\,000\,000$ \\
\hline
E2 & Fuel, stationery, electricity and running costs & $20\,000\,000$ \\
\hline
E3 & Maintenance of existing roads and markets & $15\,000\,000$ \\
\hline
E4 & Construction of a new market hall (investment) & $80\,000\,000$ \\
\hline
E5 & Purchase of a refuse-collection truck (investment) & $25\,000\,000$ \\
\hline
E6 & Social actions (support to schools, health post) & $10\,000\,000$ \\
\hline
\end{tabularx}
\end{center}

\textbf{D. The borrowing option}

Total own revenues plus the State grant do not cover all the expenditure needs. A commercial bank proposes to the council a loan of $60\,000\,000$ FCFA, repayable over $5$ years in equal annual instalments of principal, at an interest rate of $10\%$ per year calculated on the outstanding balance. Some councillors are enthusiastic: the new market hall would raise future market-fee revenues. Others fear a heavy \cle{debt burden} that would strangle future budgets.

As the council executive, you must prepare a coherent, balanced and defensible financial plan.

\courssub{Tasks}

\begin{enumerate}
\item \textbf{Classify the revenues.} For each revenue item R1 to R5, state whether it is a direct tax, an indirect tax, a fee (\emph{redevance}), or a transfer/grant. Justify each answer in one sentence, using the characteristics of a tax studied in Lesson 95.
\item \textbf{Compute the VAT payable by Nkongsamba Trading SARL.} Calculate (a) the VAT collected on sales, (b) the deductible VAT on purchases of goods and on the tricycle, (c) the net VAT payable to the tax administration for the quarter. Explain in two sentences why VAT is called an \emph{indirect} tax and who really bears it.
\item \textbf{Draw up the mini-budget.} (a) Compute total revenue (R1 to R5) and total expenditure (E1 to E6). (b) Split expenditure into recurrent and investment items. (c) Show that the budget is in deficit before financing, and state the amount to be financed.
\item \textbf{Analyse the borrowing decision.} (a) Compute the first-year debt service (instalment of principal $+$ interest) on the proposed loan of $60\,000\,000$ FCFA. (b) Express this first-year debt service as a percentage of the council's own revenues (R1 to R4). (c) Give \textbf{two arguments for} and \textbf{two arguments against} borrowing, referring to the golden rule of public finance (borrowing should finance investment, not recurrent spending) and to the risks of over-indebtedness.
\item \textbf{Propose and defend a final balanced budget.} Present your adopted budget (revenues, expenditures, loan, residual balance) in a clear table, and prepare a five-minute defence before the class acting as the municipal council. Be ready to answer: \emph{``Why not simply raise market fees?''} and \emph{``What happens if the State grant is paid late?''}
\item \textbf{Plenary discussion.} As a class, discuss: \emph{``What does this exercise teach us about the effectiveness of fiscal policy in Cameroon?''} Consider the narrow tax base, the weight of indirect taxation, delays in transfers, and the role of borrowing in financing development.
\end{enumerate}

\begin{methode}[How to organise the group work]
\begin{enumerate}
\item Form groups of 4 to 6; appoint a mayor, a finance officer, a secretary and a spokesperson.
\item Spend about 20 minutes on Tasks 1--3 (calculations), 15 minutes on Task 4, 15 minutes preparing the defence (Task 5).
\item Each group presents; the class votes on the most convincing budget; the teacher concludes with the plenary discussion (Task 6).
\end{enumerate}
\end{methode}

\courssub{Model answer}

\textbf{Task 1 — Classification of revenues.}

\begin{center}
\begin{tabularx}{\linewidth}{|l|l|X|}
\hline
\rowcolor{popblueL} \textbf{Item} & \textbf{Nature} & \textbf{Justification} \\
\hline
R1 Market fees & Fee (\emph{redevance}) & Paid in direct exchange for a service (use of a stall); the amount is linked to the service rendered. \\
\hline
R2 Business licences & Direct tax & Compulsory levy on the exercise of an activity, borne directly by the trader, with no direct counterpart. \\
\hline
R3 Property tax & Direct tax & Compulsory levy on ownership of built property, paid directly by the taxpayer to the authority. \\
\hline
R4 Parking/transit fees & Fee & Payment for a specific service (parking space, transit facilities). \\
\hline
R5 State grant & Transfer & Not a tax at all: an unconditional transfer from the central budget to the council. \\
\hline
\end{tabularx}
\end{center}

\textbf{Task 2 — VAT of Nkongsamba Trading SARL.}

\begin{align*}
\text{VAT collected on sales} &= 70\,000\,000 \times 19.25\% = 13\,475\,000\ \text{FCFA}\\
\text{Deductible VAT on goods} &= 40\,000\,000 \times 19.25\% = 7\,700\,000\ \text{FCFA}\\
\text{Deductible VAT on tricycle} &= 4\,000\,000 \times 19.25\% = 770\,000\ \text{FCFA}\\
\text{Net VAT payable} &= 13\,475\,000 - (7\,700\,000 + 770\,000) = 5\,005\,000\ \text{FCFA}
\end{align*}

VAT is an \cle{indirect tax} because it is collected by the seller but economically borne by the \textbf{final consumer}, through the price. The company is only an intermediary: it deducts the VAT it paid on its inputs and remits only the difference.

\textbf{Task 3 — The mini-budget.}

Total revenue: $48 + 30 + 22 + 10 + 60 = 170$ million FCFA.

\begin{center}
\begin{tabularx}{\linewidth}{|l|X|r|}
\hline
\rowcolor{popblueL} \textbf{Category} & \textbf{Items} & \textbf{Amount (FCFA)} \\
\hline
Recurrent expenditure & E1, E2, E3, E6 & $100\,000\,000$ \\
\hline
Investment expenditure & E4, E5 & $105\,000\,000$ \\
\hline
\textbf{Total expenditure} & E1--E6 & $205\,000\,000$ \\
\hline
\end{tabularx}
\end{center}

Balance before financing: $170 - 205 = -35$ million FCFA. The budget shows a \cle{deficit} of $35\,000\,000$ FCFA to be financed.

\textbf{Task 4 — The borrowing decision.}

Annual principal instalment: $60\,000\,000 / 5 = 12\,000\,000$ FCFA. First-year interest: $60\,000\,000 \times 10\% = 6\,000\,000$ FCFA. First-year debt service: $12 + 6 = 18$ million FCFA.

Own revenues (R1--R4) $= 110$ million FCFA, so the debt-service ratio is:
\[ \frac{18\,000\,000}{110\,000\,000} \times 100 \approx 16.4\% \]

\emph{Arguments for borrowing:} (i) the loan finances an \textbf{investment} (the market hall), respecting the golden rule, and the hall should raise future market-fee revenue (self-financing logic); (ii) spreading the cost over 5 years matches the long life of the infrastructure (intergenerational equity). \emph{Arguments against:} (i) debt service of about $16.4\%$ of own revenues is significant and reduces room for manoeuvre if revenues disappoint; (ii) if the State grant is delayed or market fees fall, the council may be forced to cut recurrent services or borrow again — the spiral of over-indebtedness observed in many developing economies.

\textbf{Task 5 — Final balanced budget (one possible proposal).}

The executive decides to take the loan of $60$ million, build the market hall ($80$ million) and the truck ($25$ million), and keep a safety margin:

\begin{center}
\begin{tabularx}{\linewidth}{|X|r|X|r|}
\hline
\rowcolor{popblueL} \textbf{Revenue} & \textbf{FCFA} & \textbf{Expenditure} & \textbf{FCFA} \\
\hline
Own revenues (R1--R4) & $110\,000\,000$ & Recurrent (E1--E3, E6) & $100\,000\,000$ \\
\hline
State grant (R5) & $60\,000\,000$ & Investment (E4, E5) & $105\,000\,000$ \\
\hline
Bank loan & $60\,000\,000$ & Debt service (year 1) & $18\,000\,000$ \\
\hline
 & & Reserve / carry-over & $7\,000\,000$ \\
\hline
\textbf{Total} & $230\,000\,000$ & \textbf{Total} & $230\,000\,000$ \\
\hline
\end{tabularx}
\end{center}

The budget is balanced at $230$ million FCFA. In defence: raising market fees beyond the current level would risk pushing traders into informality (a narrow tax base cannot absorb sharp increases — a key limit of fiscal policy in Cameroon); and the reserve of $7$ million plus the possibility of phasing the truck purchase protect the council against a late State grant.

\textbf{Task 6 — Plenary pointers.} The exercise shows that fiscal policy in Cameroon operates under real constraints: a narrow formal tax base, heavy reliance on indirect taxes such as VAT (easy to collect but potentially regressive), dependence of councils on central-government transfers, and the temptation of borrowing. Effectiveness therefore depends on widening the tax base, collecting existing taxes honestly, programming expenditure over several years, and borrowing only for productive investment.

\begin{attention}[Common mistakes to avoid]
Computing VAT on VAT-inclusive amounts without first extracting the tax; forgetting the deductible VAT on the tricycle (investment goods also carry deductible VAT); classifying the State grant as a tax; treating the loan as ordinary revenue without recording the debt service as expenditure; and presenting an unbalanced budget.
\end{attention}

\newpage

\coursec{Methods and worked examples}

\coursec{Meaning and Scope of Public Finance}

\begin{definition}[Public Finance]
Public Finance is the branch of economics that studies the \textbf{income} (revenue), \textbf{expenditure}, and \textbf{debt} operations of the public authorities (Central Government, Local Councils, Public Enterprises) and their impact on the allocation of resources, distribution of income, macroeconomic stability, and economic growth.
\end{definition}

\begin{propriete}[Public Finance vs. Private Finance]
The fundamental differences arise from the \textbf{coercive power} of the State and its \textbf{social welfare objective}.
\begin{center}
\begin{tabularx}{\linewidth}{|l|X|X|}
\hline
\rowcolor{popblueL} \textbf{Criterion} & \textbf{Public Finance (State)} & \textbf{Private Finance (Household/Firm)} \\
\hline
\textbf{Resource Mobilisation} & \textbf{Coercive/Compulsory:} Legal authority to impose taxes (e.g., IRPP, TVA in Cameroon). Non-payment $\rightarrow$ legal sanctions. & \textbf{Voluntary/Contractual:} Income from wages, profits, loans. No power to force payment. \\
\hline
\textbf{Objective} & \textbf{Social Welfare Maximisation:} Provision of public goods, redistribution, stability. Deficits acceptable for public interest. & \textbf{Private Utility/Profit Maximisation:} Maximise individual satisfaction or firm profit. Persistent deficit $\rightarrow$ bankruptcy. \\
\hline
\textbf{Budgeting Process} & \textbf{Expenditure determines Revenue:} State decides expenditure priorities first, then seeks revenue (tax rates, borrowing). Deficit financing is a standard tool. & \textbf{Revenue determines Expenditure:} Estimate income first, then plan spending within constraint. \\
\hline
\textbf{Time Horizon} & \textbf{Long/Perpetual:} Going concern; plans for generations (infrastructure, debt sustainability). & \textbf{Finite/Limited:} Lifespan (individual) or business cycle/liquidity (firm). \\
\hline
\textbf{Ability to Borrow/Print} & \textbf{High:} Can borrow domestically/externally (BEAC, IMF, Eurobonds). In CEMAC, cannot unilaterally print money but influences regional policy. & \textbf{Constrained:} Borrowing limited by creditworthiness/collateral. Cannot issue legal tender. \\
\hline
\end{tabularx}
\end{center}
\end{propriete}

\begin{aretenir}[Four Functions of the State Budget (Musgrave)]
\begin{enumerate}
    \item \textbf{Allocation Function:} Provision of \cle{Public Goods} (non-rival, non-excludable: national defence, street lighting in Yaoundé) and \cle{Merit Goods} (health, education) where market fails due to free-rider problem.
    \item \textbf{Distribution Function:} Reducing inequality via progressive taxes (IRPP brackets) and transfers (subsidies on fertilizer for farmers in the North West, \textit{Filets Sociaux} cash transfers in Far North).
    \item \textbf{Stabilisation Function:} Counter-cyclical policy. Expansionary budget (deficit) during recession (post-COVID recovery); Contractionary (surplus) during inflation.
    \item \textbf{Growth Function:} Investing in infrastructure (Kribi Deep Sea Port, Lom Pangar Dam) and human capital to shift PPF/LRAS rightward.
\end{enumerate}
\end{aretenir}

\begin{attention}[GCE Trap: "Why does the State intervene?"]
Always cite \textbf{Market Failure} (Public goods, Externalities, Monopoly power, Inequity) and \textbf{Macroeconomic Instability}. Do not just say "to build roads".
\end{attention}

\begin{exoresolu}[Public Finance vs. Private Finance and State Functions]
\textbf{Statement:}
\begin{enumerate}
    \item Distinguish between Public Finance and Private Finance using three criteria. (6 marks)
    \item Explain the \emph{Allocation} and \emph{Distribution} functions of the State budget, giving a concrete Cameroonian example for each. (8 marks)
    \item Why is the provision of national defence a classic example of the allocation function? (3 marks)
\end{enumerate}
\tcblower
\textbf{Detailed Solution:}

\textbf{1. Distinction between Public Finance and Private Finance}

\begin{center}
\begin{tabularx}{\linewidth}{|l|X|X|}
\hline
\rowcolor{popblueL} \textbf{Criterion} & \textbf{Public Finance (State)} & \textbf{Private Finance (Household/Firm)} \\
\hline
\textbf{Resource Mobilisation} & \textbf{Coercive/Compulsory:} The State uses legal authority to impose taxes (e.g., \cle{Impôt sur le Revenu des Personnes Physiques - IRPP}, \cle{TVA} in Cameroon). Non-payment leads to legal sanctions. & \textbf{Voluntary/Contractual:} Income comes from wages (labour contract), profits (voluntary exchange), or loans (contractual agreement). No power to force payment. \\
\hline
\textbf{Objective} & \textbf{Social Welfare Maximisation:} Aim is to maximise community welfare (provision of public goods, redistribution, stability). Budget deficits are acceptable if they serve public interest. & \textbf{Private Utility/Profit Maximisation:} Aim is to maximise individual satisfaction or firm profit. A persistent deficit leads to bankruptcy. \\
\hline
\textbf{Budgeting Process} & \textbf{Expenditure determines Revenue:} The State first decides on expenditure priorities (Budget preparation) then seeks revenue (tax rates, borrowing) to finance them. Deficit financing is a standard tool. & \textbf{Revenue determines Expenditure:} Agents estimate income first, then plan spending within that constraint. "Cut your coat according to your size." \\
\hline
\textbf{Time Horizon} & \textbf{Long/Perpetual:} The State is a going concern; it plans for generations (infrastructure, debt sustainability over decades). & \textbf{Finite/Limited:} Limited by lifespan (individual) or business cycle/liquidity (firm). \\
\hline
\textbf{Ability to Borrow/Print} & \textbf{High:} Can borrow domestically/externally (BEAC, IMF, Eurobonds). In a monetary union (CEMAC), it cannot unilaterally print money but influences regional policy. & \textbf{Constrained:} Borrowing limited by creditworthiness/collateral. Cannot issue legal tender. \\
\hline
\end{tabularx}
\end{center}

\textbf{2. Allocation and Distribution Functions (Cameroon Context)}

\begin{itemize}
    \item \textbf{Allocation Function:} The State provides \cle{Public Goods} (non-rival, non-excludable) and \cle{Merit Goods} which the market under-provides due to the \textbf{Free Rider Problem}.
    \begin{itemize}
        \item \textit{Cameroon Example:} Construction and maintenance of the \textbf{Yaoundé-Nsimalen Highway} or the \textbf{Kribi Deep Sea Port}. Private firms would not build these because they cannot easily charge each user (non-excludability) and one user's use doesn't reduce availability (non-rivalry). The State finances them via tax revenue for the collective benefit.
    \end{itemize}
    \item \textbf{Distribution Function:} The State modifies the \textbf{market distribution of income} (which may be highly unequal) to achieve a socially desired \textbf{equity} (vertical and horizontal).
    \begin{itemize}
        \item \textit{Cameroon Example:} The \textbf{Progressive Income Tax (IRPP)} schedule (rates from 10\% to 35\%+ surcharges) takes a larger proportion from high earners (e.g., senior executives in Douala). Revenue funds \textbf{social safety nets} like the \textbf{Filets Sociaux} cash transfers to vulnerable households in the Far North or subsidies on \textbf{fertilizer (engrais)} for smallholder cocoa/coffee farmers in the South West, effectively transferring purchasing power from rich to poor.
    \end{itemize}
\end{itemize}

\textbf{3. National Defence and the Allocation Function}
National defence is the \textbf{quintessential Pure Public Good}.
\begin{itemize}
    \item \textbf{Non-excludability:} It is impossible (or prohibitively expensive) to exclude a citizen within Cameroon's borders from the protection provided by the army (BIR, Police, Gendarmerie) once it is provided. You cannot "opt out" of national defence.
    \item \textbf{Non-rivalry:} One citizen's enjoyment of security does not diminish the security available to another.
    \item \textbf{Market Failure:} Because of non-excludability, a private firm cannot charge a price $\rightarrow$ no profit incentive $\rightarrow$ \textbf{Market Failure (Missing Market)}.
    \item \textbf{State Role:} Only the State, using its coercive tax power, can finance defence compulsorily, ensuring the socially optimal quantity is provided. This is the core of the \textbf{Allocation Function}: correcting market failure for public goods.
\end{itemize}
\end{exoresolu}

\coursec{Taxation: Meaning, Characteristics and Principles}

\begin{definition}[Tax]
A \cle{Tax} is a compulsory, unrequited payment to the general government, levied on income, wealth, consumption, or transactions, to finance public expenditure. Key characteristics: \textbf{Compulsory} (legal obligation), \textbf{No direct quid pro quo} (no specific good/service received in return), \textbf{Public purpose} (finances collective needs).
\end{definition}

\begin{propriete}[Canons of Taxation (Adam Smith + Modern)]
A good tax system should satisfy:
\begin{enumerate}
    \item \textbf{Equity (Fairness):} Horizontal (equals treated equally) and Vertical (unequals treated unequally $\rightarrow$ Ability-to-Pay principle).
    \item \textbf{Certainty:} Tax liability (amount, time, manner) must be clear and plain to the taxpayer.
    \item \textbf{Convenience:} Levied at the time/manner most convenient for payer (e.g., PAYE at source, VAT at checkout).
    \item \textbf{Economy:} Cost of collection $\ll$ Revenue collected. Low administrative burden.
    \item \textbf{Productivity (Modern):} Yields sufficient revenue for government needs.
    \item \textbf{Elasticity (Modern):} Revenue grows automatically with GDP (buoyancy).
    \item \textbf{Simplicity (Modern):} Easy to understand and comply with.
    \item \textbf{Diversity (Modern):} Mix of taxes to broaden base and reduce distortion.
\end{enumerate}
\end{propriete}

\coursec{Types of Taxes}

\begin{methode}[Classifying Taxes]
\textbf{Objective:} Classify any tax along three dimensions: \textbf{Incidence} (Direct/Indirect), \textbf{Rate Structure} (Proportional/Progressive/Regressive), \textbf{Base Measurement} (Ad Valorem/Specific).
\begin{enumerate}
    \item \textbf{Direct vs. Indirect (Incidence/Shifting):}
    \begin{itemize}
        \item \textbf{Direct Tax:} Impact = Incidence. Levied on person/income/wealth. Cannot be shifted. Examples: \cle{IRPP} (Personal Income Tax), \cle{IS} (Company Tax), Property Tax.
        \item \textbf{Indirect Tax:} Impact $\neq$ Incidence. Levied on goods/services/transactions. Can be shifted to consumer via higher prices. Examples: \cle{VAT} (19.25\%), \cle{Excise Duties} (beer, cigarettes, fuel), \cle{Customs Duties}.
    \end{itemize}
    \item \textbf{Rate Structure (Vertical Equity):}
    \begin{itemize}
        \item \textbf{Proportional (Flat):} Constant Average Tax Rate (ATR) = Marginal Tax Rate (MTR). $T = t \times Y$. (e.g., Company Tax 30\% in Cameroon).
        \item \textbf{Progressive:} ATR rises with Income ($Y$). MTR $>$ ATR. Reduces inequality. (e.g., Cameroon IRPP brackets).
        \item \textbf{Regressive:} ATR falls as $Y$ rises. MTR $<$ ATR. Hurts poor more. (e.g., Specific tax on salt/fuel if consumption doesn't rise with income; Poll tax).
    \end{itemize}
    \item \textbf{Base Measurement:}
    \begin{itemize}
        \item \textbf{Ad Valorem:} Tax $\%$ of Value ($t \times P \times Q$). Revenue rises with inflation. (e.g., VAT 19.25\%, Import Duty).
        \item \textbf{Specific (Per Unit):} Fixed amount per unit ($T \times Q$). Revenue stable in real terms only if indexed. (e.g., Excise duty: 50 FCFA/litre of beer, 200 FCFA/pack of cigarettes).
    \end{itemize}
\end{enumerate}
\end{methode}

\begin{exoresolu}[Tax Classification, Progressivity Calculation and Canon Evaluation]
\textbf{Statement:}
In Cameroon, the Income Tax on Individuals (IRPP) has the following annual brackets (simplified):
\begin{itemize}
    \item 0 -- 2,000,000 FCFA: Exempt (0\%)
    \item 2,000,001 -- 3,000,000 FCFA: 10\%
    \item 3,000,001 -- 5,000,000 FCFA: 15\%
    \item Above 5,000,000 FCFA: 25\%
\end{itemize}
A specific Excise Duty of 150 FCFA per litre is imposed on locally produced beer. VAT is 19.25\% (Standard Rate).
\begin{enumerate}
    \item Classify IRPP, Beer Excise Duty, and VAT as Direct/Indirect, and by Rate Structure (Proportional/Progressive/Regressive, Ad Valorem/Specific). (6 marks)
    \item Calculate the Total Tax Payable, Average Tax Rate (ATR), and Marginal Tax Rate (MTR) for a taxpayer earning 4,500,000 FCFA/year. Show workings. (6 marks)
    \item Evaluate the Beer Excise Duty against the Canons of \textbf{Equity}, \textbf{Elasticity}, and \textbf{Economy}. (5 marks)
\end{enumerate}
\tcblower
\textbf{Detailed Solution:}

\textbf{1. Classification Table}

\begin{center}
\begin{tabular}{|l|c|c|c|}
\hline
\rowcolor{popblueL} \textbf{Tax} & \textbf{Direct / Indirect} & \textbf{Rate Structure (Base)} & \textbf{Progressivity Type} \\
\hline
\textbf{IRPP (Income Tax)} & \textbf{Direct} (Levied on person, cannot shift) & \textbf{Ad Valorem} (\% of Income) & \textbf{Progressive} (ATR rises with Income) \\
\hline
\textbf{Beer Excise Duty} & \textbf{Indirect} (Levied on producer, shifted to consumer via price) & \textbf{Specific} (Fixed FCFA/litre) & \textbf{Regressive} (ATR falls as income rises; poor spend higher \% income on beer) \\
\hline
\textbf{VAT (19.25\%)} & \textbf{Indirect} (Levied on transaction, shifted to final consumer) & \textbf{Ad Valorem} (\% of Value Added/Price) & \textbf{Proportional} (Constant rate) but \textbf{Regressive in effect} (Poor spend higher \% of income on consumption) \\
\hline
\end{tabular}
\end{center}

\textbf{2. IRPP Calculation for Income = 4,500,000 FCFA}
\textit{Method: Apply marginal rates to income \textbf{within each bracket}.}

\begin{align*}
\text{Bracket 1: } & 0 \text{ to } 2,000,000 \quad \rightarrow \text{Tax} = 0 \text{ FCFA} \\
\text{Bracket 2: } & 2,000,001 \text{ to } 3,000,000 \quad (\text{Width} = 1,000,000) \quad \rightarrow \text{Tax} = 1,000,000 \times 10\% = 100,000 \text{ FCFA} \\
\text{Bracket 3: } & 3,000,001 \text{ to } 4,500,000 \quad (\text{Width} = 1,500,000) \quad \rightarrow \text{Tax} = 1,500,000 \times 15\% = 225,000 \text{ FCFA} \\
\text{Bracket 4: } & \text{Not reached.} \\
\hline
\textbf{Total Tax Payable (T)} & = 0 + 100,000 + 225,000 = \textbf{325,000 FCFA} \\
\textbf{Average Tax Rate (ATR)} & = \frac{T}{Y} \times 100 = \frac{325,000}{4,500,000} \times 100 = \textbf{7.22\%} \\
\textbf{Marginal Tax Rate (MTR)} & = \text{Rate on the \textbf{last franc earned}} = \text{Rate of Bracket 3} = \textbf{15\%}
\end{align*}
\textbf{Check:} MTR (15\%) $>$ ATR (7.22\%) $\rightarrow$ Confirms \textbf{Progressivity}.

\textbf{3. Evaluation of Specific Beer Excise Duty (150 FCFA/Litre)}

\begin{itemize}
    \item \textbf{Equity (Vertical): FAILS / WEAK.} It is a \textbf{Regressive} tax. A low-income earner spending 50,000 FCFA/month on beer pays the same absolute tax (150 FCFA/L) as a high-income earner. As a \% of income, the burden falls heavily on the poor. \textit{Horizontal Equity:} Treats all beer drinkers equally regardless of ability to pay.
    \item \textbf{Elasticity: FAILS / LOW.} Revenue $R = 150 \times Q$. Revenue grows \textbf{only if volume $Q$ grows}. It does \textbf{not} automatically adjust for \textbf{Inflation} or \textbf{Real Income Growth}. In Cameroon, with inflation $\approx$ 4-5\%, the real value of 150 FCFA erodes yearly. Government must legislate frequent upward revisions (politically difficult). Contrast with Ad Valorem VAT (19.25\%) which automatically captures price inflation.
    \item \textbf{Economy (Administrative Efficiency): PASSES / STRONG.} Very \textbf{cheap and easy to administer}. Tax point is clear: factory gate or customs entry. Easy to measure volume (litres) vs. valuing transactions (Ad Valorem). Low evasion (hard to hide physical production volume compared to hiding value). Low collection cost / revenue ratio.
\end{itemize}
\textbf{Conclusion:} Good for \textbf{sumptuary/discouragement} purpose (sin tax) and admin simplicity; bad for revenue buoyancy and vertical equity.
\end{exoresolu}

\coursec{VAT in Cameroon and the Structure of Taxation (Practical Work)}

\begin{methode}[Calculating VAT in Cameroon (Invoice/Credit Method)]
\textbf{Objective:} Compute Output VAT, Input VAT (Deductible), and Net VAT Payable/Refundable for a VAT-registered business in Cameroon (Regime Réel Normal). Rate: \textbf{19.25\%} (Standard), \textbf{0\%} (Exports), Exempt (no credit).
\begin{enumerate}
    \item \textbf{Identify Operations:} 
    \begin{itemize}
        \item \textbf{Taxable (Standard 19.25\%):} Most local sales of goods/services.
        \item \textbf{Zero-Rated (0\%):} Exports (goods leaving CEMAC), International transport. \textit{Key: Input VAT is DEDUCTIBLE.}
        \item \textbf{Exempt:} Medical, Education, Financial services, Unprocessed agricultural products (raw cocoa beans sold by farmer). \textit{Key: Input VAT is \textbf{NOT} deductible (cost absorbed).}
    \end{itemize}
    \item \textbf{Calculate Output VAT (Collected):} $\text{Output VAT} = \sum (\text{Sales Price}_{HT} \times 19.25\%)$ for taxable sales + $(\text{Export Sales} \times 0\%)$.
    \item \textbf{Calculate Deductible Input VAT (Paid on Inputs):} Sum VAT on \textbf{purchases of goods/services used for taxable/zero-rated operations}. \textbf{Exclude:} VAT on exempt activity inputs, VAT on non-deductible expenses (luxury cars > threshold, fuel for private use, gifts, reception costs).
    \item \textbf{Net VAT Position:} $\text{VAT Payable} = \text{Output VAT} - \text{Deductible Input VAT}$.
    \begin{itemize}
        \item If $> 0$: Pay to Tax Centre by 15th of following month (Déclaration CA3).
        \item If $< 0$: \textbf{Credit Vatable} (Reportable to next month). Refund only for Exporters (Credit $>$ 500,000 FCFA persistent) or Cessation of activity.
    \end{itemize}
    \item \textbf{Invoice Requirements:} Must show: Seller/Buyer NIU, "TVA 19.25\%", HT Amount, VAT Amount, TTC Amount. No VAT on Exempt sales.
\end{enumerate}
\end{methode}

\begin{exoresolu}[VAT Calculation for a Cameroonian Agro-Processing Firm (Practical Work 21 Style)]
\textbf{Statement:}
\textbf{SOCATRANS SARL}, a VAT-registered company in Douala (Regime Réel Normal), processes cocoa into cocoa butter and powder. In \textbf{January 2025}, it recorded the following operations (All amounts HT - Hors Taxes):
\begin{itemize}
    \item \textbf{Sales:}
    \begin{itemize}
        \item Local sale of Cocoa Butter to CHOCOCAM: 50,000,000 FCFA.
        \item Export of Cocoa Powder to Europe: 30,000,000 FCFA.
        \item Sale of Cocoa Husks (waste) to local farmers (Exempt agricultural input): 5,000,000 FCFA.
    \end{itemize}
    \item \textbf{Purchases/Expenses (Invoices received with VAT 19.25\%):}
    \begin{itemize}
        \item Raw Cocoa Beans from Cooperatives (Exempt supply): 40,000,000 FCFA (No VAT charged).
        \item Packaging Materials (Cartons, bags): 10,000,000 FCFA.
        \item Electricity \& Water (Factory): 5,000,000 FCFA.
        \item Fuel for Factory Generator: 3,000,000 FCFA.
        \item Fuel for Director's Private Car: 500,000 FCFA.
        \item Consultancy Fees (Audit): 2,000,000 FCFA.
        \item Purchase of a Luxury SUV (Price 60,000,000 FCFA) - \textit{Note: VAT on passenger vehicles non-deductible}.
    \end{itemize}
\end{itemize}
\begin{enumerate}
    \item Determine the VAT treatment (Rate, Deductibility) for each sale and purchase line. (5 marks)
    \item Calculate the \textbf{Output VAT}, \textbf{Deductible Input VAT}, and \textbf{Net VAT Payable (or Credit)} for January 2025. (10 marks)
    \item Explain why VAT on Raw Cocoa Beans is not deductible, but VAT on Packaging is. (2 marks)
\end{enumerate}
\tcblower
\textbf{Detailed Solution:}

\textbf{1. VAT Treatment Analysis}

\begin{center}
\begin{tabularx}{\linewidth}{|l|c|c|X|}
\hline
\rowcolor{popblueL} \textbf{Operation} & \textbf{Regime} & \textbf{Rate} & \textbf{Input VAT Deductible?} \\
\hline
\textbf{SALES (Output)} & & & \\
Local Sale (Butter) & Taxable & 19.25\% & N/A (Output) \\
Export (Powder) & Zero-Rated & 0\% & N/A (Output) \\
Cocoa Husks (Farmers) & Exempt & Exempt & N/A (Output) \\
\hline
\textbf{PURCHASES (Input)} & & & \\
Raw Cocoa Beans & Exempt (Supplier) & -- & \textbf{NO} (Upstream exempt $\rightarrow$ no VAT paid) \\
Packaging Materials & Taxable & 19.25\% & \textbf{YES} (Used for taxable/zero-rated output) \\
Electricity/Water (Factory) & Taxable & 19.25\% & \textbf{YES} (Industrial use) \\
Fuel (Factory Generator) & Taxable & 19.25\% & \textbf{YES} (Direct production input) \\
Fuel (Private Car) & Taxable & 19.25\% & \textbf{NO} (Private use $\rightarrow$ Art. 147 CGI non-deductible) \\
Consultancy Fees & Taxable & 19.25\% & \textbf{YES} (General overhead for taxable activity) \\
Luxury SUV & Taxable & 19.25\% & \textbf{NO} (Passenger vehicle $\rightarrow$ Art. 147 CGI non-deductible) \\
\hline
\end{tabularx}
\end{center}

\textbf{2. Calculations (Rate = 19.25\% = 0.1925)}

\textbf{A. Output VAT (Collected)}
\begin{align*}
\text{Local Sale} &: 50,000,000 \times 0.1925 = 9,625,000 \text{ FCFA} \\
\text{Export Sale} &: 30,000,000 \times 0.00 = 0 \text{ FCFA} \\
\text{Exempt Sale (Husks)} &: 0 \text{ FCFA} \\
\textbf{Total Output VAT} &: \textbf{9,625,000 FCFA}
\end{align*}

\textbf{B. Deductible Input VAT (Paid on Eligible Purchases)}
\begin{align*}
\text{Packaging} &: 10,000,000 \times 0.1925 = 1,925,000 \\
\text{Electricity/Water} &: 5,000,000 \times 0.1925 = 962,500 \\
\text{Fuel (Factory)} &: 3,000,000 \times 0.1925 = 577,500 \\
\text{Consultancy} &: 2,000,000 \times 0.1925 = 385,000 \\
\text{Raw Beans} &: 0 \text{ (Exempt purchase)} \\
\text{Fuel (Private)} &: 0 \text{ (Non-deductible)} \\
\text{SUV} &: 0 \text{ (Non-deductible)} \\
\textbf{Total Deductible Input VAT} &: 1,925,000 + 962,500 + 577,500 + 385,000 = \textbf{3,850,000 FCFA}
\end{align*}

\textbf{C. Net VAT Position}
\[
\text{Net VAT Payable} = \text{Output VAT} - \text{Deductible Input VAT} = 9,625,000 - 3,850,000 = \textbf{5,775,000 FCFA}
\]
\textbf{Action:} SOCATRANS pays \textbf{5,775,000 FCFA} to the Douala Tax Centre by \textbf{15th February 2025} (Form CA3).

\textbf{3. Explanation: Raw Beans vs. Packaging}
\begin{itemize}
    \item \textbf{Raw Cocoa Beans:} Sold by cooperative (farmer) $\rightarrow$ \textbf{Exempt Operation} (Unprocessed agricultural product, Art. 136 CGI). The farmer charges \textbf{NO VAT}. Therefore, SOCATRANS paid \textbf{0 FCFA} Input VAT on this purchase. You cannot deduct what you didn't pay. The VAT "chain" is broken here; the value added by the farmer escapes tax.
    \item \textbf{Packaging Materials:} Sold by a VAT-registered manufacturer $\rightarrow$ \textbf{Taxable Operation (19.25\%)}. SOCATRANS paid VAT (1,925,000 FCFA) on the invoice. Since packaging is used for \textbf{taxable (local) and zero-rated (export) outputs}, the Input VAT is fully \textbf{deductible} (Right to deduction, Art. 145 CGI).
\end{itemize}
\end{exoresolu}

\begin{methode}[Analysing the Structure of Taxation and Tax Burden Indicators]
\textbf{Objective:} Interpret the Tax Structure (Direct vs Indirect share, Tax/GDP ratio, Tax Buoyancy/Elasticity) and calculate Tax Incidence/Shifting.
\begin{enumerate}
    \item \textbf{Structure Ratios (Cameroon Context $\approx$ 2023/24 Budget):}
    \begin{itemize}
        \item \textbf{Tax Pressure (Taux de Pression Fiscale):} $\frac{\text{Total Tax Revenue}}{\text{GDP}} \times 100$. Cameroon $\approx$ 11-13\% (Low vs WAEMU avg $\approx$ 15\%, OECD $\approx$ 34\%). Target: 15-17\%.
        \item \textbf{Direct / Indirect Split:} Cameroon historically \textbf{Indirect-heavy} ($\approx$ 60-65\% Indirect: VAT, Customs, Excise). Direct $\approx$ 35-40\% (IRPP, IS, Property). \textit{Implication:} Structure tends \textbf{Regressive}.
        \item \textbf{Customs vs. Domestic:} High reliance on Customs Duties (Droits de Douane) $\approx$ 25-30\% of tax rev. Vulnerable to import compression / AfCFTA liberalisation.
    \end{itemize}
    \item \textbf{Buoyancy vs. Elasticity:}
    \begin{itemize}
        \item \textbf{Buoyancy ($B$):} $\frac{\% \Delta \text{Actual Tax Revenue}}{\% \Delta \text{GDP}}$. Includes discretionary changes (rate hikes, new taxes).
        \item \textbf{Elasticity ($E$):} $\frac{\% \Delta \text{Revenue (Constant Policy)}}{\% \Delta \text{GDP}}$. Measures \textbf{automatic} response. $E > 1$ = Progressive structure (good for stabilisation). $E < 1$ = Inelastic (revenue lags growth).
    \end{itemize}
    \item \textbf{Incidence/Shifting (Partial Equilibrium):}
    \begin{itemize}
        \item \textbf{Statutory Incidence:} Who legally pays (Producer vs Consumer).
        \item \textbf{Economic Incidence:} Who \textbf{actually} bears burden (Real purchasing power loss).
        \item \textbf{Shifting:} Forward (Producer $\rightarrow$ Consumer via Price $\uparrow$), Backward (Producer $\rightarrow$ Factor owners via Wage/Rent $\downarrow$).
        \item \textbf{Determinant:} \textbf{Relative Price Elasticities} ($E_d$ vs $E_s$). 
        \begin{align*}
            \text{Consumer Burden Share} &= \frac{E_s}{|E_d| + E_s} \\
            \text{Producer Burden Share} &= \frac{|E_d|}{|E_d| + E_s}
        \end{align*}
        \textit{Rule:} The \textbf{more inelastic} side bears the burden.
    \end{itemize}
\end{enumerate}
\end{methode}

\begin{exoresolu}[Tax Structure Analysis and Incidence Calculation]
\textbf{Statement:}
Cameroon's 2023 Fiscal Data (Estimates):
\begin{itemize}
    \item Nominal GDP: 30,000 Billion FCFA.
    \item Total Tax Revenue: 3,600 Billion FCFA.
    \item Direct Taxes: 1,300 Billion FCFA.
    \item Indirect Taxes: 2,300 Billion FCFA (VAT: 1,400; Customs: 600; Excise: 300).
    \item Previous Year (2022): GDP 28,000 Billion; Tax Revenue 3,200 Billion. Discretionary measures (rate hikes) added 100 Billion in 2023.
\end{itemize}
A specific tax of 200 FCFA/unit is imposed on Cement. Market: $Q_d = 1,000,000 - 50P$; $Q_s = -200,000 + 100P$ (P in FCFA/bag, Q in bags).
\begin{enumerate}
    \item Calculate the \textbf{Tax Pressure}, \textbf{Direct/Indirect \% Share}, and \textbf{Customs Dependence Ratio}. Comment on equity implications. (6 marks)
    \item Calculate \textbf{Buoyancy} and \textbf{Elasticity} for 2023. Interpret. (5 marks)
    \item Determine the \textbf{Equilibrium Price/Quantity} before tax. Calculate the \textbf{New Equilibrium} after tax (levied on Producer). Determine \textbf{Consumer Burden}, \textbf{Producer Burden}, and \textbf{Government Revenue}. (8 marks)
\end{enumerate}
\tcblower
\textbf{Detailed Solution:}

\textbf{1. Structure Indicators}
\begin{align*}
\textbf{Tax Pressure} &= \frac{3,600}{30,000} \times 100 = \textbf{12.0\%} \\
\textbf{Direct Tax Share} &= \frac{1,300}{3,600} \times 100 = \textbf{36.1\%} \\
\textbf{Indirect Tax Share} &= \frac{2,300}{3,600} \times 100 = \textbf{63.9\%} \\
\textbf{Customs Dependence} &= \frac{600}{3,600} \times 100 = \textbf{16.7\%} \text{ (of Total Tax Rev)} \quad \text{or} \quad \frac{600}{2,300} \times 100 = \textbf{26.1\%} \text{ (of Indirect Rev)}
\end{align*}
\textbf{Comment (Equity):} A \textbf{12\% Tax/GDP} is low (sub-Saharan avg $\approx$ 15-17\%), limiting fiscal space for development (infrastructure, health). The \textbf{64\% Indirect dominance} implies a \textbf{Regressive Structure}: VAT/Excise/Customs fall on consumption. Poor households spend $\approx$ 80-90\% of income on consumption $\rightarrow$ high effective tax rate. Rich save/invest $\rightarrow$ lower effective rate. Reliance on Customs (16.7\% total) creates vulnerability to \textbf{AfCFTA tariff dismantling} (revenue loss risk).

\textbf{2. Buoyancy and Elasticity}
\begin{align*}
\% \Delta GDP &= \frac{30,000 - 28,000}{28,000} \times 100 = \frac{2,000}{28,000} \times 100 = \textbf{7.14\%} \\
\% \Delta Tax Rev (Actual) &= \frac{3,600 - 3,200}{3,200} \times 100 = \frac{400}{3,200} \times 100 = \textbf{12.5\%} \\
\textbf{Buoyancy (B)} &= \frac{12.5\%}{7.14\%} = \textbf{1.75} \\
\textbf{Elasticity (E)} &: \text{Revenue growth from automatic factors only.} \\
\text{Automatic Revenue Growth} &= \text{Actual Growth} - \text{Discretionary Measures} = 400 - 100 = 300 \text{ Billion}. \\
\% \Delta Tax Rev (Auto) &= \frac{300}{3,200} \times 100 = \textbf{9.375\%} \\
\textbf{Elasticity (E)} &= \frac{9.375\%}{7.14\%} = \textbf{1.31}
\end{align*}
\textbf{Interpretation:} Both $B > 1$ and $E > 1$. The tax system is \textbf{Elastic/Buoyant}. Revenue grows faster than GDP automatically (progressive income tax brackets, Ad Valorem VAT). This is \textbf{good for Stabilisation} (automatic stabilisers work) and reduces need for frequent discretionary hikes. However, $B > E$ shows discretionary measures (100 Billion) contributed significantly (25\% of revenue growth).

\textbf{3. Tax Incidence on Cement Market}
\textbf{Step 1: Initial Equilibrium (No Tax)}
$Q_d = Q_s \Rightarrow 1,000,000 - 50P = -200,000 + 100P$
$1,200,000 = 150P \Rightarrow P_0 = \textbf{8,000 FCFA/bag}$
$Q_0 = 1,000,000 - 50(8,000) = \textbf{600,000 bags}$

\textbf{Step 2: New Equilibrium (Tax $t = 200$ FCFA on Producer)}
Supply shifts UP by 200. New Supply Price $P_s = P_d - 200$ (Producer receives $P_d - 200$).
$Q_s' = -200,000 + 100(P_d - 200) = -200,000 + 100P_d - 20,000 = -220,000 + 100P_d$
Set $Q_d = Q_s'$:
$1,000,000 - 50P_d = -220,000 + 100P_d$
$1,220,000 = 150P_d \Rightarrow P_d^* = \textbf{8,133.33 FCFA}$ (Price Consumer Pays)
$P_s^* = P_d^* - 200 = \textbf{7,933.33 FCFA}$ (Price Producer Receives)
$Q^* = 1,000,000 - 50(8,133.33) = \textbf{593,333 bags}$

\textbf{Step 3: Burden Sharing (Using Elasticities at $P_0=8000, Q_0=600,000$)}
$E_d = \frac{dQ_d}{dP} \frac{P}{Q} = -50 \times \frac{8000}{600,000} = -\frac{400,000}{600,000} = -\textbf{0.667}$ (Inelastic Demand)
$E_s = \frac{dQ_s}{dP} \frac{P}{Q} = 100 \times \frac{8000}{600,000} = \frac{800,000}{600,000} = \textbf{1.333}$ (Elastic Supply)
\textit{Check Formula:} Consumer Share $= \frac{E_s}{|E_d| + E_s} = \frac{1.333}{0.667 + 1.333} = \frac{1.333}{2.0} = \textbf{0.6665 (66.7\%)}$
Producer Share $= \frac{|E_d|}{|E_d| + E_s} = \frac{0.667}{2.0} = \textbf{0.3335 (33.3\%)}$

\textbf{Step 4: Monetary Burden \& Revenue}
\begin{align*}
\text{Consumer Burden per unit} &= P_d^* - P_0 = 8,133.33 - 8,000 = \textbf{133.33 FCFA} \\
\text{Producer Burden per unit} &= P_0 - P_s^* = 8,000 - 7,933.33 = \textbf{66.67 FCFA} \\
\text{Check:} & 133.33 + 66.67 = 200 \text{ FCFA (Tax)}. \quad \text{Ratio: } 133.33/66.67 = 2.0 \approx 66.7/33.3. \\
\text{Govt Revenue} &= t \times Q^* = 200 \times 593,333 = \textbf{118,666,600 FCFA} \approx \textbf{118.7 Million FCFA}.
\end{align*}
\textbf{Conclusion:} Consumers bear 2/3 of the tax because Demand is Inelastic (Cement is necessity, few substitutes) and Supply is Elastic (producers can switch factors easily).
\end{exoresolu}

\coursec{Economic Consequences of Taxation}

\begin{methode}[Analysing Economic Consequences of Taxation (Micro \& Macro)]
\textbf{Objective:} Analyse effects on Allocation (Efficiency/Deadweight Loss), Distribution (Equity), Stabilisation (Multipliers), and Growth (Incentives).
\begin{enumerate}
    \item \textbf{Micro: Excess Burden (Deadweight Loss - DWL).} Tax creates wedge $P_d - P_s = t$. 
    \[
    \text{DWL} \approx \frac{1}{2} \times t \times \Delta Q = \frac{1}{2} \times t \times Q_0 \times (\eta_d + \eta_s) \times \frac{t}{P_0} \quad (\text{Small tax approx})
    \]
    \textit{Rule:} DWL $\uparrow$ with Elasticities ($\eta_d, \eta_s$) and $t^2$. Tax inelastic goods (fuel, medicine) for revenue (low DWL); Tax elastic goods (luxuries, sin goods) for discouragement (high DWL accepted).
    \item \textbf{Micro: Substitution \& Income Effects.}
    \begin{itemize}
        \item \textbf{Labour Supply:} Tax on wages $\rightarrow$ Substitution Effect (Leisure cheaper $\rightarrow$ Work Less) vs Income Effect (Poorer $\rightarrow$ Work More). Net effect ambiguous. Backward bending supply curve.
        \item \textbf{Savings/Investment:} Tax on interest/dividends $\rightarrow$ Reduces return $\rightarrow$ Substitution: Consume now. Income: Save more to target wealth. Distorts Intertemporal Choice.
    \end{itemize}
    \item \textbf{Macro: Fiscal Multipliers.}
    \begin{align*}
    \Delta Y &= \frac{1}{1 - c(1-t) + m} \times \Delta G \quad \text{(Spending Multiplier)} \\
    \Delta Y &= \frac{-c}{1 - c(1-t) + m} \times \Delta T \quad \text{(Tax Multiplier)} \quad (|\text{Tax Mult}| < |\text{Spend Mult}|)
    \end{align*}
    \textit{Why?} $\Delta G$ injects full amount into AD. $\Delta T$ reduces Disposable Income ($Y_d$), only $c \times \Delta T$ reduces Consumption ($C$); rest reduces Savings/Imports (Leakage).
    \item \textbf{Laffer Curve:} Tax Revenue $R = t \times Y(t)$. $Y$ falls as $t \uparrow$ (disincentives). $R$ rises then falls. Optimal $t^*$ maximises $R$. Cameroon Corporate Tax 30\% + 10\% surcharge = 33\% vs Regional avg $\approx$ 25-30\%. Risk of base erosion (profit shifting).
\end{enumerate}
\end{methode}

\begin{exoresolu}[Deadweight Loss, Labour Supply and Fiscal Multipliers]
\textbf{Statement:}
Consider a perfectly competitive market for Beer in Cameroon.
Demand: $P = 2000 - 0.5Q$ ; Supply: $P = 500 + 0.5Q$. (P in FCFA, Q in '000 Litres).
Government imposes a Specific Tax $t = 200$ FCFA/Litre on Producers.
Assume MPC $c = 0.8$, Marginal Tax Rate $t_x = 0.2$, MPM $m = 0.1$.
\begin{enumerate}
    \item Calculate Pre-tax Equilibrium ($P_0, Q_0$), Consumer Surplus (CS), Producer Surplus (PS). (3 marks)
    \item Calculate Post-tax Equilibrium ($P_d, P_s, Q_1$), Tax Revenue, New CS, New PS, and \textbf{Deadweight Loss (DWL)}. (7 marks)
    \item Calculate the \textbf{Government Spending Multiplier} and \textbf{Tax Multiplier}. If Govt uses the Beer Tax Revenue to increase $G$, what is the \textbf{Balanced Budget Multiplier} effect on $Y$? (5 marks)
    \item Discuss the impact of a high Corporate Tax rate (33\%) on Foreign Direct Investment (FDI) in Cameroon's context (CEMAC, AfCFTA). (3 marks)
\end{enumerate}
\tcblower
\textbf{Detailed Solution:}

\textbf{1. Pre-Tax Equilibrium}
$2000 - 0.5Q = 500 + 0.5Q \Rightarrow 1500 = Q \Rightarrow Q_0 = \textbf{1,500 (thousand litres)}$
$P_0 = 500 + 0.5(1500) = \textbf{1,250 FCFA/Litre}$
$CS_0 = \frac{1}{2} \times (2000 - 1250) \times 1500 = 0.5 \times 750 \times 1500 = \textbf{562,500 (thousand FCFA)}$
$PS_0 = \frac{1}{2} \times (1250 - 500) \times 1500 = 0.5 \times 750 \times 1500 = \textbf{562,500 (thousand FCFA)}$
\textit{Note: Symmetric linear curves $\rightarrow$ Equal Surplus.}

\textbf{2. Post-Tax Equilibrium (Tax on Producer $\rightarrow$ Supply shifts UP by 200)}
New Supply: $P_s = 500 + 0.5Q + 200 = 700 + 0.5Q$ (Price Consumer Pays $P_d$)
$2000 - 0.5Q = 700 + 0.5Q \Rightarrow 1300 = Q \Rightarrow Q_1 = \textbf{1,300}$
$P_d = 2000 - 0.5(1300) = \textbf{1,350 FCFA}$ (Consumer Price)
$P_s = 1350 - 200 = \textbf{1,150 FCFA}$ (Producer Price)
\textbf{Tax Revenue} $R = t \times Q_1 = 200 \times 1,300 = \textbf{260,000 (thousand FCFA)} = 260 \text{ Million FCFA}$.
\textbf{New CS} $= \frac{1}{2} \times (2000 - 1350) \times 1300 = 0.5 \times 650 \times 1300 = \textbf{422,500}$.
\textbf{New PS} $= \frac{1}{2} \times (1150 - 500) \times 1300 = 0.5 \times 650 \times 1300 = \textbf{422,500}$.
\textbf{DWL} $= (CS_0 + PS_0) - (CS_1 + PS_1 + R)$
$= (562,500 + 562,500) - (422,500 + 422,500 + 260,000)$
$= 1,125,000 - 1,105,000 = \textbf{20,000 (thousand FCFA)} = \textbf{20 Million FCFA}$.
\textit{Check Formula:} $\frac{1}{2} \times t \times \Delta Q = 0.5 \times 200 \times (1500-1300) = 100 \times 200 = 20,000$. \textbf{Matches.}

\textbf{3. Fiscal Multipliers}
Denominator $D = 1 - c(1-t_x) + m = 1 - 0.8(1-0.2) + 0.1 = 1 - 0.8(0.8) + 0.1 = 1 - 0.64 + 0.1 = \textbf{0.46}$.
\begin{align*}
\text{Spending Multiplier } (k_G) &= \frac{1}{D} = \frac{1}{0.46} = \textbf{2.174} \\
\text{Tax Multiplier } (k_T) &= \frac{-c}{D} = \frac{-0.8}{0.46} = \textbf{-1.739} \\
\text{Balanced Budget Multiplier } (k_{BB}) &= k_G + k_T = \frac{1-c}{D} = \frac{0.2}{0.46} = \textbf{0.435}
\end{align*}
\textbf{Effect of using Beer Tax Revenue (260M) for $\Delta G$:}
$\Delta Y = k_{BB} \times \Delta G = 0.435 \times 260 \text{ Million} = \textbf{113.1 Million FCFA}$ increase in National Income.
\textit{Note:} Expansionary because $k_G > |k_T|$. The tax leaks into savings/imports; spending injects fully.

\textbf{4. Corporate Tax (33\%) and FDI in Cameroon}
\begin{itemize}
    \item \textbf{Statutory Rate:} 30\% + 10\% Council Tax = 33\%. Higher than Nigeria (30\%), Ghana (25\%), Côte d'Ivoire (25\%), Senegal (30\%). 
    \item \textbf{Effective Rate:} Lower due to incentives (Investment Code: 5-10 year tax holidays, accelerated depreciation, export deductions). 
    \item \textbf{Impact:} High statutory rate \textbf{discourages FDI} (Location decision). Profit shifting to low-tax jurisdictions (Transfer Pricing) erodes base. 
    \item \textbf{CEMAC/AfCFTA Context:} Harmonisation pressure (CEMAC directives aim for convergence). AfCFTA increases competition for capital; Cameroon must offer \textbf{competitive effective rates} + \textbf{stability} (no arbitrary changes) + \textbf{infrastructure} (Energy, Ports) to offset tax disadvantage. 
    \item \textbf{Recommendation:} Broaden base (reduce exemptions), lower statutory rate to 25-28\% to attract manufacturing FDI (AfCFTA rules of origin require local transformation).
\end{itemize}
\end{exoresolu}

\coursec{The State Budget and the Programmed Budget}

\begin{definition}[State Budget]
The \cle{State Budget} is a legal document (Finance Law) that forecasts and authorises the State's revenue and expenditure for a fiscal year (Jan 1 -- Dec 31). It embodies the \textbf{Principle of Annualité}, \textbf{Unité} (single document), \textbf{Universalité} (gross budgeting), \textbf{Spécialité} (line-item appropriation), and \textbf{Équilibre} (balance between revenue and expenditure, though deficit allowed).
\end{definition}

\begin{methode}[The Budget Cycle: Preparation, Execution, Control and Programmed Budget (BDP)]
\textbf{Objective:} Master the 4 phases of the Budget Cycle and the logic of the \textbf{Budget de Programmation (Programmed Budget / Performance Budget)} introduced in Cameroon (Law 2018/12).
\begin{enumerate}
    \item \textbf{Phase 1: Preparation (Jan--Jun).} Circular $\rightarrow$ Ministerial Conferences $\rightarrow$ Arbitration (Minfi) $\rightarrow$ Draft Finance Law (PLF) $\rightarrow$ Council of Ministers $\rightarrow$ Parliament (Nov/Dec vote).
    \item \textbf{Phase 2: Execution (Jan--Dec).} \textit{Ordonnancement} (Commitment) $\rightarrow$ \textit{Liquidation} (Verification) $\rightarrow$ \textit{Paiement} (Payment). \textbf{Separation of Authorities:} Ordonnateur (Minister/Manager) $\neq$ Comptable (Treasurer/Paymaster). \textbf{Principle:} Annualité (Jan 1--Dec 31), Unité (Single doc), Universalité (Gross), Spécialité (Line items), Équilibre (Balance).
    \item \textbf{Phase 3: Control (Ex-post).} 
    \begin{itemize}
        \item \textbf{Administrative:} Inspection Générale des Finances (IGF), Cour des Comptes (audit).
        \item \textbf{Parliamentary:} Vote \textbf{Loi de Règlement} (Settlement Law) $\approx$ June Year N+1. Certifies accounts.
        \item \textbf{Judicial:} Chambre des Comptes (Supreme Court) sanctions mismanagement.
    \end{itemize}
    \item \textbf{Phase 4: Evaluation (Performance).} \textbf{Programmed Budget (BDP)} Logic:
    \begin{itemize}
        \item \textbf{Programmes} (Major policy: e.g., "Universal Health Coverage") $\rightarrow$ \textbf{Actions} $\rightarrow$ \textbf{Activities}.
        \item \textbf{Indicateurs de Performance:} \textit{Output} (Quantity: \# schools built), \textit{Outcome} (Quality: Enrolment rate), \textit{Impact} (Literacy rate).
        \item \textbf{Credits:} \textit{Crédits de Paiement} (Cash limit for year) vs \textit{Autorisations d'Engagement} (Multi-year commitment ceiling for investment).
    \end{itemize}
    \item \textbf{GCE Trap:} Distinguish \textbf{Virement} (Transfer within same Programme/Action, Minister signature) vs \textbf{Transfert} (Between Programmes/Ministries, PM/President decree). \textbf{Report} (Carry forward unused AE to next year) vs \textbf{Rejet} (Unused CP cancelled).
\end{enumerate}
\end{methode}

\begin{exoresolu}[Budget Cycle, Programmed Budget Logic and Virements]
\textbf{Statement:}
In the 2025 Finance Law, the \textbf{Ministry of Public Health (MINSANTÉ)} has Programme 152 "Disease Control".
\begin{itemize}
    \item \textbf{Action 1:} Malaria Control. \textit{Credit AE: 5 Billion, CP: 4 Billion}.
    \item \textbf{Action 2:} Vaccination. \textit{Credit AE: 3 Billion, CP: 3 Billion}.
\end{itemize}
During execution (July 2025), Malaria Control faces a shortage of 1 Billion FCFA (CP) due to price hike of insecticide-treated nets. Vaccination has a surplus of 1 Billion FCFA (CP) due to delayed delivery of vaccines.
\begin{enumerate}
    \item Define \textbf{Autorisation d'Engagement (AE)} and \textbf{Crédit de Paiement (CP)}. Why can AE $>$ CP for Investment? (4 marks)
    \item Can the Minister of Health transfer 1 Billion CP from Vaccination to Malaria Control by simple signature? Justify legally. (4 marks)
    \item Define \textbf{Output Indicator} and \textbf{Outcome Indicator} for "Malaria Control". Give one example each. (4 marks)
    \item If the 1 Billion CP in Vaccination is not spent by Dec 31, 2025, what happens to it? What about the unused AE? (3 marks)
    \item Explain the role of the \textbf{Loi de Règlement 2025} (voted in 2026). (2 marks)
\end{enumerate}
\tcblower
\textbf{Detailed Solution:}

\textbf{1. AE vs CP}
\begin{itemize}
    \item \textbf{Autorisation d'Engagement (AE):} The \textbf{legal ceiling} for \textbf{commitments} (Ordonnancement) a Minister can sign during the year. It represents the \textbf{multi-year contractual obligation} (e.g., signing a 3-year construction contract). For investment, AE covers the \textbf{total contract value}.
    \item \textbf{Crédit de Paiement (CP):} The \textbf{cash limit} for \textbf{actual payments} (Paiement) that can be made \textbf{within the fiscal year (Jan 1--Dec 31)}. It covers the \textbf{annual tranche} of the contract (e.g., down payment + 2025 work certificates).
    \item \textbf{Why AE $>$ CP for Investment?} Because investment projects (Dams, Roads, Hospitals) span multiple years. The State commits (AE) to the full contract now (legal obligation), but pays (CP) progressively as work is done each year. \textit{Example:} Malaria nets contract 5 Billion total (AE), but only 4 Billion payable in 2025 (CP).
\end{itemize}

\textbf{2. Transfer Legality (Virement vs Transfert)}
\textbf{NO}, the Minister \textbf{cannot} do this by simple signature.
\begin{itemize}
    \item \textbf{Virement:} Transfer \textbf{within the same Programme} (same Ministry, same Programme code). Authorised by \textbf{Ministerial Order (Arrêté Ministériel)}. Requires favourable opinion of Minfi (Budget Dept).
    \item \textbf{Transfert:} Transfer \textbf{between Programmes} (e.g., Prog 152 Action 1 $\rightarrow$ Action 2 is technically same Programme? Usually Actions are sub-programmes. If distinct "Actions" under same Programme $\rightarrow$ Virement. If distinct "Programmes" $\rightarrow$ Transfert).
    \item \textbf{Here:} Malaria (Action 1) and Vaccination (Action 2) are \textbf{different Actions} under \textbf{Programme 152}.
    \item \textbf{Rule (LOLF Art 34-35):} Virements \textbf{between Actions} of the same Programme are allowed by \textbf{Ministerial Order} (with Minfi visa). \textbf{However}, Virements \textbf{on Personnel Credits (Titre 2)} are forbidden. Virements \textbf{from Investment (Titre 3/4) to Operating (Titre 2)} forbidden.
    \item \textbf{Conclusion:} Since both are likely Operating/Investment mix (Titre 2/3), a \textbf{Virement} via \textbf{Arrêté Ministériel} (Minister + Minfi visa) is required, \textbf{NOT} a simple signature. A "Transfert" (Decree) would be needed if moving credits to a \textbf{different Ministry/Programme}.
\end{itemize}

\textbf{3. Performance Indicators (Malaria Control)}
\begin{itemize}
    \item \textbf{Output Indicator (Produit/Activité):} Measures \textbf{quantity of goods/services delivered} directly by the action.
    \textit{Example:} "\textbf{Number of Insecticide-Treated Nets (ITNs) distributed}" (Target: 5 Million). Or "\textbf{Number of households sprayed}".
    \item \textbf{Outcome Indicator (Résultat/Effet):} Measures the \textbf{intermediate effect} on the target population (behaviour/coverage).
    \textit{Example:} "\textbf{Percentage of children under 5 sleeping under ITN}" (Target: 80\%). Or "\textbf{Malaria incidence rate per 1000 population}".
    \item \textit{(Impact Indicator = Long term: Under-5 Mortality Rate).}
\end{itemize}

\textbf{4. Year-End Fate of Unused Credits (Dec 31, 2025)}
\begin{itemize}
    \item \textbf{Unused CP (Crédits de Paiement):} \textbf{Cancelled (Annulés / Rejetés)}. Principle of \textbf{Annualité}. Money returns to Treasury Single Account (CUT). Cannot be spent in 2026. \textit{Exception:} "Reports de crédits" for specific justified cases (rare, requires decree).
    \item \textbf{Unused AE (Autorisations d'Engagement):} \textbf{Reported (Reportés)} automatically to 2026 (and subsequent years until project end). The legal commitment capacity remains valid for the multi-year contract. The 1 Billion AE surplus in Vaccination stays available to sign contracts in 2026.
\end{itemize}

\textbf{5. Loi de Règlement 2025 (Voted June 2026)}
It is the \textbf{ex-post validation} of the 2025 budget execution.
\begin{itemize}
    \item \textbf{Function:} Parliament (Assemblée Nationale) examines the \textbf{Compte Administratif} (Minfi's account) and \textbf{Compte de Gestion} (Treasurer's account) certified by \textbf{Cour des Comptes}.
    \item \textbf{Outcome:} Votes the \textbf{Loi de Règlement} which: (1) Approves/Rejects accounts, (2) Grants \textbf{Quitus} (Discharge) to Ordonnateurs/Comptables, (3) Certifies the \textbf{Deficit/Surplus} and \textbf{Debt} level for 2025, (4) Authorises \textbf{Reports d'AE} officially. It closes the budget cycle for Year N.
\end{itemize}
\end{exoresolu}

\coursec{The Public Debt and the Effectiveness of Fiscal Policy in Cameroon}

\begin{methode}[Public Debt Dynamics: Sustainability, Indicators and Restructuring]
\textbf{Objective:} Calculate Debt/GDP, Debt Service/Revenue, Grant Element; assess Sustainability (Solvency vs Liquidity); understand HIPC/CCRF, Eurobonds, and Domestic vs External debt.
\begin{enumerate}
    \item \textbf{Key Stock Indicators (Solvency):}
    \begin{itemize}
        \item \textbf{Debt / GDP:} Thresholds (IMF DSA): Low income $\approx$ 50\% (PV), Market access $\approx$ 70\%. Cameroon $\approx$ 45-50\% (Gross), PV $\approx$ 35-40\%. \textit{Trend rising.}
        \item \textbf{Debt / Revenue:} Better for repayment capacity. Threshold $\approx$ 300\% (PV). Cameroon $\approx$ 250-300\%.
        \item \textbf{Present Value (PV) of Debt:} Discount future flows at \textbf{Unified Discount Rate (5\%)} (IMF/WEO). Concessional loans (IDA, AfDB) have low PV (High Grant Element). Commercial loans (Eurobonds, Syndicated) have PV $\approx$ Face Value.
    \end{itemize}
    \item \textbf{Key Flow Indicators (Liquidity):}
    \begin{itemize}
        \item \textbf{Debt Service / Revenue:} Threshold $\approx$ 18-22\%. Cameroon often $\approx$ 20-25\% (Pressure). Includes Principal + Interest.
        \item \textbf{Debt Service / Exports:} Threshold $\approx$ 15-20\%. Critical for FX availability.
    \end{itemize}
    \item \textbf{Grant Element (GE):} $GE = 1 - \frac{PV}{Face Value}$. Concessional if $GE \ge 35\%$ (IDA: $\approx$ 50-60\%). Eurobond $GE \approx 0\%$.
    \item \textbf{Debt Dynamics Equation:} $\Delta (D/Y) \approx (r - g) \frac{D}{Y} + \frac{PD}{Y}$. 
    \begin{itemize}
        \item $r$ = Avg real interest rate, $g$ = Real GDP growth.
        \item If $r > g$ (Cameroon recently: $r \approx 3-4\%$, $g \approx 3-4\%$ $\rightarrow$ tight), Primary Surplus ($PD<0$) needed to stabilise ratio.
    \end{itemize}
    \item \textbf{Restructuring Tools:} Reprofiling (Maturity extension), Haircut (Principal reduction - rare for sovereign), Debt-for-Nature/Development Swaps. \textbf{Common Framework (G20)} for LICs (Chad, Zambia, Ethiopia $\rightarrow$ Cameroon not eligible yet but monitored).
\end{enumerate}
\end{methode}

\begin{exoresolu}[Debt Sustainability Analysis (DSA) for a Hypothetical CEMAC Country]
\textbf{Statement:}
\textbf{Country "Equatoria"} (CEMAC member, CFA Franc zone). 2024 Data:
\begin{itemize}
    \item Nominal GDP: 20,000 Billion FCFA.
    \item Total Public Debt (Face Value): 10,000 Billion FCFA.
    \item \textbf{Composition:} 
    \begin{itemize}
        \item Multilateral (IDA, AfDB, IMF): 4,000 Billion (Avg Interest 1\%, Maturity 30y, Grace 10y).
        \item Bilateral (Paris Club, China): 2,000 Billion (Avg Interest 2\%, Maturity 20y).
        \item Commercial / Eurobond: 2,500 Billion (Coupon 6.5\%, Maturity 10y, Bullet).
        \item Domestic (T-Bills/Bonds BEAC market): 1,500 Billion (Avg Yield 5.5\%, Avg Maturity 3y).
    \end{itemize}
    \item Budget Revenue (Excl Grants): 2,000 Billion FCFA.
    \item Exports (Goods/Services): 4,000 Billion FCFA.
    \item Debt Service Paid in 2024: 450 Billion FCFA (Int: 300, Principal: 150).
    \item Real GDP Growth ($g$): 4.0\%. Avg Nominal Interest ($i$): 3.5\%. Inflation ($\pi$): 2.5\%.
\end{itemize}
\begin{enumerate}
    \item Calculate \textbf{Debt/GDP}, \textbf{Debt/Revenue}, \textbf{Debt Service/Revenue}, \textbf{Debt Service/Exports}. Assess against IMF LIC DSA Thresholds (PV Debt/GDP 50\%, PV Debt/Rev 300\%, DS/Rev 18\%, DS/Exp 15\%). Assume PV $\approx$ 85\% Face Value (Concessional share high). (6 marks)
    \item Calculate the \textbf{Grant Element} of the Eurobond (Face 2,500, Coupon 6.5\%, 10y Bullet, Discount Rate 5\%). Is it concessional? (5 marks)
    \item Use the \textbf{Debt Dynamics Equation} to estimate the \textbf{Primary Balance (\% GDP)} required to \textbf{stabilise Debt/GDP} at current level (50\%). (4 marks)
    \item Discuss the specific \textbf{Rollover Risk} of the Domestic Debt profile (Short maturity 3y) vs Eurobond (Bullet 10y) in the CEMAC context (BEAC refinancing). (3 marks)
\end{enumerate}
\tcblower
\textbf{Detailed Solution:}

\textbf{1. Debt Indicators (Face Value vs PV)}
\textbf{Face Value Ratios:}
\begin{align*}
\text{Debt/GDP} &= 10,000 / 20,000 = \textbf{50.0\%} \\
\text{Debt/Revenue} &= 10,000 / 2,000 = \textbf{500\%} \\
\text{DS/Revenue} &= 450 / 2,000 = \textbf{22.5\%} \\
\text{DS/Exports} &= 450 / 4,000 = \textbf{11.25\%}
\end{align*}
\textbf{Present Value (PV) Ratios (PV = 0.85 $\times$ Face = 8,500 Billion):}
\begin{align*}
\text{PV Debt/GDP} &= 8,500 / 20,000 = \textbf{42.5\%} \quad (\textbf{< 50\% Threshold} \rightarrow \textbf{Moderate Risk}) \\
\text{PV Debt/Revenue} &= 8,500 / 2,000 = \textbf{425\%} \quad (\textbf{> 300\% Threshold} \rightarrow \textbf{High Risk / Breach}) \\
\text{DS/Revenue} &= \textbf{22.5\%} \quad (\textbf{> 18-22\% Threshold} \rightarrow \textbf{High Risk / Liquidity Pressure}) \\
\text{DS/Exports} &= \textbf{11.25\%} \quad (\textbf{< 15\% Threshold} \rightarrow \textbf{Moderate Risk})
\end{align*}
\textbf{Assessment:} \textbf{High Overall Risk of Debt Distress.} While PV/GDP is okay, the \textbf{Revenue-based indicators are breached}. Equatoria has a \textbf{Revenue Mobilisation Problem} (Tax/GDP = 10\%). High Debt Service (22.5\% Rev) crowds out priority spending. Liquidity risk is acute.

\textbf{2. Grant Element of Eurobond}
\textbf{Parameters:} Face $F=2500$, Coupon $c=6.5\%$ (Annual 162.5), Maturity $n=10$, Bullet (Principal at end), Discount Rate $d=5\%$.
\textbf{PV Calculation:}
\[
PV = \sum_{t=1}^{10} \frac{162.5}{(1.05)^t} + \frac{2500}{(1.05)^{10}}
\]
Annuity Factor (10y, 5\%) $= \frac{1 - 1.05^{-10}}{0.05} = 7.7217$.
Discount Factor (10y) $= 1.05^{-10} = 0.6139$.
\[
PV = 162.5 \times 7.7217 + 2500 \times 0.6139 = 1,254.8 + 1,534.8 = \textbf{2,789.6 Billion}
\]
\textit{Wait:} Coupon (6.5\%) > Discount Rate (5\%) $\rightarrow$ PV $>$ Face Value. This is a \textbf{Premium Bond}.
$PV \approx 2,790$ Billion. Face = 2,500.
\textbf{Grant Element} $GE = 1 - \frac{PV}{Face} = 1 - \frac{2790}{2500} = 1 - 1.116 = \textbf{-11.6\%}$.
\textbf{Conclusion:} \textbf{Negative Grant Element.} Not concessional (Threshold $\ge$ 35\%). It is a \textbf{Commercial Loan} at market terms (actually slightly expensive vs 5\% benchmark). Increases PV of debt stock significantly.

\textbf{3. Primary Balance for Debt Stabilisation}
\textbf{Debt Dynamics:} $\Delta (d) \approx (r - g) d + p$ where $d = D/Y$, $p = PD/Y$ (Primary Deficit $>0$, Surplus $<0$).
Target: $\Delta d = 0$ (Stabilise at $d=50\%$).
Real Interest Rate $r = i - \pi = 3.5\% - 2.5\% = \textbf{1.0\%}$.
Real Growth $g = \textbf{4.0\%}$.
$r - g = 1.0\% - 4.0\% = \textbf{-3.0\%}$.
\[
0 = (-0.03) \times 0.50 + p \Rightarrow p = 0.015 = \textbf{1.5\% of GDP (Primary Surplus)}
\]
\textbf{Interpretation:} Because $g > r$ (Favourable growth-interest differential), Equatoria can run a \textbf{Primary Deficit of 1.5\% GDP} and still stabilise Debt/GDP at 50\%. 
\textit{Current Reality Check:} If current Primary Deficit is $> 1.5\%$ GDP (likely given DS/Rev pressure), Debt/GDP will \textbf{rise}. Need fiscal consolidation.

\textbf{4. Rollover Risk: Domestic (3y) vs Eurobond (10y Bullet)}
\begin{itemize}
    \item \textbf{Domestic Debt (1,500B, 3y Avg Mat):} \textbf{High Rollover Risk.} $\approx$ 500 Billion FCFA matures \textbf{every year}. Must be refinanced on BEAC market (Weekly T-Bill auctions, Monthly Bond syndications). 
    \begin{itemize}
        \item \textbf{Risk:} If BEAC tightens liquidity (raises Taux Directeur / Taux de Pension) to defend CFA parity / fight inflation $\rightarrow$ Yields spike $\rightarrow$ Roll-over cost $\uparrow$ sharply. 
        \item \textbf{Crowding Out:} High domestic borrowing absorbs bank liquidity $\rightarrow$ Credit to Private Sector $\downarrow$ (Investment $\downarrow$).
    \end{itemize}
    \item \textbf{Eurobond (2,500B, 10y Bullet):} \textbf{ZERO Rollover Risk for 10 years.} No principal repayment until 2034. Only coupon payments (162.5B/yr). 
    \begin{itemize}
        \item \textbf{Risk:} \textbf{FX Risk} (CFA/EUR fixed, but USD/EUR moves $\rightarrow$ if bond in USD, huge risk. Assumed EUR). \textbf{Refinancing Risk in 2034} (Bullet repayment 2,500B = 12.5\% GDP!). Must start sinking fund / liability management now.
        \item \textbf{Market Access Risk:} If Cameroon loses market access (rating downgrade), cannot issue new Eurobond to refinance 2034 maturity.
    \end{itemize}
    \item \textbf{Strategy:} Lengthen Domestic Maturity (Issue 5-7y, 10y Bonds on UMOA/CEMAC market), Build Sinking Fund for Eurobond, Improve Tax Revenue to reduce borrowing need.
\end{itemize}
\end{exoresolu}

\coursec{Effectiveness of Fiscal Policy in Cameroon}

\begin{methode}[Evaluating Fiscal Policy Effectiveness in Cameroon]
\textbf{Objective:} Structure an essay answer on "Effectiveness of Fiscal Policy in Cameroon" using the \textbf{Transmission Channels}, \textbf{Constraints (Structural/Institutional)}, and \textbf{Reforms}.
\begin{enumerate}
    \item \textbf{Transmission Channels (How it \textit{should} work):}
    \begin{itemize}
        \item \textbf{AD Channel:} $G \uparrow / T \downarrow \rightarrow AD \uparrow \rightarrow Y \uparrow$ (Multiplier). $G \downarrow / T \uparrow \rightarrow AD \downarrow \rightarrow Inflation \downarrow$.
        \item \textbf{Supply Side:} Productive $G$ (Infrastructure, Education) $\rightarrow$ LRAS $\uparrow$ (Potential Growth).
        \item \textbf{Expectations/Credibility:} Rules-based policy (WAEMU/CEMAC Convergence Criteria: Deficit $\le$ 3\% GDP, Debt $\le$ 70\% GDP, Inflation $\le$ 3\%, Arrears=0) $\rightarrow$ Anchors expectations $\rightarrow$ Lower Risk Premium $\rightarrow$ Investment $\uparrow$.
    \end{itemize}
    \item \textbf{Major Constraints (Why it \textit{fails} / is weak in Cameroon):}
    \begin{enumerate}
        \item \textbf{Narrow Tax Base / Low Revenue:} Tax/GDP $\approx$ 12-13\%. High Informality ($\approx$ 90\% employment), Exemptions (Investment Code, Agreements), Evasion. $\rightarrow$ Low Fiscal Space. Reliance on Oil/Customs (Volatile).
        \item \textbf{Rigid Expenditure Structure:} Wage Bill ($\approx$ 35-40\% Rev), Debt Service ($\approx$ 20-25\% Rev), Subsidies (Fuel $\approx$ 500-1000B/yr) $\rightarrow$ \textbf{Crowding out Capital Expenditure (CAPEX)}. CAPEX/Total Exp often $< 20\%$.
        \item \textbf{Execution Bottlenecks:} Slow procurement (Marchés Publics), Cash rationing (Trésorerie), "Comptabilité Matières" weak $\rightarrow$ Low Absorption Capacity (Budget execution rate $\approx$ 80-85\%, but CAPEX execution $\approx$ 60-70\%).
        \item \textbf{CEMAC Straitsjacket:} Fixed Exchange Rate $\rightarrow$ No Monetary Policy autonomy. Fiscal Policy \textbf{MUST} do the adjustment. Convergence Criteria often breached (Deficit $> 3\%$, Arrears accumulation). No Lender of Last Resort for State (BEAC Statutes Art. 23 - no direct financing).
        \item \textbf{Governance/Leakages:} Ghost workers, Over-invoicing, Weak PFM (Public Financial Management) - PEFA scores low on predictability/control.
        \item \textbf{External Shocks:} Oil Price, COVID, Russia-Ukraine (Wheat/Fertilizer/Fuel prices), Climate (Floods North, Droughts).
    \end{enumerate}
    \item \textbf{Recent Reforms (Effectiveness Boosters):}
    \begin{itemize}
        \item \textbf{Programmed Budget (BDP) 2013/2018:} Shift from Means $\rightarrow$ Results.
        \item \textbf{Treasury Single Account (CUT / Compte Unique du Trésor):} Cash management, reduces idle balances/borrowing cost.
        \item \textbf{Digitalisation:} E-tax (Tele-declaration), E-procurement (COLEPS), ASYCUDA World (Customs), SIGIF (Forestry).
        \item \textbf{Medium Term Expenditure Framework (CDMT / MTEF):} 3-year planning.
        \item \textbf{Debt Management:} MTDS (Medium Term Debt Strategy), DMFAS, Eurobond liability management (Buybacks/Switches).
    \end{itemize}
    \item \textbf{Evaluation Grid (For Essay Conclusion):} 
    \begin{center}
    \begin{tabularx}{\linewidth}{|l|X|X|}
    \hline
    \rowcolor{popblueL} \textbf{Dimension} & \textbf{Assessment} & \textbf{Evidence} \\
    \hline
    \textbf{Stabilisation (Counter-cyclical)} & \textbf{Low / Pro-cyclical often} & Deficit cuts in recession (IMF programs); Fuel subsidies pro-cyclical. \\
    \hline
    \textbf{Allocation (Efficiency)} & \textbf{Moderate / Improving} & CAPEX share rising but execution low; "White elephants" (stadiums vs rural roads). \\
    \hline
    \textbf{Distribution (Equity)} & \textbf{Very Low} & Regressive tax structure; Subsidies benefit rich (fuel); Social spending low. \\
    \hline
    \textbf{Growth (Supply Side)} & \textbf{Moderate} & Infrastructure stock rising (Ports, Energy, Roads) but Human Capital (Health/Edu outcomes) lagging; Business climate constraints limit private crowding-in. \\
    \hline
    \end{tabularx}
    \end{center}
\end{enumerate}
\end{methode}

\begin{exoresolu}[Essay: "Assess the Effectiveness of Fiscal Policy in Cameroon since 2010"]
\textbf{Statement:}
"Despite the adoption of the Programmed Budget and various digitalisation reforms, fiscal policy in Cameroon has struggled to achieve its objectives of macroeconomic stability, equitable redistribution, and sustainable growth." 
\textbf{Discuss this assertion, analysing the transmission channels, binding constraints, and recent reform efforts.} (25 marks - Essay Style)

\tcblower
\textbf{Detailed Solution (Essay Plan \& Key Arguments)}

\textbf{Introduction}
\begin{itemize}
    \item \textbf{Hook:} Cameroon's fiscal policy operates within the \textbf{CEMAC fixed exchange rate regime} (CFA Franc), making it the \textbf{sole macroeconomic stabilisation instrument} (Mundell-Fleming: Monetary policy ineffective/constrained).
    \item \textbf{Definitions:} Fiscal Policy = Use of Govt Revenue (Taxes, Oil Rent, Grants) and Expenditure (Current, Capital) to influence AD/AS.
    \item \textbf{Thesis:} While the \textbf{legal/technical framework} has modernised significantly (BDP, CUT, Digitalisation, CEMAC Convergence), \textbf{structural rigidities} (Revenue narrowness, Expenditure rigidity, Governance leaks, External shocks) have severely blunted the \textbf{transmission mechanism}, resulting in \textbf{pro-cyclicality, low multipliers, and missed development targets}.
    \item \textbf{Plan:} 1. Theoretical Channels \& Objectives. 2. Binding Constraints (Revenue, Expenditure, Institutional, External). 3. Reform Efforts \& Results. 4. Synthesis/Evaluation.
\end{itemize}

\textbf{Body Paragraph 1: Transmission Channels \& Objectives (The "What")} 
\begin{itemize}
    \item \textbf{Stabilisation (AD Management):} Counter-cyclical impulse. Multiplier $k = \frac{1}{1-c(1-t)+m}$. In Cameroon: $c \approx 0.7$, $t \approx 0.15$, $m \approx 0.3$ (High import propensity) $\rightarrow$ $k \approx \frac{1}{1-0.7(0.85)+0.3} = \frac{1}{0.705} \approx 1.4$. \textbf{Low Multiplier} due to high leakage (Imports/Savings).
    \item \textbf{Allocation (Efficiency):} Shift from Current $\rightarrow$ Capital Expenditure. Target: CAPEX/Total Exp $> 30\%$ (WAEMU norm). Reality: Often $< 20\%$.
    \item \textbf{Distribution (Equity):} Progressive Tax + Targeted Transfers. Reality: Regressive Tax Structure (Indirect 64\%), Universal Fuel Subsidy (Rich benefit more), Weak Social Registry.
    \item \textbf{Growth (Supply Side):} Public Investment $\rightarrow$ Infrastructure $\rightarrow$ Productivity $\uparrow$. "Big Push" (Kribi Port, Lom Pangar, Autoroute Yaoundé-Douala).
\end{itemize}

\textbf{Body Paragraph 2: Binding Constraints (The "Why it Fails")} 
\begin{enumerate}
    \item \textbf{Revenue Constraint (The "Fiscal Space" Trap):} 
    \begin{itemize}
        \item Tax/GDP $\approx$ 12-13\% (SSA Avg 15-17\%, WAEMU Target 20\%). 
        \item \textit{Causes:} Informal Sector ($\approx$ 90\% jobs, $\approx$ 30-40\% GDP), Narrow Base (Formal sector $\approx$ 500 large firms pay 80\% IS), Tax Expenditures (Exemptions $\approx$ 2-3\% GDP/yr), Administration Capacity (DGI/DGD staffing).
        \item \textit{Consequence:} High reliance on Oil Revenue (Volatile) $\rightarrow$ \textbf{Pro-cyclicality}: Spend windfall in boom, brutal cut in bust (2014-2016, 2020).
    \end{itemize}
    \item \textbf{Expenditure Rigidity (The "Crowding Out"):} 
    \begin{itemize}
        \item Wage Bill: $\approx$ 38\% Tax Rev (WAEMU Norm $\le$ 35\%). Hiring freezes ineffective (demographic pressure).
        \item Debt Service: $\approx$ 22\% Rev (Rising fast due to Eurobonds 2015, 2017, 2021 + Domestic costly).
        \item Fuel Subsidy: $\approx$ 600-1000B/yr (2-4\% GDP). Untargeted. Crowds out Health/Education/CAPEX.
    \end{itemize}
    \item \textbf{Execution \& Governance (The "Leakage"):} 
    \begin{itemize}
        \item Low CAPEX Execution Rate ($\approx$ 65-75\%). Causes: Procurement delays (ARMP procedures), Cash Rationing (Trésorerie prioritises Wages/Debt), Weak Project Preparation (Feasibility studies missing).
        \item PFM Weaknesses (PEFA 2022): Low score on \textbf{Predictability} (In-year budget adjustments frequent), \textbf{Control} (Arrears accumulation $\approx$ 2-3\% GDP/yr), \textbf{Accounting} (Accrual adoption slow).
    \end{itemize}
    \item \textbf{CEMAC Institutional Straitsjacket:} 
    \begin{itemize}
        \item Convergence Criteria (Deficit $\le$ 3\%, Debt $\le$ 70\%, Inflation $\le$ 3\%, Arrears=0, Non-accumulation Domestic/External Arrears). 
        \item \textbf{Systemic Breach:} Deficit often $> 3\%$ (2023: $\approx$ 3.5-4\%), Arrears persistent. 
        \item \textbf{No Monetary Financing:} BEAC Statutes Art 23 $\rightarrow$ Hard Budget Constraint $\rightarrow$ Sudden Stops / Cash Rationing if markets close.
    \end{itemize}
\end{enumerate}

\textbf{Body Paragraph 3: Reform Efforts \& Marginal Gains (The "What is Changing")} 
\begin{itemize}
    \item \textbf{Budget de Programmation (BDP) / LOLF 2018:} Introduced Performance Logic (Programmes $\rightarrow$ Actions $\rightarrow$ Indicators). \textit{Limit:} Still input-focused in practice; Indicators often "Output" not "Outcome"; Parliament oversight weak on performance.
    \item \textbf{Cash Management Reform (CUT / TSA):} Single Treasury Account (BEAC). Reduced idle balances, lowered borrowing costs, improved cash forecasting. \textit{Success:} Visible reduction in Treasury Bills stock cost.
    \item \textbf{Digitalisation:}\begin{itemize}
        \item \textbf{DGI:} Tele-declaration (IRPP, IS, TVA), E-invoicing (Facture Normalisée) rollout $\rightarrow$ VAT buoyancy improving.
        \item \textbf{DGD:} ASYCUDA World, Scanner deployment (Douala/Kribi), Electronic Single Window (GUCE) $\rightarrow$ Customs revenue resilience. \textit{Limit:} Connectivity/Power issues inland; Resistance to change; Informal sector largely untouched.
    \end{itemize}
    \item \textbf{Debt Management (MTDS):} Shift to longer maturities (10-12y domestic bonds), Eurobond liability management (2021 buyback of 2025 maturity). \textit{Success:} Reduced rollover risk, extended profile. \textit{Risk:} Rising interest burden (Domestic yields $\uparrow$ due to BEAC tightening).
    \item \textbf{IMF Programs (ECF/EFF 2021-2024):} Anchored consolidation (Fuel subsidy reform $\rightarrow$ targeted cash transfers "Filets Sociaux", Tax administration gains). \textit{Result:} Deficit narrowed from 3.6\% (2021) to $\approx$ 1.5\% (2023 est.), but at cost of social tension (fuel price hikes) and compressed CAPEX.
\end{itemize}

\textbf{Body Paragraph 4: Synthesis \& Evaluation (The "Verdict")} 
\begin{center}
\begin{tabularx}{\linewidth}{|l|X|X|}
\hline
\rowcolor{popblueL} \textbf{Objective} & \textbf{Effectiveness} & \textbf{Key Determinant} \\
\hline
\textbf{Stabilisation} & \textbf{Low / Pro-cyclical} & Revenue volatility (Oil) + Rigid Expenditure + CEMAC Rules $\rightarrow$ Forced austerity in downturns. \\
\hline
\textbf{Allocation (Efficiency)} & \textbf{Moderate / Improving} & BDP/CUT help; but Low CAPEX execution \& Project selection politicisation (White elephants) persist. \\
\hline
\textbf{Distribution (Equity)} & \textbf{Very Low} & Regressive Tax + Universal Subsidies + Weak Targeting (Social Registry coverage $< 20\%$ poor). \\
\hline
\textbf{Growth (Supply Side)} & \textbf{Moderate} & Infrastructure stock rising (Ports, Energy, Roads) but Human Capital (Health/Edu outcomes) lagging; Business climate constraints (Doing Business) limit private crowding-in. \\
\hline
\end{tabularx}
\end{center}

\textbf{Conclusion}
\begin{itemize}
    \item Fiscal policy in Cameroon is \textbf{necessary but insufficient}. The \textbf{instrument} (Budget) is structurally constrained by a \textbf{narrow revenue base} and \textbf{rigid current spending}, amplified by the \textbf{CEMAC fixed exchange rate} which forbids monetary accommodation.
    \item Reforms (BDP, Digitalisation, CUT, MTDS) are \textbf{necessary conditions} for effectiveness (they fix the "plumbing"), but \textbf{not sufficient}. 
    \item \textbf{Game Changers Required:} 
    \begin{enumerate}
        \item \textbf{Tax Transition:} Shift from Oil/Customs $\rightarrow$ Broad-based Domestic Consumption (VAT efficiency) + Income/Property (Progressivity). Formalise Informal (Simplified Regime enforcement).
        \item \textbf{Subsidy Reform:} Full liberalisation of Fuel prices + Scaling up \textbf{Filets Sociaux} (Digital ID + Mobile Money targeting).
        \item \textbf{Wage Bill Rule:} Strict ceiling (35\% Rev) + Integrated HR/Payroll (SIGIPES) to stop ghost workers.
        \item \textbf{Investment Quality:} Mandatory Cost-Benefit Analysis (CBA) for all projects $> 5$B FCFA; Independent Project Appraisal Unit (Prime Minister's Office).
        \item \textbf{Regional Coordination:} Push for CEMAC "Fiscal Union" elements: Common Tax Policy (VAT/Excise harmonisation), Regional Stabilisation Fund (Rainy day fund), BEAC Lender of Last Resort facility (conditional on reforms).
    \end{enumerate}
    \item \textbf{Final Judgement:} Effectiveness remains \textbf{sub-optimal}. The "fiscal multiplier" is low, the "fiscal space" is tight, and the "fiscal contract" (Taxes for Services) is broken. Without a \textbf{Second Generation Structural Transformation} of the Revenue-Expenditure architecture, fiscal policy will remain a tool of \textbf{short-term crisis management} rather than \textbf{long-term development steering}.
\end{itemize}

\textbf{Examiner's Marking Guide (Indicative):}
\begin{itemize}
    \item \textbf{Introduction/Definitions/Plan:} 3 marks.
    \item \textbf{Transmission Channels (Theory + Cameroon params):} 4 marks.
    \item \textbf{Constraints Analysis (Revenue, Expenditure, Execution, CEMAC):} 8 marks (2 each, need specific data/examples).
    \item \textbf{Reforms Assessment (BDP, Digital, Debt, IMF):} 5 marks (Balanced view: Gains vs Limits).
    \item \textbf{Synthesis Table / Evaluation Grid:} 3 marks.
    \item \textbf{Conclusion / Policy Recommendations:} 2 marks.
    \item \textbf{Total:} 25 marks.
\end{itemize}
\end{exoresolu}

\newpage

\coursec{Exercises}

\courssub{I check my knowledge}

\begin{exercice}[MCQ: Public finance components]
Which of the following is NOT a component of public finance? 
\begin{enumerate}
\item Public revenue
\item Public expenditure
\item Public debt
\item Private investment
\end{enumerate}
\end{exercice}

\begin{exercice}[True/False with justification]
State whether each statement is true or false and justify your answer briefly.
\begin{enumerate}
\item A progressive tax system always reduces income inequality.
\item Value Added Tax (VAT) in Cameroon is a direct tax.
\item The programmed budget in Cameroon is prepared on a multi‑annual basis.
\item A budget deficit necessarily leads to an increase in the public debt.
\end{enumerate}
\end{exercice}

\begin{exercice}[Definitions]
Define the following terms and give a Cameroonian example for each:
\begin{enumerate}
\item Public finance
\item Direct tax
\item Indirect tax
\item Budget deficit
\end{enumerate}
\end{exercice}

\begin{exercice}[Direct calculation: VAT]
A wholesaler in Douala purchases goods for 1\,200\,000 FCFA (VAT‑exclusive) and sells them for 1\,800\,000 FCFA (VAT‑exclusive). The VAT rate in Cameroon is 19.25\%. Calculate:
\begin{enumerate}
\item Output VAT
\item Input VAT
\item VAT payable to the State
\end{enumerate}
\end{exercice}

\courssub{I practise}

\begin{exercice}[Tax classification and evaluation]
Distinguish between direct and indirect taxes. Provide two Cameroonian examples of each category and state one advantage and one disadvantage for each type.
\end{exercice}

\begin{exercice}[Tax incidence using supply‑demand diagram]
Using a supply and demand diagram, illustrate the effect of an indirect tax of 500 FCFA per unit on a good. Show the new equilibrium price, the price received by producers, and explain how the tax burden is shared between consumers and producers. (A sketch with labelled axes and curves is sufficient.)
\end{exercice}

\begin{exercice}[Progressive income tax computation]
Cameroon’s personal income tax brackets (annual) are as follows:
\begin{center}
\begin{tabular}{|c|c|}
\hline
\textbf{Taxable income (FCFA)} & \textbf{Rate} \\ \hline
0 – 2\,000\,000 & 10\% \\ \hline
2\,000\,001 – 5\,000\,000 & 15\% \\ \hline
5\,000\,001 – 10\,000\,000 & 20\% \\ \hline
Above 10\,000\,000 & 25\% \\ \hline
\end{tabular}
\end{center}
A taxpayer earns 7\,500\,000 FCFA per year. Compute:
\begin{enumerate}
\item Total tax payable
\item Average tax rate
\item Marginal tax rate
\end{enumerate}
\end{exercice}

\begin{exercice}[Budgetary process and principles]
Outline the main stages of the budgetary process in Cameroon. Then state four fundamental principles that govern the State budget.
\end{exercice}

\courssub{I prepare for the GCE}

\begin{exercice}[Structural question: Public finance components (10 marks)]
Define public finance and distinguish its four main components. Illustrate each component with a concrete example from the Cameroonian economy.
\end{exercice}

\begin{exercice}[Essay: Effectiveness of fiscal policy in Cameroon (20 marks)]
\textbf{``Fiscal policy has been effective in Cameroon.''} Discuss this statement. In your answer, present the achievements of fiscal policy, its limitations, and suggest measures to improve its effectiveness. (GCE Paper 2 style)
\end{exercice}

\begin{exercice}[Public debt analysis (15 marks)]
\begin{enumerate}
\item Distinguish between a budget deficit and the public debt. (4 marks)
\item Explain two major causes of the rise in Cameroon’s public debt. (6 marks)
\item Analyse two economic consequences of a high public debt for Cameroon. (5 marks)
\end{enumerate}
\end{exercice}

\newpage

\coursec{Solutions to the exercises}

\begin{corrige}[Exercise 1 — MCQ: Public finance components]
The correct answer is \textbf{d) Private investment}.

\textbf{Justification:} Public finance is the branch of economics that studies the financial activities of the \cle{State} and other public authorities. Its four classical components are:
\begin{enumerate}
\item \cle{Public revenue} (taxes, duties, fees, non-tax revenue such as oil royalties);
\item \cle{Public expenditure} (spending on administration, investment, transfers);
\item \cle{The budget} (the annual instrument that authorises revenue and expenditure);
\item \cle{Public debt} (borrowing used to finance deficits).
\end{enumerate}
\cle{Private investment} belongs to the \emph{private} sector (households and firms); it is studied in microeconomics and macroeconomics, not in public finance. Hence option (d) is the element that is \textbf{NOT} a component of public finance.
\end{corrige}

\begin{corrige}[Exercise 2 — True/False with justification]
\begin{enumerate}
\item \textbf{``A progressive tax system always reduces income inequality.'' — FALSE.} A progressive tax (rate rising with income) \emph{tends} to reduce inequality, but the word ``always'' makes the statement false. The redistributive effect can be neutralised by tax evasion and avoidance, generous exemptions for high incomes, loopholes, or by regressive indirect taxes (such as VAT) that weigh more heavily on the poor. Progressivity reduces inequality only if the tax is actually collected and the revenue is redistributed through social spending.
\item \textbf{``VAT in Cameroon is a direct tax.'' — FALSE.} VAT is an \cle{indirect tax}: it is levied on \emph{consumption expenditure}, collected by sellers (the taxable persons) but economically borne by the final consumer. Direct taxes (personal income tax, company tax) are paid directly by the person on whom they are imposed.
\item \textbf{``The programmed budget in Cameroon is prepared on a multi-annual basis.'' — TRUE.} The programmed (programme) budget links appropriations to multi-annual programmes and objectives of performance, typically framed over a three-year medium-term expenditure framework, instead of the traditional annual line-by-line (means-based) budget.
\item \textbf{``A budget deficit necessarily leads to an increase in the public debt.'' — FALSE.} A deficit increases the debt only if it is financed by \emph{borrowing}. It can also be financed without new debt: by drawing on Treasury reserves, by privatisation receipts, by grants, or (where permitted) by monetary financing. Borrowing is the \emph{usual} method, but it is not \emph{necessary}.
\end{enumerate}
\end{corrige}

\begin{corrige}[Exercise 3 — Definitions]
\begin{enumerate}
\item \textbf{Public finance:} the branch of economics that studies how the State and other public authorities raise revenue, allocate expenditure, prepare the budget and manage public borrowing in order to achieve economic and social objectives. \emph{Cameroonian example:} the annual State budget voted by the National Assembly in Yaoundé, which fixes expected revenue (taxes, oil royalties) and expenditure (salaries, roads, hospitals).
\item \textbf{Direct tax:} a tax imposed directly on the income or wealth of a person (natural or legal) and paid directly by that person to the tax authority; the taxpayer and the tax bearer are the same. \emph{Cameroonian example:} the personal income tax (IRPP) deducted from the salary of a civil servant in Bamenda, or the company tax paid by a firm such as a brewery in Douala.
\item \textbf{Indirect tax:} a tax levied on goods, services or transactions; it is collected by an intermediary (seller, importer) but its burden is shifted to the final consumer through prices. \emph{Cameroonian example:} VAT at 19.25\% charged on goods sold in a supermarket, or customs duties collected at the Port of Douala on imported vehicles.
\item \textbf{Budget deficit:} the excess of total public expenditure over total public revenue during a budgetary year. \emph{Cameroonian example:} when the State spends more on infrastructure (roads, stadiums, energy) and salaries than it collects in taxes and oil revenue, the gap is covered by borrowing on the CEMAC market or from external lenders.
\end{enumerate}
\end{corrige}

\begin{corrige}[Exercise 4 — Direct calculation: VAT]
Recall the VAT mechanism: the firm collects \cle{output VAT} on its sales and deducts the \cle{input VAT} paid on its purchases; only the difference is paid to the State. Rate: $t = 19.25\% = 0.1925$.

\begin{enumerate}
\item \textbf{Output VAT} (collected on sales):
\[ \text{Output VAT} = 1\,800\,000 \times 0.1925 = 346\,500 \text{ FCFA}. \]
\item \textbf{Input VAT} (paid on purchases, deductible):
\[ \text{Input VAT} = 1\,200\,000 \times 0.1925 = 231\,000 \text{ FCFA}. \]
\item \textbf{VAT payable to the State}:
\[ \text{VAT payable} = 346\,500 - 231\,000 = 115\,500 \text{ FCFA}. \]
\end{enumerate}

\textbf{Interpretation:} the wholesaler pays to the Treasury only the tax on the \emph{value added}, i.e. $19.25\% \times (1\,800\,000 - 1\,200\,000) = 0.1925 \times 600\,000 = 115\,500$ FCFA. Both methods give the same result, which confirms that VAT is a tax on value added.

\textbf{Common mistake:} Do not compute VAT on the VAT-inclusive price, and do not pay the full output VAT: forgetting to deduct input VAT is the classic error in this calculation.
\end{corrige}

\begin{corrige}[Exercise 5 — Tax classification and evaluation]
\textbf{Distinction.} A \cle{direct tax} is imposed on the income or wealth of a person and paid directly by that person to the tax authority: the legal taxpayer and the economic bearer coincide. An \cle{indirect tax} is levied on goods, services or transactions; it is collected by an intermediary (trader, importer) who shifts the burden to the consumer through the selling price.

\begin{center}
\begin{tabularx}{\linewidth}{|l|X|X|}
\hline
\rowcolor{popblueL} & \textbf{Direct taxes} & \textbf{Indirect taxes} \\ \hline
Cameroonian examples & Personal income tax (IRPP) on salaries; company tax (IS) on firms' profits & VAT (19.25\%) on goods and services; customs duties at the Port of Douala; excise duties on drinks and tobacco \\ \hline
Advantage & Equitable: can be made progressive and adjusted to the taxpayer's ability to pay & Easy and cheap to collect; difficult to evade since it is built into prices; large and regular yield \\ \hline
Disadvantage & Costly to administer, encourages evasion and avoidance; may discourage work, saving and investment if rates are too high & Regressive: the poor spend a larger share of their income, so the tax weighs relatively more on them; raises prices (inflationary effect) \\ \hline
\end{tabularx}
\end{center}

\textbf{Conclusion:} a good tax system combines both types so as to balance equity (direct taxes) and yield/collectability (indirect taxes). In Cameroon, indirect taxes provide the larger share of tax revenue, which raises an equity concern.
\end{corrige}

\begin{corrige}[Exercise 6 — Tax incidence using supply-demand diagram]
\textbf{Mechanism.} A specific (per-unit) tax of 500 FCFA shifts the supply curve vertically \emph{upwards} by exactly 500 FCFA: for each quantity, producers now require the old price plus the tax. The new supply curve $S'$ is parallel to $S$.

\begin{popfigure}
\begin{center}
\begin{tikzpicture}[scale=1.0]
\draw[->, thick] (0,0) -- (7,0) node[below left] {$Q$};
\draw[->, thick] (0,0) -- (0,5.5) node[below left] {$P$};
\draw[very thick, popblue] (0.5,0.6) -- (6,4.6) node[right] {$S$};
\draw[very thick, poppurple] (0.5,1.6) -- (6,5.1) node[right] {$S'$ ($S+500$)};
\draw[very thick, poporange] (0.5,4.8) -- (6,1.0) node[right] {$D$};
\coordinate (E0) at (3.55,2.85);
\coordinate (E1) at (3.0,3.25);
\coordinate (P1) at (3.0,2.35);
\fill[popink] (E0) circle (2pt);
\fill[popink] (E1) circle (2pt);
\fill[popink] (P1) circle (2pt);
\draw[dashed, gray] (E0) -- (0,2.85) node[left] {$P_0$};
\draw[dashed, gray] (E1) -- (0,3.25) node[left] {$P_c$};
\draw[dashed, gray] (P1) -- (0,2.35) node[left] {$P_p$};
\draw[dashed, gray] (E0) -- (3.55,0) node[below] {$Q_0$};
\draw[dashed, gray] (E1) -- (3.0,0) node[below] {$Q_1$};
\draw[<->, popdark, thick] (3.35,2.35) -- (3.35,3.25) node[midway, right, font=\small] {500};
\end{tikzpicture}
\end{center}
\legende{Incidence of a specific indirect tax of 500 FCFA per unit}
\end{popfigure}

\textbf{Reading the diagram:}
\begin{enumerate}
\item Before the tax, equilibrium is at $(Q_0, P_0)$ where $D$ meets $S$.
\item The tax shifts supply from $S$ to $S'$. The new equilibrium quantity falls to $Q_1$.
\item \textbf{Price paid by consumers} rises to $P_c$ (on the demand curve).
\item \textbf{Price received by producers} falls to $P_p = P_c - 500$ (on the original supply curve).
\end{enumerate}

\textbf{Sharing of the burden:}
\begin{itemize}
\item Consumers bear $P_c - P_0$ per unit (the rise in the price they pay);
\item Producers bear $P_0 - P_p$ per unit (the fall in the price they receive);
\item The two shares sum to the tax: $(P_c - P_0) + (P_0 - P_p) = P_c - P_p = 500$ FCFA.
\end{itemize}

\textbf{Key remark:} the side of the market that is \emph{less elastic} bears the larger share. If demand is relatively inelastic (necessities such as kerosene or basic foodstuffs), consumers bear most of the tax; if supply is less elastic, producers bear most of it. The government collects $500 \times Q_1$ of revenue, and the fall from $Q_0$ to $Q_1$ represents the deadweight loss of the tax.
\end{corrige}

\begin{corrige}[Exercise 7 — Progressive income tax computation]
Taxable income: $7\,500\,000$ FCFA. In a progressive (bracket) system, each slice of income is taxed at its own rate.

\textbf{a) Total tax payable:}
\begin{align*}
\text{Bracket 1: } & 2\,000\,000 \times 10\% = 200\,000 \\
\text{Bracket 2: } & (5\,000\,000 - 2\,000\,000) \times 15\% = 3\,000\,000 \times 0.15 = 450\,000 \\
\text{Bracket 3: } & (7\,500\,000 - 5\,000\,000) \times 20\% = 2\,500\,000 \times 0.20 = 500\,000
\end{align*}
The income does not reach the 25\% bracket (above 10\,000\,000).
\[ \text{Total tax} = 200\,000 + 450\,000 + 500\,000 = \boxed{1\,150\,000 \text{ FCFA}}. \]

\textbf{b) Average tax rate:}
\[ \text{Average rate} = \frac{\text{Total tax}}{\text{Total income}} \times 100 = \frac{1\,150\,000}{7\,500\,000} \times 100 \approx 15.33\%. \]

\textbf{c) Marginal tax rate:} the rate applied to the \emph{last} franc earned, i.e. the rate of the highest bracket reached: $\boxed{20\%}$.

\textbf{Comment:} the average rate (15.33\%) is lower than the marginal rate (20\%): this is the signature of a progressive tax. The marginal rate is the one that matters for decisions about earning additional income.

\textbf{Common mistake:} Do NOT apply 20\% to the whole income ($7\,500\,000 \times 0.20 = 1\,500\,000$): that would be a proportional tax. Each bracket is taxed separately.
\end{corrige}

\begin{corrige}[Exercise 8 — Budgetary process and principles]
\textbf{Main stages of the budgetary process in Cameroon:}
\begin{enumerate}
\item \textbf{Preparation (elaboration):} the Ministry of Finance sends a budget circular to all ministries; each ministry prepares its estimates of revenue and expenditure; the draft finance bill is consolidated and adopted by the Government (Council of Ministers).
\item \textbf{Presentation and adoption (vote):} the finance bill is tabled before Parliament (National Assembly and Senate), examined in committee, debated and voted before the start of the budgetary year; it is then promulgated by the President of the Republic.
\item \textbf{Execution:} revenue is collected by the tax and customs administrations; expenditure is committed, verified, ordered and paid through the Treasury, under the control of the Ministry of Finance.
\item \textbf{Control:} internal control (ministerial and Treasury controls), parliamentary control (through the settlement/review law), and external control by the Audit Bench (Supreme State Control / Chamber of Accounts), which certifies the accounts.
\end{enumerate}

\textbf{Four fundamental budgetary principles:}
\begin{enumerate}
\item \textbf{Annuality:} the budget is voted and executed for one budgetary year (in Cameroon, the calendar year).
\item \textbf{Universality (unity of the budget):} all revenue and all expenditure must appear in a single document; revenue is not earmarked to specific expenditure (non-affectation rule).
\item \textbf{Balance (equilibrium):} voted expenditure must be covered by voted revenue; the budget is presented in balance.
\item \textbf{Speciality (specification):} appropriations are authorised for specific services and purposes and cannot be transferred freely between chapters without authorisation.
\end{enumerate}
(One may also cite \emph{sincerity} — the budget must give a faithful picture of the State's finances.)
\end{corrige}

\begin{corrige}[Exercise 9 — Structural question: Public finance components (10 marks)]
\textbf{Indicative mark scheme:} definition (2 marks); identification of the four components ($4 \times 1 = 4$ marks); Cameroonian illustration of each ($4 \times 1 = 4$ marks).

\textbf{Definition (2 marks).} \cle{Public finance} is the branch of economics that studies the financial operations of the State and other public authorities: how they obtain resources (revenue), how they use them (expenditure), how these operations are authorised and planned (the budget) and how gaps are financed (public debt), with the aim of allocating resources, redistributing income and stabilising the economy.

\textbf{The four components, with Cameroonian examples (8 marks):}
\begin{enumerate}
\item \textbf{Public revenue:} all resources of the State — fiscal revenue (taxes, duties), non-fiscal revenue (fees, rents, dividends) and exceptional revenue (grants, borrowing). \emph{Example:} VAT and company tax collected by the Directorate General of Taxes, customs duties at the Port of Douala, and oil royalties from crude oil sales.
\item \textbf{Public expenditure:} all spending of public authorities — current expenditure (salaries, goods and services, interest on debt) and capital expenditure (public investment). \emph{Example:} salaries of civil servants and teachers, construction of roads (e.g. the Yaoundé–Douala motorway project), hospitals, schools and sports infrastructure.
\item \textbf{The budget:} the annual legislative act (finance law) that forecasts and authorises public revenue and expenditure for the year. \emph{Example:} the finance law voted each year by the National Assembly and the Senate and promulgated by the President of the Republic.
\item \textbf{Public debt:} the stock of loans contracted by the State (internal and external) to finance budget deficits or major investments. \emph{Example:} Treasury bonds issued on the CEMAC regional market (BEAC/BVMAC) and external borrowing from partners such as the World Bank, the African Development Bank or bilateral lenders.
\end{enumerate}

\textbf{Examiner expectations:} a precise definition mentioning both revenue and expenditure; four clearly distinguished components (not confused with one another); examples that are genuinely Cameroonian and correctly classified (e.g. VAT is revenue, not expenditure). A short introduction and conclusion are appreciated at this level.
\end{corrige}

\begin{corrige}[Exercise 10 — Essay: Effectiveness of fiscal policy in Cameroon (20 marks)]
\textbf{Indicative mark scheme:} introduction (2); achievements (6); limitations (6); suggested measures (4); conclusion (2).

\textbf{Introduction (2 marks).} Fiscal policy is the use of public revenue (taxation) and public expenditure to influence economic activity, employment, prices and income distribution. In Cameroon it is implemented through the annual finance law, tax reforms and public investment programmes. The statement claims it has been \emph{effective}: this must be discussed by weighing achievements against limitations before giving a balanced verdict.

\textbf{I. Achievements of fiscal policy in Cameroon (6 marks — any three, well developed):}
\begin{enumerate}
\item \textbf{Financing of public investment:} taxation and borrowing have financed major infrastructure (roads, ports, hydroelectric dams, stadiums), supporting growth and integration of the national territory.
\item \textbf{Modernisation of tax administration:} creation of large-taxpayer offices, computerisation, online declaration and payment, and the generalisation of VAT have broadened the tax base and increased non-oil revenue.
\item \textbf{Counter-cyclical role:} during shocks (fall in oil prices, security and health crises), public spending sustained demand, protected employment in the public sector and supported affected households and firms.
\item \textbf{Redistribution and social spending:} budgetary transfers finance education, health and subsidies on some basic products, attenuating inequalities and poverty.
\end{enumerate}

\textbf{II. Limitations (6 marks — any three, well developed):}
\begin{enumerate}
\item \textbf{Narrow tax base and evasion:} the very large informal sector escapes taxation; fraud, smuggling and corruption reduce revenue, so the tax burden falls heavily on formal firms and salaried workers.
\item \textbf{Regressive structure:} revenue relies heavily on indirect taxes (VAT, customs duties) which weigh relatively more on the poor, limiting the redistributive effect.
\item \textbf{Weak execution and inefficiency of expenditure:} slow disbursement of investment budgets, overruns, and sometimes unproductive current spending reduce the multiplier effect of public expenditure.
\item \textbf{Dependence and rigidity:} dependence on oil revenue and on a few large taxpayers makes revenue volatile; membership of CEMAC (CFA franc zone) limits monetary-fiscal room for manoeuvre; rising debt service crowds out social and capital spending.
\end{enumerate}

\textbf{III. Measures to improve effectiveness (4 marks — any two or three):}
\begin{enumerate}
\item Broaden the tax base by gradually formalising the informal sector (simplified tax regimes, incentives) rather than raising rates on the same taxpayers.
\item Strengthen tax administration: digitalisation, fight against corruption and smuggling, better control of tax exemptions.
\item Improve the quality of expenditure: programme budgeting with performance indicators, rigorous project selection and evaluation, faster execution of investment budgets.
\item Make taxation more equitable (stronger progressivity, relief on basic necessities) and manage public debt prudently to preserve budgetary room.
\end{enumerate}

\textbf{Conclusion (2 marks).} Fiscal policy in Cameroon has recorded real achievements — modernisation of revenue collection and financing of structuring investment — but its effectiveness remains limited by a narrow base, evasion, inefficiency of spending and debt pressure. It is therefore only \emph{partially} effective; the suggested reforms could make it a genuine instrument of growth and redistribution.

\textbf{Examiner expectations:} a clear thesis and plan; balanced treatment (both sides of the discussion); Cameroonian illustrations; correct economic vocabulary (multiplier, tax base, progressivity, crowding-out); a personal, justified conclusion. Mere listing without explanation earns half marks.
\end{corrige}

\begin{corrige}[Exercise 11 — Public debt analysis (15 marks)]
\textbf{a) Budget deficit vs public debt (4 marks).}
The \cle{budget deficit} is a \emph{flow}: the excess of public expenditure over public revenue during \emph{one} budgetary year. The \cle{public debt} is a \emph{stock}: the accumulation, at a given date, of all loans contracted by the State and not yet repaid. The link between them is dynamic: when deficits are financed by borrowing, they add to the debt year after year:
\[ \text{Debt}_t = \text{Debt}_{t-1} + \text{Deficit financed by borrowing}_t . \]
Thus the deficit is like the water flowing into a bathtub each year, while the debt is the total quantity of water accumulated in the tub. (2 marks per concept correctly defined with the flow/stock distinction.)

\textbf{b) Two major causes of the rise in Cameroon's public debt (6 marks — 3 each):}
\begin{enumerate}
\item \textbf{Persistent budget deficits linked to ambitious investment programmes:} the financing of major structuring projects (roads, ports, hydroelectric dams, sports infrastructure) exceeded domestic revenue, forcing the State to borrow massively on the CEMAC regional market (Treasury bonds) and externally.
\item \textbf{Exogenous shocks and fall in revenue:} the decline in world oil prices reduced oil revenue, while security crises (Far North, North-West and South-West) increased security and humanitarian expenditure and depressed economic activity and tax receipts; the gap was covered by new borrowing.
\end{enumerate}
(Other acceptable causes: exchange-rate and refinancing risks on external debt, borrowing to repay maturing debt — ``rolling over'' — and support to struggling public enterprises.)

\textbf{c) Two economic consequences of a high public debt (5 marks):}
\begin{enumerate}
\item \textbf{Debt service crowds out productive spending:} interest payments and amortisation absorb a growing share of the budget, reducing the resources available for health, education and investment — the \emph{crowding-out} effect within the budget. This slows long-term growth.
\item \textbf{Higher taxation and loss of credibility:} to service the debt, the State may raise taxes (discouraging work, investment and consumption) and faces higher risk premiums from lenders; excessive debt can threaten debt sustainability, force adjustment programmes and reduce future budgetary room for manoeuvre.
\end{enumerate}
(Acceptable alternatives: inflationary pressure if deficits are monetised; burden shifted to future generations — intergenerational inequity.)

\textbf{Examiner expectations:} a rigorous flow/stock distinction in (a); causes that are explained (mechanism), not merely named; consequences analysed with their transmission mechanism; Cameroonian context (CEMAC market, oil revenue, security crises) rewarded.
\end{corrige}

\newpage

\coursec{Summary sheet}

\begin{popfigure}
\begin{tikzpicture}[scale=0.8, every node/.style={transform shape}]
% Central node
\node[circle, draw=popdark, fill=popblue, text=white, minimum size=1.8cm, font=\bfseries] (center) {Public Finance};
% Branch nodes
\node[rounded corners, fill=popblue, text=white, text width=2.5cm, align=center, font=\small] (b1) at (0:4.5) {Meaning \& Scope};
\node[rounded corners, fill=popgreen, text=white, text width=2.5cm, align=center, font=\small] (b2) at (51.4:4.5) {Taxation: Principles};
\node[rounded corners, fill=poporange, text=white, text width=2.5cm, align=center, font=\small] (b3) at (102.9:4.5) {Types of Taxes};
\node[rounded corners, fill=poppurple, text=white, text width=2.5cm, align=center, font=\small] (b4) at (154.3:4.5) {VAT in Cameroon};
\node[rounded corners, fill=poppink, text=white, text width=2.5cm, align=center, font=\small] (b5) at (205.7:4.5) {Economic Consequences};
\node[rounded corners, fill=popgold, text=white, text width=2.5cm, align=center, font=\small] (b6) at (257.1:4.5) {Budget \& Programmed Budget};
\node[rounded corners, fill=popdark, text=white, text width=2.5cm, align=center, font=\small] (b7) at (308.6:4.5) {Public Debt \& Fiscal Policy};
% Edges
\foreach \n in {b1,b2,b3,b4,b5,b6,b7} {
  \draw[popline, thick] (center) -- (\n);
}
\end{tikzpicture}
\legende{Mind map of the Public Finance chapter}
\end{popfigure}

\begin{bilan}
\textbf{Key Definitions}
\begin{itemize}
\item \cle{Public Finance}: The branch of economics that studies government revenue, expenditure, and debt management.
\item \cle{Tax}: A compulsory, unrequited levy imposed by the state on individuals or legal entities.
\item \cle{VAT (Value Added Tax)}: A multi-stage consumption tax levied on the value added at each stage of production and distribution.
\item \cle{Budget}: A legal document forecasting government revenue and authorising expenditure for a fiscal year.
\item \cle{Programmed Budget}: A budget that links expenditures to specific objectives, programmes, and performance indicators.
\item \cle{Public Debt}: The stock of financial obligations incurred by the government through borrowing.
\end{itemize}

\textbf{Canons of Taxation (Adam Smith + Modern)}
\begin{itemize}
\item \cle{Equity}: Horizontal equity (equals treated equally) and vertical equity (unequals treated differently).
\item \cle{Certainty}: Tax liability must be clear, predictable, and not arbitrary.
\item \cle{Convenience}: Payment should be easy and timed appropriately for the taxpayer.
\item \cle{Economy}: Collection costs should be minimal relative to revenue yielded.
\item \cle{Flexibility}: The tax system should adapt to changing economic conditions.
\item \cle{Neutrality}: Taxes should minimise distortion of economic decisions.
\end{itemize}

\textbf{Classification of Taxes}
\begin{itemize}
\item \cle{Direct Taxes}: Levied on income, wealth, or capital (e.g., personal income tax, corporate tax, property tax). Incidence cannot be shifted.
\item \cle{Indirect Taxes}: Levied on expenditure (e.g., VAT, excise duties, customs duties). Incidence can be shifted to consumers.
\item \cle{By rate structure}: \cle{Proportional} (constant rate), \cle{Progressive} (rate rises with base), \cle{Regressive} (rate falls with base).
\item \cle{By base}: Income taxes, consumption taxes, wealth taxes, trade taxes.
\end{itemize}

\textbf{VAT in Cameroon (Practical Work 21)}
\begin{itemize}
\item Standard rate: \textbf{19.25\%} (since 2019).
\item Mechanism: Output VAT (on sales) minus Input VAT (on purchases) = VAT payable to the tax administration.
\item Example: A firm sells goods for \textbf{1,000,000 FCFA} (VAT inclusive). VAT-exclusive price = 1,000,000 / 1.1925 ≈ 838,600 FCFA. VAT = 1,000,000 – 838,600 = \textbf{161,400 FCFA}.
\item Exemptions: Basic foodstuffs, pharmaceutical products, educational services, exports (zero-rated).
\item Filing: Monthly declaration (CA3) by the 15th of the following month.
\end{itemize}

\textbf{Structure of Taxation in Cameroon (Practical Work 22)}
\begin{itemize}
\item Dominance of indirect taxes (VAT, customs duties, excise) ≈ 60\% of tax revenue.
\item Direct taxes (corporate tax, personal income tax) ≈ 30\%.
\item Non-tax revenue (oil royalties, dividends from state enterprises) significant but volatile.
\item Tax/GDP ratio low (≈ 13-15\%), below CEMAC convergence criterion (17\%).
\item Ongoing reforms: broadening base, digitalising collection (e-tax), reducing exemptions.
\end{itemize}

\textbf{Economic Consequences of Taxation}
\begin{itemize}
\item \cle{Allocation effect}: Taxes distort relative prices → resource misallocation (e.g., excise on alcohol reduces consumption).
\item \cle{Distribution effect}: Progressive taxes reduce income inequality; regressive taxes worsen it.
\item \cle{Stabilisation effect}: Progressive income tax acts as automatic stabiliser (revenues fall in recession, rise in boom).
\item \cle{Incentive effect}: High marginal rates may discourage labour supply, saving, investment, entrepreneurship.
\item \cle{Incidence}: Statutory incidence (who writes the cheque) vs. economic incidence (who bears the burden). Depends on price elasticities of demand and supply.
\item \cle{Excess burden (deadweight loss)}: Loss of consumer/producer surplus beyond revenue collected; larger with elastic curves.
\end{itemize}

\textbf{The State Budget and Programmed Budget}
\begin{itemize}
\item \cle{Budget cycle}: Preparation (MINEPAT), Parliamentary approval, Execution (MINFI), Audit (Court of Auditors).
\item \cle{Programmed budget}: Introduced 2013; classifies spending by \cle{programmes} (e.g., education, health) with objectives, indicators, targets.
\item \cle{Budget balance}: Surplus (revenue > expenditure), Deficit (expenditure > revenue), Balanced.
\item \cle{Deficit financing}: Domestic borrowing (treasury bills, bonds), External borrowing (multilateral, bilateral), Monetary financing (limited by CEMAC statutes).
\item \cle{CEMAC convergence criteria}: Deficit ≤ 3\% of GDP, Debt ≤ 70\% of GDP, Inflation ≤ 3\%, Non-accumulation of arrears.
\end{itemize}

\textbf{Public Debt and Fiscal Policy Effectiveness in Cameroon}
\begin{itemize}
\item \cle{Composition}: Domestic debt (treasury bills, bonds) vs. External debt (IDA, AfDB, bilateral, Eurobonds).
\item \cle{Sustainability indicators}: Debt-to-GDP, Debt-service-to-revenue, Present value of debt-to-exports.
\item \cle{Debt management}: Refinancing, rescheduling, HIPC/MDRI relief (completed 2006), current IMF ECF programme.
\item \cle{Effectiveness of fiscal policy}: Constrained by CFA franc fixed parity (no independent monetary policy), narrow tax base, large informal sector, rigid expenditure (wages, subsidies), weak revenue administration.
\item \cle{Policy mix}: Need for revenue mobilisation (VAT efficiency, property tax), expenditure rationalisation (subsidy reform), and structural reforms to boost growth.
\end{itemize}

\textbf{Key Formulas / Relationships}
\begin{itemize}
\item VAT payable = Output VAT – Input VAT.
\item Tax incidence on consumers (specific tax) = $\frac{E_s}{E_s+E_d} \times t$ where $E_s$, $E_d$ are supply and demand elasticities.
\item Laffer curve: Tax revenue $R = \tau \times B(\tau)$; maximum at $\tau^*$ where $\frac{dR}{d\tau}=0$.
\item Budget deficit = Total Expenditure – Total Revenue.
\item Debt-to-GDP ratio = $\frac{\text{Total Public Debt}}{\text{GDP}} \times 100$.
\item Primary deficit = Total Expenditure – Interest Payments – Total Revenue.
\end{itemize}

\textbf{Common Mistakes (GCE Traps)}
\begin{itemize}
\item Confusing \emph{statutory incidence} with \emph{economic incidence}.
\item Treating VAT as a single-stage sales tax; forgetting the credit-invoice mechanism.
\item Mislabeling a proportional tax as progressive because the \emph{amount} rises with income.
\item Assuming a budget deficit automatically means unsustainable debt (ignores financing composition and growth).
\item Applying the Laffer curve as a precise policy rule rather than a theoretical concept.
\item Not distinguishing between \emph{domestic} and \emph{external} debt sustainability thresholds.
\item Forgetting CEMAC constraints (no monetary financing, convergence criteria) when analysing Cameroon's fiscal policy.
\end{itemize}

\textbf{GCE Examination Tips}
\begin{itemize}
\item \textbf{Definitions}: Give precise, textbook definitions (e.g., "Public finance is the study of how the government raises and spends money...").
\item \textbf{Diagrams}: Draw AD/AS shifts for tax changes, Laffer curve, Lorenz curve for distribution, VAT credit mechanism flowchart.
\item \textbf{VAT calculations}: Show all steps: find VAT-exclusive price, compute VAT amount, determine payable/refund.
\item \textbf{Essay structure}: Introduction (definitions, context), Analysis (theory + diagrams + Cameroon examples), Evaluation (advantages/limits, short vs. long run), Conclusion (policy recommendation).
\item \textbf{Cameroon context}: Mention CFA franc constraints, SONARA subsidies, cocoa/coffee export taxes, AfCFTA implications, IMF programmes.
\item \textbf{Time management}: Allocate ~25 minutes per essay question; 10 minutes for short-answer/calculation questions.
\item \textbf{Command words}: "Explain" requires reasons; "Analyse" requires diagrams and mechanisms; "Evaluate" requires pros/cons and judgement.
\end{itemize}
\end{bilan}

\findecours

\end{document}
